\documentclass[11pt,oneside,a4paper]{report}

% ---------- course configuration ----------
\def\utprunninghead{Operating Systems Notes}
\def\utplabtotal{10}
\def\utpcoursetitle{Operating Systems}
\def\utpdocumentsubtitle{Teaching Assistant Guide}
\def\utpcourseoutline{%
  A step-by-step path from the Linux command line and process model through CPU scheduling, concurrency, memory management, and file systems, with one practical lab paired to each core chapter.%
}

% ---------- shared UTP teaching-note framework ----------
% The files in ../.shared are supplied separately and are intentionally
% not modified by this course.
% =====================================================================
%  utpnotes.tex
%  Shared preamble for the UTP teaching-assistant course notes.
%
%  Every book includes this file, and it is the ONLY file a book needs
%  from .shared/: the cover, the logos and the optional diagram
%  libraries are all reached from here. It owns everything the books
%  have in common, so a change here reaches every course at once.
%
%  ------------------------------------------------------------------
%  HOW A COURSE USES THIS FILE
%  ------------------------------------------------------------------
%      \documentclass[11pt,oneside,a4paper]{report}
%
%      \def\utprunninghead{Machine Learning Notes}   % section 1
%      \def\utplabtotal{8}
%      \def\utpcoursetitle{Machine Learning}
%      \def\utpcourseoutline{One or two lines for the cover.}
%
%      % =====================================================================
%  utpnotes.tex
%  Shared preamble for the UTP teaching-assistant course notes.
%
%  Every book includes this file, and it is the ONLY file a book needs
%  from .shared/: the cover, the logos and the optional diagram
%  libraries are all reached from here. It owns everything the books
%  have in common, so a change here reaches every course at once.
%
%  ------------------------------------------------------------------
%  HOW A COURSE USES THIS FILE
%  ------------------------------------------------------------------
%      \documentclass[11pt,oneside,a4paper]{report}
%
%      \def\utprunninghead{Machine Learning Notes}   % section 1
%      \def\utplabtotal{8}
%      \def\utpcoursetitle{Machine Learning}
%      \def\utpcourseoutline{One or two lines for the cover.}
%
%      % =====================================================================
%  utpnotes.tex
%  Shared preamble for the UTP teaching-assistant course notes.
%
%  Every book includes this file, and it is the ONLY file a book needs
%  from .shared/: the cover, the logos and the optional diagram
%  libraries are all reached from here. It owns everything the books
%  have in common, so a change here reaches every course at once.
%
%  ------------------------------------------------------------------
%  HOW A COURSE USES THIS FILE
%  ------------------------------------------------------------------
%      \documentclass[11pt,oneside,a4paper]{report}
%
%      \def\utprunninghead{Machine Learning Notes}   % section 1
%      \def\utplabtotal{8}
%      \def\utpcoursetitle{Machine Learning}
%      \def\utpcourseoutline{One or two lines for the cover.}
%
%      \input{../.shared/utpnotes}
%
%      \utpcontentspage{\small}{\normalsize}{0.65em}{0.07em}{0.16em}
%
%      \begin{document}
%      \utpfrontmatter                     % cover, contents, map
%      \part{Orientation}
%      \input{chapters/01-where-to-begin}
%      \end{document}
%
%  Configuration goes BEFORE the \input; overrides go AFTER it.
%  Paths are relative, so compile from inside the course folder.
%
%  ------------------------------------------------------------------
%  NAMING
%  ------------------------------------------------------------------
%  One name per idea, used by every book. Anything a book can use is
%  either an unprefixed learning box (keyidea, pitfall, labproject,
%  ...) or carries the utp prefix (utpfigure, utptable, utpfig,
%  utpsteps, utppython, UTPNavy, utpcell, ...). Course prefixes such
%  as ml-, ds- or erp- are not used: a chapter moved from one book to
%  another compiles unchanged. .shared/README.md lists every name.
%
%  Adding a new course? Copy .shared/TEMPLATE-main.tex and follow
%  .shared/ADDING-A-COURSE.md. Every step is a numbered checklist.
%
%  ------------------------------------------------------------------
%  CONTENTS
%  ------------------------------------------------------------------
%    1  Course configuration     the knobs a book may set
%    2  Packages                 the shared package set and fonts
%    3  Colour palette           the one UTP palette
%    4  Page layout and tables   margins, lists, table styling
%    5  Table of contents        parts, chapters, labs, figures
%    6  Headings                 chapter badge, part pages, sections
%    7  Page furniture           running heads and footers
%    8  Learning boxes           key idea, pitfall, try it, and friends
%    9  Small commands           \term, \cmd, \chapterguide, utpsteps
%   10  Diagram node styles      the TikZ vocabulary every book shares
%   11  Figure frames            utpfigure, utpfig, source images
%   12  Lab sheets               the lab banner and the lab toolkit
%   13  Stage divider            the part-opening page in longer books
%   14  Front matter             \utpfrontmatter: cover, contents, map
%   15  Listing styles           code, terminal and pseudocode
% =====================================================================

% --------------------------------------------------------------------
%  1.  COURSE CONFIGURATION
% --------------------------------------------------------------------
%  The knobs below, plus \utpstyleparts in section 5 and the box titles
%  in section 8, are the whole configuration surface. A book sets what
%  it needs BEFORE the \input and inherits the rest.
%
%  \providecommand means "define unless the book already did", which is
%  what makes the \def-before-\input pattern work.
%
%  Running text
%    \utprunninghead   the left running head on every page.
%    \utpchapterlabel  the word in the chapter badge. CHAPTER by
%                      default; a book whose chapters are labs sets LAB.
%    \utplabtotal      the N in the "Lab 3 of N" banners.

\providecommand{\utprunninghead}{Course Notes}
\providecommand{\utpchapterlabel}{CHAPTER}
\providecommand{\utplabtotal}{10}

%  Cover, printed by section 14
%    \utpcoursetitle       the title on the title page.
%    \utpdocumentsubtitle  the line under the title; Teaching Assistant Guide by default.
%    \utpcourseoutline     the sentence or two below it; \\ breaks a line.
%    \utpoutlineheading    the small heading above that sentence.
%    \utpauthor            the name in the PREPARED BY block.
%    \utpauthorrole        the line under the name.

\providecommand{\utpcoursetitle}{\utprunninghead}
\providecommand{\utpdocumentsubtitle}{Teaching Assistant Guide}
\providecommand{\utpcourseoutline}{}
\providecommand{\utpoutlineheading}{COURSE OUTLINE}
\providecommand{\utpauthor}{Reeyan Afzal Khan}
\providecommand{\utpauthorrole}{Graduate Research Assistant}

%  Files and folders
%    \utpextragraphicspath  image folders to search before the standard
%                           ones, as brace groups: {{screenshots/}}.
%                           A course's own figures/ folder and the
%                           shared logos are always searched.
%    \utpcoursemap          the overview chapter \utpfrontmatter inputs
%                           after the contents. Empty means none.

\providecommand{\utpextragraphicspath}{}
\providecommand{\utpcoursemap}{chapters/course-map}

%  Contents entries for chapters (section 5 applies them). A book whose
%  chapters are labs says so once here rather than restating the three
%  tocloft settings after the input.
%    \utpchaptertocword      Chapter~ by default, Lab~ in one book.
%    \utpchaptertocindent    left indent of the entry.
%    \utpchaptertocnumwidth  room reserved for the word and number.

\providecommand{\utpchaptertocword}{Chapter~}
\providecommand{\utpchaptertocindent}{1.5em}
\providecommand{\utpchaptertocnumwidth}{6em}

%  Monospace face (section 2)
%    \utpmonofont   sourcecodepro by default. inconsolata is the
%                   alternative, for a book whose pages are mostly
%                   device transcripts.

\providecommand{\utpmonofont}{sourcecodepro}

% --------------------------------------------------------------------
%  2.  PACKAGES
% --------------------------------------------------------------------
%  The union of what the books need. A course that needs something
%  extra loads it in its own file, under the heading "packages only
%  this book needs". Pseudocode is the one optional package the shared
%  file knows about: \utpusealgorithms in section 15 loads and styles it.
%
%  The tikz and pgfplots library lists are the union across all books,
%  so a diagram or chart copied from one course compiles in another.

\usepackage[T1]{fontenc}
\IfFileExists{glyphtounicode.tex}{\input{glyphtounicode}\pdfgentounicode=1}{}
\usepackage[american]{babel}
\IfFileExists{sourcesanspro.sty}
  {\usepackage[default]{sourcesanspro}}
  {\usepackage[scaled]{helvet}\renewcommand{\familydefault}{\sfdefault}}
\DeclareFontShape{T1}{cmss}{b}{n}{<->ssub*cmss/bx/n}{}
\DeclareFontShape{TS1}{cmss}{b}{n}{<->ssub*cmss/bx/n}{}
% Monospace. Without this the tt default stays Computer Modern Typewriter,
% a serif face that clashes with the sans body text in every listing and
% every \cmd. Source Code Pro is the companion to Source Sans Pro; it is
% scaled down because mono digits otherwise sit large next to the sans.
%
% Inconsolata, chosen with \def\utpmonofont{inconsolata}, distinguishes
% 0/O and 1/l/I at a glance, which matters when a reader copies an
% interface name or an address out of a transcript. Only one mono package
% is ever loaded: a second would win the family but inherit the first's
% series names, and every \ttfamily would then resolve by substitution.
\def\utpmonoinconsolata{inconsolata}
\ifx\utpmonofont\utpmonoinconsolata
  \IfFileExists{inconsolata.sty}{\usepackage[varqu,varl,scaled=1.02]{inconsolata}}{}
\else
  \IfFileExists{sourcecodepro.sty}
    {\usepackage[scaled=0.88]{sourcecodepro}}
    {\usepackage[scaled=0.86]{beramono}}
  % Source Sans names its normal weight "regular", but the Computer Modern
  % fallbacks call it "m". Map one onto the other so a \ttfamily reached
  % while the regular series is active resolves instead of warning.
  \DeclareFontShape{T1}{cmtt}{regular}{n}{<->ssub*cmtt/m/n}{}
  \DeclareFontShape{TS1}{cmtt}{regular}{n}{<->ssub*cmtt/m/n}{}
  \DeclareFontShape{T1}{cmtt}{regular}{it}{<->ssub*cmtt/m/it}{}
\fi
\usepackage[margin=21mm,headheight=22pt]{geometry}
\usepackage{amsmath,amssymb,mathtools,bm}
\usepackage{microtype}
\usepackage{enumitem}
\usepackage{booktabs}
\usepackage{tabularx}
\usepackage{array}
\usepackage{ragged2e}
\usepackage{etoolbox}
\usepackage{titlesec}
\usepackage{tocloft}
\usepackage{fancyhdr}
\usepackage[table]{xcolor}
\usepackage{tcolorbox}
\tcbuselibrary{breakable,skins}
\usepackage[hidelinks]{hyperref}
\usepackage{bookmark}
% Document properties, so the PDF identifies itself in a reader's
% title bar, in a library catalogue, and in search results. Set from
% the same configuration macros the cover uses.
\hypersetup{
  pdftitle={\utpcoursetitle},
  pdfauthor={\utpauthor},
  pdfsubject={Universiti Teknologi PETRONAS course notes and laboratory guide},
  pdfkeywords={UTP, course notes, laboratory manual, \utpcoursetitle},
  pdfdisplaydoctitle=true,
  bookmarksnumbered=true,
  bookmarksopenlevel=1
}
\usepackage{listings}
\usepackage{graphicx}
\usepackage{caption}
\usepackage{tikz}
\usetikzlibrary{angles,arrows.meta,backgrounds,calc,chains,decorations.pathmorphing,decorations.pathreplacing,fit,matrix,patterns,positioning,quotes,shadows.blur,shapes.geometric,shapes.misc,shapes.symbols}
\usepackage{pgfplots}
% statistics gives boxplot prepared; groupplots gives multi-panel charts.
\usepgfplotslibrary{statistics}
\usepgfplotslibrary{groupplots}

% Images: any folders the book names first, then its own figures/
% folder, then the shared logos used by the cover.
\expandafter\graphicspath\expandafter{\utpextragraphicspath{figures/}{../.shared/logos/}}

% --------------------------------------------------------------------
%  3.  COLOUR PALETTE
% --------------------------------------------------------------------
%  One palette, defined once under UTP* names, and used by every book,
%  box, diagram and chart. Five accents (blue, teal, green, orange, red)
%  with a pale tint of each for box backgrounds, plus navy for headings,
%  gold for the cover, ink and muted for text, and a sixth accent,
%  purple, for charts and layered diagrams that need one more series.

\definecolor{UTPNavy}{HTML}{18324A}
\definecolor{UTPBlue}{HTML}{246B9E}
\definecolor{UTPTeal}{HTML}{138A80}
\definecolor{UTPGreen}{HTML}{2D7D46}
\definecolor{UTPOrange}{HTML}{C86B24}
\definecolor{UTPRed}{HTML}{B34843}
\definecolor{UTPPurple}{HTML}{6B4FA0}
\definecolor{UTPGold}{HTML}{D4A13E}
\definecolor{UTPInk}{HTML}{263238}
\definecolor{UTPMuted}{HTML}{60717C}
\definecolor{UTPPaleBlue}{HTML}{EEF6FB}
\definecolor{UTPPaleTeal}{HTML}{EEF8F6}
\definecolor{UTPPaleOrange}{HTML}{FFF5EA}
\definecolor{UTPPaleRed}{HTML}{FFF0EF}
\definecolor{UTPPalePurple}{HTML}{F0ECF9}
\definecolor{UTPPaleGray}{HTML}{F4F6F7}
\color{UTPInk}

% --------------------------------------------------------------------
%  4.  PAGE LAYOUT AND TABLES
% --------------------------------------------------------------------
%  Paragraph spacing, list indents, and the table look: alternating
%  row tint, top-aligned paragraph columns, ragged-right cells.
%
%  Two extra column types for tabularx tables:
%      D{width}   a fixed-width, ragged-right, vertically centred column
%      Y          a ragged-right column that fills the remaining width
%
%  Every table in every book has the same shape, the one Algorithms &
%  Data Structures set: a utptable frame around a full-width tabularx,
%  booktabs rules, a tinted header row of \utpth cells (centred, bold,
%  navy), and a first column in bold navy when it names the rows.
%  Cells are vertically centred, so use m{} rather than p{} columns.
%
%      \begin{utptable}
%      \begin{tabularx}{\linewidth}{@{}>{\bfseries\leavevmode\color{UTPNavy}}m{0.22\linewidth}X@{}}
%      \toprule
%      \rowcolor{UTPPaleBlue!92}
%      \utpth{Operation} & \utpth{What it does} \\
%      \midrule
%      Insertion & Add a new element. \\
%      \bottomrule
%      \end{tabularx}
%      \end{utptable}
%
%  A book whose tables carry dense text asks for \utpcompacttables after
%  the input: tabularx tables then set in \footnotesize with tighter
%  padding and vertically centred cells.

\setlength{\parindent}{0pt}
\setlength{\parskip}{0.55em}
\setlist[itemize]{leftmargin=1.45em,itemsep=0.22em,topsep=0.25em}
\setlist[enumerate]{leftmargin=1.75em,itemsep=0.28em,topsep=0.3em}
\setlist[description]{font=\bfseries\color{UTPNavy},itemsep=0.25em,topsep=0.3em}
\renewcommand{\arraystretch}{1.34}
\setlength{\tabcolsep}{6.5pt}
\arrayrulecolor{UTPBlue!45}
% Keep paragraph-style columns ragged and vertically centred. Tables use
% fixed m columns beside flexible X columns; giving X the same centring
% keeps a short cell level with a tall neighbour in a multi-line row.
\renewcommand{\tabularxcolumn}[1]{m{#1}}
\makeatletter
\g@addto@macro\@arrayparboxrestore{\RaggedRight\arraybackslash}
\makeatother
% A header cell: centred, bold and navy, on the tinted header row.
\newcommand{\utpth}[1]{\multicolumn{1}{c}{\bfseries\color{UTPNavy}#1}}
\AtBeginEnvironment{tabular}{\small\setlength{\extrarowheight}{0.6pt}\rowcolors{2}{UTPPaleBlue!72}{white}}
\AtBeginEnvironment{tabularx}{\small\setlength{\extrarowheight}{0.6pt}\rowcolors{2}{UTPPaleBlue!72}{white}}
\AtBeginEnvironment{figure}{\centering}
\newcolumntype{D}[1]{>{\raggedright\arraybackslash}m{#1}}
\newcolumntype{Y}{>{\raggedright\arraybackslash}X}
\newcommand{\utpcompacttables}{%
  \renewcommand{\tabularxcolumn}[1]{m{##1}}%
  \AtBeginEnvironment{tabularx}{%
    \footnotesize
    \setlength{\tabcolsep}{4pt}%
    \renewcommand{\arraystretch}{1.25}%
    \setlength{\extrarowheight}{0.6pt}%
  }%
}
\emergencystretch=1.5em
\hbadness=10000
\hfuzz=8pt

% Page-breaking discipline. Without these a paragraph can leave a single
% line stranded at the foot of one page or the head of the next, which is
% the most visible sign of unattended typesetting.
\widowpenalty=10000          % no first line of a paragraph left alone at a page foot
\clubpenalty=10000           % no last line of a paragraph left alone at a page head
\displaywidowpenalty=10000   % the same around displayed equations
\brokenpenalty=10000         % do not break a page after a hyphenated line
\predisplaypenalty=10000     % keep the line before a display with it
% Pages are built from boxes of many heights, so let the last line float
% up rather than stretching the gaps between boxes to force a flush foot.
\raggedbottom
% Prefer a slightly loose line to an overfull one, but do not hyphenate
% freely -- short sans lines look ragged when over-hyphenated.
\tolerance=1200
\hyphenpenalty=200

% Character protrusion and font expansion. Both are what stop a justified
% sans column from showing visible rivers down the page.
\microtypesetup{protrusion=true,expansion=true,tracking=true}

\pgfplotsset{
  compat=1.18,
  every axis/.append style={
    axis line style={draw=UTPMuted!85,line width=0.55pt},
    tick style={draw=UTPMuted!85},
    tick label style={font=\scriptsize\color{UTPMuted}},
    label style={font=\small\color{UTPNavy}},
    title style={font=\small\bfseries\color{UTPNavy}},
    legend style={font=\scriptsize,draw=UTPBlue!22,fill=white,rounded corners=1pt}
  },
  every axis plot/.append style={line width=1.15pt}
}
\tikzset{every picture/.append style={line cap=round,line join=round}}

% --------------------------------------------------------------------
%  5.  TABLE OF CONTENTS
% --------------------------------------------------------------------
%  Formatting for part, chapter and lab entries, and for a list of
%  figures if a book prints one. The lab level is a custom tocloft list
%  so lab banners appear in the contents.
%
%  Part styling is skipped when \utpstyleparts is 0, which is what a
%  book without \part divisions should set.
%
%  The contents page is meant to fill exactly one page, and how loosely
%  it is set depends on how many entries a book has. One line after the
%  input tunes it:
%
%      \utpcontentspage{<base size>}{<part size>}%
%                      {<space before a part>}{<before a chapter>}{<before a lab>}
%      \utpcontentspage{\normalsize}{\large}{1.15em}{0.38em}{0.30em}   % the default

\addto\captionsamerican{\renewcommand{\contentsname}{Learning Path}}
\newlistentry{lab}{toc}{-1}
\providecommand{\utpstyleparts}{1}
\ifnum\utpstyleparts=1
\renewcommand{\cftpartpresnum}{}
\renewcommand{\cftpartaftersnum}{}
\setlength{\cftpartnumwidth}{5.2em}
\setlength{\cftbeforepartskip}{0.38em}
\renewcommand{\cftpartfont}{\normalsize\bfseries\color{UTPNavy}}
\renewcommand{\cftpartpagefont}{\normalsize\bfseries\color{UTPNavy}}
\renewcommand{\cftpartleader}{\hfill}
\fi
\renewcommand{\cftchappresnum}{\utpchaptertocword}
\renewcommand{\cftchapaftersnum}{\quad}
\setlength{\cftchapindent}{\utpchaptertocindent}
\setlength{\cftchapnumwidth}{\utpchaptertocnumwidth}
\setlength{\cftbeforechapskip}{0pt}
\renewcommand{\cftchapfont}{\normalfont}
\renewcommand{\cftchappagefont}{\normalfont}
\renewcommand{\cftchapleader}{\cftdotfill{\cftdotsep}}
\setlength{\cftlabindent}{3em}
\setlength{\cftlabnumwidth}{0em}
\setlength{\cftbeforelabskip}{0.10em}
\renewcommand{\cftlabfont}{\bfseries\color{UTPOrange!90!black}}
\renewcommand{\cftlabpagefont}{\bfseries\color{UTPOrange!90!black}}
\renewcommand{\cftlableader}{\color{UTPOrange!65!UTPMuted}\cftdotfill{\cftlabdotsep}}

% A list of figures, for a book that prints one (\listoffigures).
\renewcommand{\cftfigfont}{\normalfont\small}
\renewcommand{\cftfigpagefont}{\normalfont\small}
\renewcommand{\cftfigpresnum}{Figure~}
\renewcommand{\cftfigaftersnum}{\quad}
\setlength{\cftfigindent}{1.5em}
\setlength{\cftfignumwidth}{5.6em}
\setlength{\cftbeforefigskip}{0pt}
\renewcommand{\cftfigleader}{\cftdotfill{\cftdotsep}}

% The contents page itself, printed by \utpfrontmatter (section 14).
\newcommand{\utpcontentspage}[5]{%
  \def\utptocsize{#1}\def\utptocpartsize{#2}%
  \def\utptocpartskip{#3}\def\utptocchapterskip{#4}\def\utptoclabskip{#5}}
\utpcontentspage{\normalsize}{\large}{1.15em}{0.38em}{0.30em}
\newcommand{\utptableofcontents}{%
  \begingroup
  \utptocsize
  \setlength{\parskip}{0pt}%
  \setlength{\cftbeforepartskip}{\utptocpartskip}%
  \setlength{\cftbeforechapskip}{\utptocchapterskip}%
  \setlength{\cftbeforelabskip}{\utptoclabskip}%
  \renewcommand{\cftpartfont}{\utptocpartsize\bfseries\color{UTPNavy}}%
  \renewcommand{\cftpartpagefont}{\utptocpartsize\bfseries\color{UTPNavy}}%
  \tableofcontents
  \endgroup
}

% --------------------------------------------------------------------
%  6.  HEADINGS
% --------------------------------------------------------------------
%  The chapter badge (a coloured block with a word and the number),
%  part pages, and section spacing.
%
%  The word in the badge is \utpchapterlabel; Data Communication sets
%  it to LAB because each of its chapters is a lab.

\titleformat{\chapter}[display]
  {\bfseries\color{UTPNavy}}
  {\filleft\colorbox{UTPBlue}{\parbox{2.55cm}{\centering\color{white}\Large
    \utpchapterlabel\\[-0.15em]\Huge\thechapter}}}
  {0.65em}{\Huge\raggedright}
  [\vspace{0.25em}{\color{UTPBlue}\titlerule[1.2pt]}]
\renewcommand{\thepart}{Stage~\arabic{part}}
\titleformat{\part}[display]
  {\centering\color{UTPNavy}\bfseries}
  {\Large\color{UTPTeal}\thepart}{0.45em}{\Huge}
  [\vspace{0.55em}\color{UTPTeal}\rule{0.45\linewidth}{1.5pt}]
\titleformat{\section}{\Large\bfseries\color{UTPNavy}}{\thesection}{0.55em}{}
\titleformat{\subsection}{\large\bfseries\color{UTPBlue}}{\thesubsection}{0.5em}{}
\titlespacing*{\chapter}{0pt}{-1.8em}{1.0em}
\titlespacing*{\part}{0pt}{0pt}{1em}
\titlespacing*{\section}{0pt}{1.2em}{0.35em}
\titlespacing*{\subsection}{0pt}{0.95em}{0.2em}

% --------------------------------------------------------------------
%  7.  PAGE FURNITURE
% --------------------------------------------------------------------
%  Running heads and footers. The left header is \utprunninghead;
%  the right one follows the current chapter.

\pagestyle{fancy}
\fancyhf{}
\lhead{\small\color{UTPMuted}\utprunninghead}
\rhead{\small\color{UTPMuted}\nouppercase{\leftmark}}
\cfoot{\small\color{UTPMuted}\thepage}
\renewcommand{\headrulewidth}{0.5pt}
\renewcommand{\headrule}{\hbox to\headwidth{\color{UTPPaleBlue}\leaders\hrule height \headrulewidth\hfill}}
\renewcommand{\chaptermark}[1]{\markboth{#1}{}}
\fancypagestyle{plain}{%
  \fancyhf{}%
  \cfoot{\small\color{UTPMuted}\thepage}%
  \renewcommand{\headrulewidth}{0pt}%
}

% --------------------------------------------------------------------
%  8.  LEARNING BOXES
% --------------------------------------------------------------------
%  The coloured callouts the chapters use. Every title is a macro, so
%  a book can rename a box without restating its geometry:
%
%      \def\utptitlepitfall{Technical note}   % before the input
%
%  To change how a box LOOKS in one book, use one of the named
%  variations at the end of this section, or \renewtcolorbox after the
%  input. To add a box only one book needs, declare it in that book
%  with \utpnewcallout or \utpnewquietbox.

% Box titles -- override any of these before the input line.
\providecommand{\utptitlekeyidea}{Key idea}
\providecommand{\utptitleworked}{Worked example}
\providecommand{\utptitlepitfall}{Common mistake}
\providecommand{\utptitletryit}{Try it yourself}
\providecommand{\utptitlequickcheck}{Check your understanding}
\providecommand{\utptitlerecap}{Chapter recap}
\providecommand{\utptitledeliverable}{Lab deliverable}
\providecommand{\utptitlelabmode}{How much is scaffolded?}
\providecommand{\utptitleintuition}{Plain idea}
\providecommand{\utptitlefirstread}{Core path}
\providecommand{\utptitleglance}{Chapter at a glance}
\providecommand{\utptitledebug}{Debugging task}
\providecommand{\utptitleoptional}{Beyond the core}
\providecommand{\utptitleverification}{Verification checklist}
\providecommand{\utptitletroubleshoot}{Troubleshooting}
\providecommand{\utptitlesourcebasis}{Source basis}

% Explaining an idea
\newtcolorbox{keyidea}{enhanced,breakable,colback=UTPPaleTeal,colframe=UTPTeal, boxrule=0.8pt,arc=2mm,left=2.5mm,right=2.5mm,top=1.8mm,bottom=1.8mm, fonttitle=\bfseries\color{UTPNavy},title=\utptitlekeyidea}
\newtcolorbox{workedexample}{enhanced,breakable,colback=UTPPaleBlue,colframe=UTPBlue, boxrule=0.8pt,arc=2mm,left=2.5mm,right=2.5mm,top=1.8mm,bottom=1.8mm, fonttitle=\bfseries\color{UTPNavy},title=\utptitleworked}
\newtcolorbox{pitfall}{enhanced,breakable,colback=UTPPaleRed,colframe=UTPRed, boxrule=0.8pt,arc=2mm,left=2.5mm,right=2.5mm,top=1.8mm,bottom=1.8mm, fonttitle=\bfseries\color{UTPNavy},title=\utptitlepitfall}
% Asking the reader to do something
\newtcolorbox{tryit}{enhanced,breakable,colback=white,colframe=UTPTeal, boxrule=0.9pt,arc=2mm,left=2.5mm,right=2.5mm,top=1.8mm,bottom=1.8mm, fonttitle=\bfseries\color{white},colbacktitle=UTPTeal,title=\utptitletryit}
\newtcolorbox{quickcheck}{enhanced,breakable,colback=white,colframe=UTPOrange, boxrule=0.9pt,arc=2mm,left=2.5mm,right=2.5mm,top=1.8mm,bottom=1.8mm, fonttitle=\bfseries\color{white},colbacktitle=UTPOrange,title=\utptitlequickcheck}
\newtcolorbox{chapterrecap}{enhanced,breakable,colback=UTPPaleTeal,colframe=UTPGreen, boxrule=0.9pt,arc=2mm,left=2.5mm,right=2.5mm,top=1.8mm,bottom=1.8mm, fonttitle=\bfseries\color{white},colbacktitle=UTPGreen,title=\utptitlerecap}
% Lab structure
\newtcolorbox{labdeliverable}{enhanced,breakable,colback=white,colframe=UTPGreen, boxrule=0.9pt,arc=2mm,left=2.6mm,right=2.6mm,top=1.9mm,bottom=1.9mm, fonttitle=\bfseries\color{white},colbacktitle=UTPGreen,title=\utptitledeliverable}
\newtcolorbox{labmode}{enhanced,breakable,colback=UTPPaleBlue!52,colframe=UTPBlue, boxrule=0.75pt,arc=2mm,left=2.6mm,right=2.6mm,top=1.7mm,bottom=1.7mm, fonttitle=\bfseries\color{white},colbacktitle=UTPBlue,title=\utptitlelabmode}
% Softer, lighter-weight asides
\newtcolorbox{intuition}{enhanced,breakable,colback=white,colframe=UTPOrange!42, boxrule=0.5pt,arc=2mm,left=2.5mm,right=2.5mm,top=1.5mm,bottom=1.5mm, fonttitle=\bfseries\color{UTPOrange!85!black},title=\utptitleintuition}
\newtcolorbox{firstread}{enhanced,breakable,colback=white,colframe=UTPBlue!42, boxrule=0.5pt,arc=2mm,left=2.5mm,right=2.5mm,top=1.5mm,bottom=1.5mm, fonttitle=\bfseries\color{UTPBlue}, title=\utptitlefirstread}
\newtcolorbox{debugtask}{enhanced,breakable,colback=UTPPaleRed!58,colframe=UTPRed, boxrule=0.85pt,arc=2mm,left=2.6mm,right=2.6mm,top=1.9mm,bottom=1.9mm, fonttitle=\bfseries\color{white},colbacktitle=UTPRed,title=\utptitledebug}
\newtcolorbox{analogy}{enhanced,breakable,colback=white,colframe=UTPOrange!38, boxrule=0.5pt,arc=2mm,left=2.5mm,right=2.5mm,top=1.5mm,bottom=1.5mm, fonttitle=\bfseries\color{UTPOrange!85!black},title=Analogy}
\newtcolorbox{coursechoice}{enhanced,breakable,colback=white,colframe=UTPNavy, boxrule=0.8pt,arc=2mm,left=2.8mm,right=2.8mm,top=2mm,bottom=2mm, fonttitle=\bfseries\color{white},colbacktitle=UTPNavy,title=Course choice}
\newtcolorbox{decisionmemo}{enhanced,breakable,colback=UTPPaleOrange!55,colframe=UTPOrange, boxrule=0.85pt,arc=2mm,left=2.6mm,right=2.6mm,top=1.9mm,bottom=1.9mm, fonttitle=\bfseries\color{white},colbacktitle=UTPOrange,title=Decision memo}
\newtcolorbox{libnote}{enhanced,breakable,colback=white,colframe=UTPMuted!48, boxrule=0.5pt,arc=2mm,left=2.5mm,right=2.5mm,top=1.5mm,bottom=1.5mm, fonttitle=\bfseries\color{UTPNavy},title=Reference note}
\newtcolorbox{miniproject}{enhanced,breakable,colback=UTPPaleTeal!52,colframe=UTPGreen, boxrule=0.9pt,arc=2mm,left=2.6mm,right=2.6mm,top=1.9mm,bottom=1.9mm, fonttitle=\bfseries\color{white},colbacktitle=UTPGreen,title=Mini-project}
% Where the material came from: the lecturer's slides, manual or tutorial
% that a chapter is built on.
\newtcolorbox{sourcebasis}{enhanced,breakable,colback=white,colframe=UTPMuted!55, boxrule=0.5pt,arc=2mm,left=2.5mm,right=2.5mm,top=1.5mm,bottom=1.5mm, fonttitle=\bfseries\color{UTPMuted},title=\utptitlesourcebasis}
% Boxes that take a title argument
\newtcolorbox{modelcard}[1]{enhanced,breakable,colback=white,colframe=UTPBlue, boxrule=0.9pt,arc=2mm,left=2.5mm,right=2.5mm,top=1.8mm,bottom=1.8mm, fonttitle=\bfseries\color{white},colbacktitle=UTPBlue,title=#1}
\newtcolorbox{optionaltopic}{enhanced,breakable,colback=white,colframe=UTPMuted!58, boxrule=0.5pt,arc=2mm,left=2.5mm,right=2.5mm,top=1.5mm,bottom=1.5mm, fonttitle=\bfseries\color{UTPMuted},title=\utptitleoptional}
\newtcolorbox{secondpass}{enhanced,breakable,colback=UTPPaleGray!82,colframe=UTPMuted, boxrule=0.75pt,arc=2mm,left=2.6mm,right=2.6mm,top=1.9mm,bottom=1.9mm, fonttitle=\bfseries\color{white},colbacktitle=UTPMuted,title=Second pass}
\newtcolorbox{labproject}[1]{enhanced,breakable,colback=white,colframe=UTPBlue, boxrule=0.9pt,arc=2mm,left=2.6mm,right=2.6mm,top=1.9mm,bottom=1.9mm, fonttitle=\bfseries\color{white},colbacktitle=UTPBlue,title={#1}}
% Evidence that a lab worked, and what to do when it did not
\newtcolorbox{verification}{enhanced,breakable,colback=UTPPaleTeal!65,colframe=UTPGreen, boxrule=0.8pt,arc=2mm,left=2.6mm,right=2.6mm,top=1.8mm,bottom=1.8mm, fonttitle=\bfseries\color{white},colbacktitle=UTPGreen,title=\utptitleverification}
\newtcolorbox{troubleshoot}{enhanced,breakable,colback=UTPPaleOrange!70,colframe=UTPOrange, boxrule=0.8pt,arc=2mm,left=2.6mm,right=2.6mm,top=1.8mm,bottom=1.8mm, fonttitle=\bfseries\color{white},colbacktitle=UTPOrange,title=\utptitletroubleshoot}

% Named variations a book asks for after the input. Each restyles boxes
% that already exist, so the chapters do not change.
%
%   \utpheavycallouts          a solid colour bar behind the titles of
%                              keyidea, workedexample, intuition and
%                              firstread, with a firmer frame.
%   \utpbreakableframes        let utpfigure and utptable split over a
%                              page break, in a lighter blue (section 11).
%   \utplabprojectcolour{...}  the labproject box in another colour, with
%                              a little more room inside.
\newcommand{\utplabprojectcolour}[1]{%
  \renewtcolorbox{labproject}[1]{enhanced,breakable,colback=white,colframe=#1, boxrule=0.9pt,arc=2mm,left=2.7mm,right=2.7mm,top=2mm,bottom=2mm, fonttitle=\bfseries\color{white},colbacktitle=#1,title={##1}}%
}
\newcommand{\utpheavycallouts}{%
  \renewtcolorbox{keyidea}{enhanced,breakable,colback=UTPPaleTeal,colframe=UTPTeal, boxrule=0.85pt,arc=2mm,left=2.5mm,right=2.5mm,top=1.8mm,bottom=1.8mm, fonttitle=\bfseries\color{white},colbacktitle=UTPTeal,title=\utptitlekeyidea}%
  \renewtcolorbox{workedexample}{enhanced,breakable,colback=UTPPaleBlue,colframe=UTPBlue, boxrule=0.8pt,arc=2mm,left=2.5mm,right=2.5mm,top=1.8mm,bottom=1.8mm, fonttitle=\bfseries\color{white},colbacktitle=UTPBlue,title=\utptitleworked}%
  \renewtcolorbox{intuition}{enhanced,breakable,colback=UTPPaleOrange,colframe=UTPOrange, boxrule=0.8pt,arc=2mm,left=2.5mm,right=2.5mm,top=1.8mm,bottom=1.8mm, fonttitle=\bfseries\color{white},colbacktitle=UTPOrange,title=\utptitleintuition}%
  \renewtcolorbox{firstread}{enhanced,breakable,colback=white,colframe=UTPBlue, boxrule=0.9pt,arc=2mm,left=2.5mm,right=2.5mm,top=1.8mm,bottom=1.8mm, fonttitle=\bfseries\color{white},colbacktitle=UTPBlue,title=\utptitlefirstread}%
}

% Declaring a box a single book needs. Both constructors take the house
% geometry and let a trailing key list override any of it, so a new box
% is one line and still matches the others:
%
%   \utpnewcallout{odoobox}{white}{UTPBlue}{Odoo 19 Community}
%   \utpnewcallout[boxrule=0.9pt]{name}{back}{frame}{Title}
%   \utpnewquietbox{referencebox}{UTPMuted!58}{UTPMuted}{Reference note}
%
% A callout is a tinted box under a solid title bar, for something the
% reader has to act on. A quiet box is a light frame with coloured
% title text, for an aside they can read past. If a second book ever
% wants the same box, move its line into this section.
\newcommand{\utpnewcallout}[5][]{%
  \newtcolorbox{#2}{enhanced,breakable,colback=#3,colframe=#4, boxrule=0.8pt,arc=2mm,left=2.6mm,right=2.6mm,top=1.8mm,bottom=1.8mm, fonttitle=\bfseries\color{white},colbacktitle=#4,title={#5},#1}%
}
\newcommand{\utpnewquietbox}[5][]{%
  \newtcolorbox{#2}{enhanced,breakable,colback=white,colframe=#3, boxrule=0.5pt,arc=2mm,left=2.5mm,right=2.5mm,top=1.5mm,bottom=1.5mm, fonttitle=\bfseries\color{#4},title={#5},#1}%
}

% --------------------------------------------------------------------
%  9.  SMALL COMMANDS
% --------------------------------------------------------------------
%  \term for a defined term, \cmd for inline code, \chapref and
%  \secref for cross-references, and the utpsteps numbered-step list.
%
%      \term{overfitting}          a term being defined
%      \cmd{git status}            inline code, in any language
%      \chapref{ch:trees}          "Chapter 12"
%      \secref{sec:pruning}        "Section 12.3"
%      \chapterguide{prior}{outcomes}   the box that opens a chapter
%      \begin{utpsteps} ... \end{utpsteps}  a numbered Step 1/Step 2 list
%
%  The second group is the shorthand the chapters use in running text:
%
%      \figref{fig:tree}           "Figure 4.2"
%      \chapsref{ch:a}{ch:b}       "Chapters 3--5"
%      \lib{pandas}                a library name
%      \cxg{x}, \cxo{x}, \cxr{x}   maths coloured green/orange/red
%      \given                      the bar in P(a \given b)
%      \coretag, \latertag         read-now / come-back-to markers
%      \kagglesource{url}{name}    the dataset a lab is built on

\newcommand{\term}[1]{\textbf{\color{UTPNavy}#1}}
\newcommand{\cmd}[1]{\texttt{#1}}
\newcommand{\chapref}[1]{Chapter~\ref{#1}}
\newcommand{\secref}[1]{Section~\ref{#1}}
\newcommand{\chapterguide}[2]{%
  \begin{tcolorbox}[enhanced,breakable,colback=white,colframe=UTPNavy,boxrule=0.8pt,
    arc=2mm,left=2.5mm,right=2.5mm,top=2mm,bottom=2mm,
    title=\utptitleglance,fonttitle=\bfseries\color{white},colbacktitle=UTPNavy]
    \textbf{Before you start:} #1\par
    \textbf{By the end, you can:} #2
  \end{tcolorbox}}
\newenvironment{utpsteps}
  {\begin{enumerate}[label=\textcolor{UTPBlue}{\bfseries Step~\arabic*:},leftmargin=4.5em,itemsep=0.38em]}
  {\end{enumerate}}

\newcommand{\figref}[1]{Figure~\ref{#1}}
\newcommand{\chapsref}[2]{Chapters~\ref{#1}--\ref{#2}}
\newcommand{\lib}[1]{\textbf{\color{UTPBlue}\texttt{#1}}}
\newcommand{\cxg}[1]{{\color{UTPGreen}\bfseries $#1$}}
\newcommand{\cxo}[1]{{\color{UTPOrange}\bfseries $#1$}}
\newcommand{\cxr}[1]{{\color{UTPRed}\bfseries $#1$}}
\newcommand{\given}{\mathrel{|}}
\DeclareMathOperator*{\argmin}{arg\,min}
\DeclareMathOperator*{\argmax}{arg\,max}
\newcommand{\coretag}{\par\smallskip\noindent\colorbox{UTPPaleBlue}{\strut\small\bfseries\color{UTPBlue}CORE NOW}\par\smallskip}
\newcommand{\latertag}{\par\smallskip\noindent\colorbox{UTPPaleGray}{\strut\small\bfseries\color{UTPMuted}RETURN LATER}\par\smallskip}
\newcommand{\kagglesource}[2]{%
  \par\smallskip\noindent\textbf{Kaggle source:} \href{#1}{#2}\par
}

% --------------------------------------------------------------------
%  10.  DIAGRAM NODE STYLES
% --------------------------------------------------------------------
%  The vocabulary the hand-drawn TikZ diagrams are built from: a cell
%  in an array, the one cell under the cursor, a circular graph node,
%  the arrow between two of them. Every book has the whole set:
%
%      utpcell, utphot, utpdone, utpdim, utpwall        array cells
%      utpnode, utpnodehot, utpnodedone, utpnodedim     graph nodes
%      utparrow, utpptr, utpdash, utpedge               connectors
%      utplabel, utptag, utpbrace                       annotation
%      utpstep, utpstage, utppanel, utpbox              boxes
%
%  A cell is 1.0cm by 0.8cm. A book whose diagrams were drawn to a
%  different grid resizes the cells once, after the input:
%
%      \utpdiagramsize{1.05cm}{0.85cm}
%
%  Larger vocabularies for one kind of drawing live in .shared/diagrams/
%  and are loaded by name, after the input:
%
%      \utpusediagrams{network}         devices, links, packet headers,
%                                       sequence diagrams, state machines
%      \utpusediagrams{process}         business-process flows, documents,
%                                       swimlanes, status chips
%      \utpusediagrams{data-objects}    vectors, lists and data frames
%                                       drawn as cells; flowcharts
%      \utpusediagrams{network,process} several at once

\newcommand{\utpdiagramsize}[2]{%
  \tikzset{
    utpcell/.style={draw=UTPBlue,fill=UTPPaleBlue,minimum width=#1,minimum height=#2,
                    font=\small\bfseries\color{UTPNavy},inner sep=1pt},
    utphot/.style={draw=UTPOrange,fill=UTPPaleOrange,minimum width=#1,minimum height=#2,
                   font=\small\bfseries\color{UTPNavy},line width=1pt,inner sep=1pt},
    utpdone/.style={draw=UTPGreen,fill=UTPPaleTeal,minimum width=#1,minimum height=#2,
                    font=\small\bfseries\color{UTPNavy},inner sep=1pt},
    utpdim/.style={draw=UTPMuted!45,fill=UTPPaleGray,minimum width=#1,minimum height=#2,
                   font=\small\color{UTPMuted},inner sep=1pt},
    utpwall/.style={draw=UTPNavy,fill=UTPNavy!82,minimum width=#1,minimum height=#2,
                    inner sep=1pt},
    utpnode/.style={circle,draw=UTPBlue,fill=UTPPaleBlue,minimum size=8mm,
                    font=\small\bfseries\color{UTPNavy},inner sep=0.5pt},
    utpnodehot/.style={circle,draw=UTPOrange,fill=UTPPaleOrange,minimum size=8mm,line width=1pt,
                       font=\small\bfseries\color{UTPNavy},inner sep=0.5pt},
    utpnodedone/.style={circle,draw=UTPGreen,fill=UTPPaleTeal,minimum size=8mm,
                        font=\small\bfseries\color{UTPNavy},inner sep=0.5pt},
    utpnodedim/.style={circle,draw=UTPMuted!45,fill=UTPPaleGray,minimum size=8mm,
                       font=\small\color{UTPMuted},inner sep=0.5pt},
    utparrow/.style={-{Stealth[length=2.2mm]},thick,draw=UTPTeal},
    utpptr/.style={-{Stealth[length=2mm]},semithick,draw=UTPBlue},
    utpdash/.style={-{Stealth[length=2mm]},semithick,draw=UTPRed,dashed},
    utpedge/.style={thick,draw=UTPTeal},
    utplabel/.style={font=\scriptsize\color{UTPMuted},inner sep=1pt},
    utptag/.style={font=\scriptsize\bfseries\color{UTPOrange},inner sep=1pt},
    utpbrace/.style={decorate,decoration={brace,amplitude=4pt,raise=2pt},draw=UTPMuted,semithick},
    utpstep/.style={draw=UTPNavy,fill=white,rounded corners=1.5mm,minimum height=0.75cm,
                    font=\scriptsize\color{UTPNavy},align=center,inner sep=2pt},
    utpstage/.style={draw=UTPNavy,fill=white,rounded corners=1.5mm,minimum height=0.75cm,
                     font=\scriptsize\color{UTPNavy},align=center,inner sep=2pt},
    utppanel/.style={draw=UTPBlue!30,fill=white,rounded corners=1.5mm,inner sep=2.5mm},
    utpbox/.style={draw=UTPBlue,fill=UTPPaleBlue,rounded corners=2mm,align=center,
                   font=\small\color{UTPNavy},inner sep=2.2mm}
  }%
}
\utpdiagramsize{1.0cm}{0.8cm}

% \@for, not \foreach: \foreach runs each pass in a group, which would
% throw away the styles a library declares.
\makeatletter
\newcommand{\utpusediagrams}[1]{%
  \@for\utp@library:=#1\do{\input{../.shared/diagrams/\utp@library}}%
}
\makeatother

% --------------------------------------------------------------------
%  11.  FIGURE FRAMES
% --------------------------------------------------------------------
%  Three ways to put a drawing, chart or screenshot on the page:
%
%      \begin{utpfigure} ... \end{utpfigure}
%          a centred frame with no caption, for a diagram the text
%          talks about directly. An optional title puts a bar on top:
%          \begin{utpfigure}[Where the VLAN tag is added] ...
%
%      \begin{utptable} ... \end{utptable}
%          the same frame in teal, around a tabular or tabularx.
%
%      \begin{utpfig}[short entry]{The caption.}{fig:label}
%        \begin{tikzpicture} ... \end{tikzpicture}
%      \end{utpfig}
%          a frame with a numbered caption that \figref can point to.
%          The optional first argument is the short form for a list of
%          figures. A page never breaks between the drawing and its
%          caption.
%
%  A screenshot or slide lifted from the lecturer's material goes in a
%  titled frame, so the reader can tell it from a drawing made for the
%  notes; the optional width defaults to 0.9 and 0.78 of the line:
%
%      \sourceimage[0.8\linewidth]{Static-routing topology}{static-topology.png}
%      \compactsourceimage{The same, kept to a smaller height}{static-topology.png}
%
%  \fignote{...} adds a short muted line under a drawing, inside its frame.

\newtcolorbox{utpfigure}[1][]{enhanced,width=\linewidth,colback=white,colframe=UTPBlue!28, boxrule=0.7pt,arc=2mm, left=2.2mm,right=2.2mm,top=2.2mm,bottom=1.8mm, before skip=0.75em,after skip=0.85em, before upper={\centering}, halign=center, fonttitle=\bfseries\color{white},colbacktitle=UTPBlue,title={#1}}
\newtcolorbox{utptable}{enhanced,width=\linewidth,colback=white,colframe=UTPTeal!32, boxrule=0.7pt,arc=2mm, left=2.2mm,right=2.2mm,top=1.8mm,bottom=1.8mm, before skip=0.75em,after skip=0.85em, halign=center}
\newcommand{\utpbreakableframes}{%
  \renewtcolorbox{utpfigure}[1][]{enhanced,breakable,colback=white,colframe=UTPBlue!35, boxrule=0.5pt,arc=2mm,left=2.5mm,right=2.5mm,top=2.2mm,bottom=2.2mm, halign=center, fonttitle=\bfseries\color{white},colbacktitle=UTPBlue,title={##1}}%
  \renewtcolorbox{utptable}{enhanced,breakable,colback=white,colframe=UTPBlue!28, boxrule=0.45pt,arc=2mm,left=2.2mm,right=2.2mm,top=1.8mm,bottom=1.8mm}%
}

\captionsetup{
  hypcap=false,
  font={small},
  labelfont={bf,color=UTPNavy},
  textfont={color=UTPInk},
  labelsep=period,
  justification=centering,
  skip=4pt
}

\newenvironment{utpfig}[3][\relax]
  {\def\utpfigshort{#1}\def\utpfigcap{#2}\def\utpfiglab{#3}%
   \par\addvspace{0.55em}%
   \begin{tcolorbox}[enhanced,breakable,colback=white,colframe=UTPBlue!32,
     boxrule=0.5pt,arc=2mm,left=2.5mm,right=2.5mm,top=2.6mm,bottom=2.2mm]
   \centering}
  {\par\nopagebreak\vspace{0.45em}%
   % The caption is boxed: \captionof starts with a write and a zero
   % rule, and the glue after them would otherwise be a place where the
   % page could break, leaving the caption alone at the top of the next.
   % Inside the box the caption is told it sits below the drawing, which
   % it can no longer see, so the gap stays above it.
   \nointerlineskip
   \vbox{\setlength{\parskip}{0pt}\captionsetup{position=below}%
   \ifx\utpfigshort\relax
     \captionof{figure}{\utpfigcap}%
   \else
     \captionof{figure}[\utpfigshort]{\utpfigcap}%
   \fi
   \label{\utpfiglab}}%
   \end{tcolorbox}%
   \par\addvspace{0.55em}}

\newcommand{\sourceimage}[3][0.9\linewidth]{%
  \begin{utpfigure}[Source visual --- #2]%
    \includegraphics[width=#1,height=0.56\textheight,keepaspectratio]{#3}%
  \end{utpfigure}}
\newcommand{\compactsourceimage}[3][0.78\linewidth]{%
  \begin{utpfigure}[Source visual --- #2]%
    \includegraphics[width=#1,height=0.42\textheight,keepaspectratio]{#3}%
  \end{utpfigure}}

\newcommand{\fignote}[1]{\par\smallskip{\scriptsize\color{UTPMuted}#1\par}}

% --------------------------------------------------------------------
%  12.  LAB SHEETS
% --------------------------------------------------------------------
%  Every lab sheet opens with the banner, a full-width page break that
%  also puts the lab in the contents:
%
%      \labsection{3}{Scheduling Policies}{sec:lab03}
%
%  The "of N" comes from \utplabtotal, so a book states its lab count
%  once in its configuration rather than in every lab file.
%
%  The body of the sheet is a labproject box (section 8), and inside it
%  the same toolkit serves every course:
%
%      \labflow{chapter link}{goal}{what to explain}
%          the three-line brief under the lab title.
%      \labsourcecard{file}{what it does}{what to inspect}
%          the heading over a listing of one supplied source file.
%      \begin{labquestion}{3}{Short title} ... \end{labquestion}
%      \begin{labanswer} ... \expectedoutput ... \end{labanswer}
%          a question as the student meets it, then the worked code, the
%          output it should print, and what that output means.
%      \begin{checkanswers} ... \end{checkanswers}
%          answers to a quickcheck, set apart so they can be skipped.
%      labmode, verification, troubleshoot, labdeliverable (section 8)
%          how much is scaffolded, the evidence it worked, what to do if
%          it did not, and what to hand in.

\newcommand{\labsection}[3]{%
  \clearpage
  \phantomsection
  \addcontentsline{toc}{lab}{Lab #1 of \utplabtotal\space--- #2}%
  \markboth{Lab #1 of \utplabtotal\space--- #2}{}%
  \begin{tcolorbox}[enhanced,width=\linewidth,colback=UTPNavy,colframe=UTPNavy,
    boxrule=0pt,arc=2mm,left=5mm,right=5mm,top=5mm,bottom=5mm,
    before skip=0pt,after skip=1.2em]
    {\large\bfseries\color{UTPTeal!35!white} PRACTICAL MILESTONE\par}
    \vspace{0.25em}
    {\Huge\bfseries\color{white} Lab #1 of \utplabtotal\par}
    \vspace{0.20em}
    {\Large\bfseries\color{white} #2\par}
  \end{tcolorbox}
  \label{#3}%
}

\newcommand{\labflow}[3]{%
  \par\smallskip\noindent
  \textbf{Chapter connection.} #1\par
  \textbf{Goal.} #2\par
  \textbf{Explain.} #3\par\medskip}
\newcommand{\labsourcecard}[3]{%
  \par\medskip\noindent\textbf{#1 --- #2}\par
  \textbf{Inspect.} #3\par\smallskip}

\newtcolorbox{labquestion}[2]{enhanced,breakable,colback=UTPPaleOrange!55,colframe=UTPOrange,
  boxrule=0.8pt,arc=2mm,left=2.6mm,right=2.6mm,top=1.8mm,bottom=1.8mm,
  fonttitle=\bfseries\color{white},colbacktitle=UTPOrange,title={Question #1 --- #2}}
\newtcolorbox{labanswer}{enhanced,breakable,colback=white,colframe=UTPTeal,
  boxrule=0.8pt,arc=2mm,left=2.6mm,right=2.6mm,top=1.8mm,bottom=1.8mm,
  fonttitle=\bfseries\color{white},colbacktitle=UTPTeal,title=Worked answer}
\newcommand{\expectedoutput}{\par\smallskip\noindent\textbf{\color{UTPNavy}Reference output.}\nopagebreak\par\nopagebreak}
\utpnewcallout[boxrule=0.65pt,left=2.5mm,right=2.5mm,top=1.6mm,bottom=1.6mm,colbacktitle=UTPMuted]{checkanswers}{UTPPaleGray!72}{UTPMuted!62}{Quick-check answers}

% --------------------------------------------------------------------
%  13.  STAGE DIVIDER
% --------------------------------------------------------------------
%  A full-page divider used by the longer books to open a part:
%  \substage{3A}{Supervised Workflow}{one-line summary}.
%  Harmless to leave unused.

\newcommand{\substage}[3]{%
  \cleardoublepage
  \thispagestyle{empty}
  \phantomsection
  \addcontentsline{toc}{chapter}{Stage #1 --- #2}
  \begin{center}
    \vspace*{0.20\textheight}
    {\Large\bfseries\color{UTPTeal} Stage #1\par}
    \vspace{0.45em}
    {\Huge\bfseries\color{UTPNavy} #2\par}
    \vspace{0.65em}
    {\color{UTPTeal}\rule{0.46\linewidth}{1.5pt}\par}
    \vspace{1.1em}
    \begin{minipage}{0.76\linewidth}
      \centering\large\color{UTPMuted} #3
    \end{minipage}
  \end{center}
  \cleardoublepage
}

% --------------------------------------------------------------------
%  14.  FRONT MATTER
% --------------------------------------------------------------------
%  The opening every book shares: the cover, the contents in roman
%  numerals, the course map, then chapter one on an arabic page 1.
%
%      \begin{document}
%      \utpfrontmatter
%
%  What it prints comes from the configuration in section 1 --- the
%  cover text from \utpcoursetitle and \utpcourseoutline, the overview
%  chapter from \utpcoursemap --- and the contents page is set as
%  \utpcontentspage (section 5) describes. A book with no overview
%  chapter sets \def\utpcoursemap{} and the step is skipped.
%
%  Page anchors are off across the cover so hyperref does not warn
%  about the unnumbered pages, and back on before the contents links
%  are used.
%
%  The number of front-matter pages is recorded in the .aux file as
%  \utpfrontpages, which is how .shared/build.py turns the page numbers
%  in the contents into physical pages when it splits the book.

\newcommand{\utpfrontmatter}{%
  \hypersetup{pageanchor=false}%
  \input{../.shared/cover}%
  \cleardoublepage
  \pagenumbering{roman}%
  \setcounter{tocdepth}{0}%
  \utptableofcontents
  \newpage
  \hypersetup{pageanchor=true}%
  \ifx\utpcoursemap\empty\else\input{\utpcoursemap}\fi
  \cleardoublepage
  \ifdefined\ReadonlyShipoutCounter
    \immediate\write\csname @auxout\endcsname
      {\string\gdef\string\utpfrontpages{\the\ReadonlyShipoutCounter}}%
  \fi
  \pagenumbering{arabic}%
}

% --------------------------------------------------------------------
%  15.  LISTING STYLES
% --------------------------------------------------------------------
%  Two base looks, and one named style per language built on them, so
%  the colours and frame are the same in every book:
%
%      utpcode          the base for source code
%      utpterminal      the base for commands typed and output printed
%      utppython        Python
%      utpr             R
%      utpc             C
%      utpcpp           C++
%      utpbash          shell scripts and commands
%      utpcnumbered     C with line numbers, for a complete source file
%      utpbashnumbered  shell with line numbers, likewise
%
%      \begin{lstlisting}[style=utppython] ... \end{lstlisting}
%      \lstinputlisting[style=utpcnumbered]{all_experiments/Lab04/fcfs.c}
%
%  A book that sets most of its listings in one style makes it the
%  default once, after the input:  \lstset{style=utpcpp}

\lstdefinestyle{utpcode}{
  basicstyle=\ttfamily\scriptsize,
  keywordstyle=\bfseries\color{UTPBlue},
  commentstyle=\itshape\color{UTPMuted},
  stringstyle=\color{UTPGreen!70!black},
  backgroundcolor=\color{UTPPaleGray!72},
  frame=single,
  rulecolor=\color{UTPBlue!28},
  framesep=5pt,
  breaklines=true,
  breakatwhitespace=false,
  showstringspaces=false,
  keepspaces=true,
  columns=fullflexible,
  upquote=true,
  tabsize=4,
  aboveskip=0.7em,
  belowskip=0.8em
}
\lstdefinestyle{utpterminal}{
  basicstyle=\ttfamily\scriptsize,
  backgroundcolor=\color{UTPPaleGray!72},
  frame=single,
  rulecolor=\color{UTPMuted!35},
  framesep=5pt,
  breaklines=true,
  showstringspaces=false,
  keepspaces=true,
  columns=fullflexible,
  upquote=true,
  aboveskip=0.6em,
  belowskip=0.7em
}

\lstdefinestyle{utppython}{style=utpcode,language=Python}
\lstdefinestyle{utpr}{style=utpcode,language=R,tabsize=2}
\lstdefinestyle{utpc}{style=utpcode,language=C}
\lstdefinestyle{utpcpp}{style=utpcode,language=C++}
\lstdefinestyle{utpbash}{style=utpcode,language=bash,
  rulecolor=\color{UTPMuted!35},aboveskip=0.6em,belowskip=0.7em}
\lstdefinestyle{utpcnumbered}{style=utpc,numbers=left,numberstyle=\tiny\color{UTPMuted},
  numbersep=7pt,xleftmargin=1.7em,framexleftmargin=1.4em}
\lstdefinestyle{utpbashnumbered}{style=utpbash,numbers=left,numberstyle=\tiny\color{UTPMuted},
  numbersep=7pt,xleftmargin=1.7em,framexleftmargin=1.4em}

% Pseudocode. A book that sets algorithms asks for them once, after the
% input, and gets algorithm2e loaded and coloured to match:
%
%     \utpusealgorithms
%
% \algostyle then sets one algorithm environment a size smaller than the
% text around it.
\newcommand{\utpalgorithmstyle}{%
  \SetAlCapNameFnt{\bfseries\color{UTPNavy}}%
  \SetAlCapFnt{\bfseries\color{UTPNavy}}%
  \SetAlgoCaptionSeparator{.}%
  \SetKwInput{KwIn}{Input}%
  \SetKwInput{KwOut}{Output}%
  \SetNlSty{color}{UTPMuted}{}%
  \SetCommentSty{\itshape\color{UTPMuted}}%
  \SetKwSty{bfseries\color{UTPBlue}}%
  \SetArgSty{}%
  \providecommand{\algostyle}{\small\SetAlgoNlRelativeSize{-1}\DontPrintSemicolon\SetAlgoVlined}%
}
\newcommand{\utpusealgorithms}{%
  \usepackage[linesnumbered,ruled,vlined]{algorithm2e}%
  \utpalgorithmstyle
}

%
%      \utpcontentspage{\small}{\normalsize}{0.65em}{0.07em}{0.16em}
%
%      \begin{document}
%      \utpfrontmatter                     % cover, contents, map
%      \part{Orientation}
%      \input{chapters/01-where-to-begin}
%      \end{document}
%
%  Configuration goes BEFORE the \input; overrides go AFTER it.
%  Paths are relative, so compile from inside the course folder.
%
%  ------------------------------------------------------------------
%  NAMING
%  ------------------------------------------------------------------
%  One name per idea, used by every book. Anything a book can use is
%  either an unprefixed learning box (keyidea, pitfall, labproject,
%  ...) or carries the utp prefix (utpfigure, utptable, utpfig,
%  utpsteps, utppython, UTPNavy, utpcell, ...). Course prefixes such
%  as ml-, ds- or erp- are not used: a chapter moved from one book to
%  another compiles unchanged. .shared/README.md lists every name.
%
%  Adding a new course? Copy .shared/TEMPLATE-main.tex and follow
%  .shared/ADDING-A-COURSE.md. Every step is a numbered checklist.
%
%  ------------------------------------------------------------------
%  CONTENTS
%  ------------------------------------------------------------------
%    1  Course configuration     the knobs a book may set
%    2  Packages                 the shared package set and fonts
%    3  Colour palette           the one UTP palette
%    4  Page layout and tables   margins, lists, table styling
%    5  Table of contents        parts, chapters, labs, figures
%    6  Headings                 chapter badge, part pages, sections
%    7  Page furniture           running heads and footers
%    8  Learning boxes           key idea, pitfall, try it, and friends
%    9  Small commands           \term, \cmd, \chapterguide, utpsteps
%   10  Diagram node styles      the TikZ vocabulary every book shares
%   11  Figure frames            utpfigure, utpfig, source images
%   12  Lab sheets               the lab banner and the lab toolkit
%   13  Stage divider            the part-opening page in longer books
%   14  Front matter             \utpfrontmatter: cover, contents, map
%   15  Listing styles           code, terminal and pseudocode
% =====================================================================

% --------------------------------------------------------------------
%  1.  COURSE CONFIGURATION
% --------------------------------------------------------------------
%  The knobs below, plus \utpstyleparts in section 5 and the box titles
%  in section 8, are the whole configuration surface. A book sets what
%  it needs BEFORE the \input and inherits the rest.
%
%  \providecommand means "define unless the book already did", which is
%  what makes the \def-before-\input pattern work.
%
%  Running text
%    \utprunninghead   the left running head on every page.
%    \utpchapterlabel  the word in the chapter badge. CHAPTER by
%                      default; a book whose chapters are labs sets LAB.
%    \utplabtotal      the N in the "Lab 3 of N" banners.

\providecommand{\utprunninghead}{Course Notes}
\providecommand{\utpchapterlabel}{CHAPTER}
\providecommand{\utplabtotal}{10}

%  Cover, printed by section 14
%    \utpcoursetitle       the title on the title page.
%    \utpdocumentsubtitle  the line under the title; Teaching Assistant Guide by default.
%    \utpcourseoutline     the sentence or two below it; \\ breaks a line.
%    \utpoutlineheading    the small heading above that sentence.
%    \utpauthor            the name in the PREPARED BY block.
%    \utpauthorrole        the line under the name.

\providecommand{\utpcoursetitle}{\utprunninghead}
\providecommand{\utpdocumentsubtitle}{Teaching Assistant Guide}
\providecommand{\utpcourseoutline}{}
\providecommand{\utpoutlineheading}{COURSE OUTLINE}
\providecommand{\utpauthor}{Reeyan Afzal Khan}
\providecommand{\utpauthorrole}{Graduate Research Assistant}

%  Files and folders
%    \utpextragraphicspath  image folders to search before the standard
%                           ones, as brace groups: {{screenshots/}}.
%                           A course's own figures/ folder and the
%                           shared logos are always searched.
%    \utpcoursemap          the overview chapter \utpfrontmatter inputs
%                           after the contents. Empty means none.

\providecommand{\utpextragraphicspath}{}
\providecommand{\utpcoursemap}{chapters/course-map}

%  Contents entries for chapters (section 5 applies them). A book whose
%  chapters are labs says so once here rather than restating the three
%  tocloft settings after the input.
%    \utpchaptertocword      Chapter~ by default, Lab~ in one book.
%    \utpchaptertocindent    left indent of the entry.
%    \utpchaptertocnumwidth  room reserved for the word and number.

\providecommand{\utpchaptertocword}{Chapter~}
\providecommand{\utpchaptertocindent}{1.5em}
\providecommand{\utpchaptertocnumwidth}{6em}

%  Monospace face (section 2)
%    \utpmonofont   sourcecodepro by default. inconsolata is the
%                   alternative, for a book whose pages are mostly
%                   device transcripts.

\providecommand{\utpmonofont}{sourcecodepro}

% --------------------------------------------------------------------
%  2.  PACKAGES
% --------------------------------------------------------------------
%  The union of what the books need. A course that needs something
%  extra loads it in its own file, under the heading "packages only
%  this book needs". Pseudocode is the one optional package the shared
%  file knows about: \utpusealgorithms in section 15 loads and styles it.
%
%  The tikz and pgfplots library lists are the union across all books,
%  so a diagram or chart copied from one course compiles in another.

\usepackage[T1]{fontenc}
\IfFileExists{glyphtounicode.tex}{\input{glyphtounicode}\pdfgentounicode=1}{}
\usepackage[american]{babel}
\IfFileExists{sourcesanspro.sty}
  {\usepackage[default]{sourcesanspro}}
  {\usepackage[scaled]{helvet}\renewcommand{\familydefault}{\sfdefault}}
\DeclareFontShape{T1}{cmss}{b}{n}{<->ssub*cmss/bx/n}{}
\DeclareFontShape{TS1}{cmss}{b}{n}{<->ssub*cmss/bx/n}{}
% Monospace. Without this the tt default stays Computer Modern Typewriter,
% a serif face that clashes with the sans body text in every listing and
% every \cmd. Source Code Pro is the companion to Source Sans Pro; it is
% scaled down because mono digits otherwise sit large next to the sans.
%
% Inconsolata, chosen with \def\utpmonofont{inconsolata}, distinguishes
% 0/O and 1/l/I at a glance, which matters when a reader copies an
% interface name or an address out of a transcript. Only one mono package
% is ever loaded: a second would win the family but inherit the first's
% series names, and every \ttfamily would then resolve by substitution.
\def\utpmonoinconsolata{inconsolata}
\ifx\utpmonofont\utpmonoinconsolata
  \IfFileExists{inconsolata.sty}{\usepackage[varqu,varl,scaled=1.02]{inconsolata}}{}
\else
  \IfFileExists{sourcecodepro.sty}
    {\usepackage[scaled=0.88]{sourcecodepro}}
    {\usepackage[scaled=0.86]{beramono}}
  % Source Sans names its normal weight "regular", but the Computer Modern
  % fallbacks call it "m". Map one onto the other so a \ttfamily reached
  % while the regular series is active resolves instead of warning.
  \DeclareFontShape{T1}{cmtt}{regular}{n}{<->ssub*cmtt/m/n}{}
  \DeclareFontShape{TS1}{cmtt}{regular}{n}{<->ssub*cmtt/m/n}{}
  \DeclareFontShape{T1}{cmtt}{regular}{it}{<->ssub*cmtt/m/it}{}
\fi
\usepackage[margin=21mm,headheight=22pt]{geometry}
\usepackage{amsmath,amssymb,mathtools,bm}
\usepackage{microtype}
\usepackage{enumitem}
\usepackage{booktabs}
\usepackage{tabularx}
\usepackage{array}
\usepackage{ragged2e}
\usepackage{etoolbox}
\usepackage{titlesec}
\usepackage{tocloft}
\usepackage{fancyhdr}
\usepackage[table]{xcolor}
\usepackage{tcolorbox}
\tcbuselibrary{breakable,skins}
\usepackage[hidelinks]{hyperref}
\usepackage{bookmark}
% Document properties, so the PDF identifies itself in a reader's
% title bar, in a library catalogue, and in search results. Set from
% the same configuration macros the cover uses.
\hypersetup{
  pdftitle={\utpcoursetitle},
  pdfauthor={\utpauthor},
  pdfsubject={Universiti Teknologi PETRONAS course notes and laboratory guide},
  pdfkeywords={UTP, course notes, laboratory manual, \utpcoursetitle},
  pdfdisplaydoctitle=true,
  bookmarksnumbered=true,
  bookmarksopenlevel=1
}
\usepackage{listings}
\usepackage{graphicx}
\usepackage{caption}
\usepackage{tikz}
\usetikzlibrary{angles,arrows.meta,backgrounds,calc,chains,decorations.pathmorphing,decorations.pathreplacing,fit,matrix,patterns,positioning,quotes,shadows.blur,shapes.geometric,shapes.misc,shapes.symbols}
\usepackage{pgfplots}
% statistics gives boxplot prepared; groupplots gives multi-panel charts.
\usepgfplotslibrary{statistics}
\usepgfplotslibrary{groupplots}

% Images: any folders the book names first, then its own figures/
% folder, then the shared logos used by the cover.
\expandafter\graphicspath\expandafter{\utpextragraphicspath{figures/}{../.shared/logos/}}

% --------------------------------------------------------------------
%  3.  COLOUR PALETTE
% --------------------------------------------------------------------
%  One palette, defined once under UTP* names, and used by every book,
%  box, diagram and chart. Five accents (blue, teal, green, orange, red)
%  with a pale tint of each for box backgrounds, plus navy for headings,
%  gold for the cover, ink and muted for text, and a sixth accent,
%  purple, for charts and layered diagrams that need one more series.

\definecolor{UTPNavy}{HTML}{18324A}
\definecolor{UTPBlue}{HTML}{246B9E}
\definecolor{UTPTeal}{HTML}{138A80}
\definecolor{UTPGreen}{HTML}{2D7D46}
\definecolor{UTPOrange}{HTML}{C86B24}
\definecolor{UTPRed}{HTML}{B34843}
\definecolor{UTPPurple}{HTML}{6B4FA0}
\definecolor{UTPGold}{HTML}{D4A13E}
\definecolor{UTPInk}{HTML}{263238}
\definecolor{UTPMuted}{HTML}{60717C}
\definecolor{UTPPaleBlue}{HTML}{EEF6FB}
\definecolor{UTPPaleTeal}{HTML}{EEF8F6}
\definecolor{UTPPaleOrange}{HTML}{FFF5EA}
\definecolor{UTPPaleRed}{HTML}{FFF0EF}
\definecolor{UTPPalePurple}{HTML}{F0ECF9}
\definecolor{UTPPaleGray}{HTML}{F4F6F7}
\color{UTPInk}

% --------------------------------------------------------------------
%  4.  PAGE LAYOUT AND TABLES
% --------------------------------------------------------------------
%  Paragraph spacing, list indents, and the table look: alternating
%  row tint, top-aligned paragraph columns, ragged-right cells.
%
%  Two extra column types for tabularx tables:
%      D{width}   a fixed-width, ragged-right, vertically centred column
%      Y          a ragged-right column that fills the remaining width
%
%  Every table in every book has the same shape, the one Algorithms &
%  Data Structures set: a utptable frame around a full-width tabularx,
%  booktabs rules, a tinted header row of \utpth cells (centred, bold,
%  navy), and a first column in bold navy when it names the rows.
%  Cells are vertically centred, so use m{} rather than p{} columns.
%
%      \begin{utptable}
%      \begin{tabularx}{\linewidth}{@{}>{\bfseries\leavevmode\color{UTPNavy}}m{0.22\linewidth}X@{}}
%      \toprule
%      \rowcolor{UTPPaleBlue!92}
%      \utpth{Operation} & \utpth{What it does} \\
%      \midrule
%      Insertion & Add a new element. \\
%      \bottomrule
%      \end{tabularx}
%      \end{utptable}
%
%  A book whose tables carry dense text asks for \utpcompacttables after
%  the input: tabularx tables then set in \footnotesize with tighter
%  padding and vertically centred cells.

\setlength{\parindent}{0pt}
\setlength{\parskip}{0.55em}
\setlist[itemize]{leftmargin=1.45em,itemsep=0.22em,topsep=0.25em}
\setlist[enumerate]{leftmargin=1.75em,itemsep=0.28em,topsep=0.3em}
\setlist[description]{font=\bfseries\color{UTPNavy},itemsep=0.25em,topsep=0.3em}
\renewcommand{\arraystretch}{1.34}
\setlength{\tabcolsep}{6.5pt}
\arrayrulecolor{UTPBlue!45}
% Keep paragraph-style columns ragged and vertically centred. Tables use
% fixed m columns beside flexible X columns; giving X the same centring
% keeps a short cell level with a tall neighbour in a multi-line row.
\renewcommand{\tabularxcolumn}[1]{m{#1}}
\makeatletter
\g@addto@macro\@arrayparboxrestore{\RaggedRight\arraybackslash}
\makeatother
% A header cell: centred, bold and navy, on the tinted header row.
\newcommand{\utpth}[1]{\multicolumn{1}{c}{\bfseries\color{UTPNavy}#1}}
\AtBeginEnvironment{tabular}{\small\setlength{\extrarowheight}{0.6pt}\rowcolors{2}{UTPPaleBlue!72}{white}}
\AtBeginEnvironment{tabularx}{\small\setlength{\extrarowheight}{0.6pt}\rowcolors{2}{UTPPaleBlue!72}{white}}
\AtBeginEnvironment{figure}{\centering}
\newcolumntype{D}[1]{>{\raggedright\arraybackslash}m{#1}}
\newcolumntype{Y}{>{\raggedright\arraybackslash}X}
\newcommand{\utpcompacttables}{%
  \renewcommand{\tabularxcolumn}[1]{m{##1}}%
  \AtBeginEnvironment{tabularx}{%
    \footnotesize
    \setlength{\tabcolsep}{4pt}%
    \renewcommand{\arraystretch}{1.25}%
    \setlength{\extrarowheight}{0.6pt}%
  }%
}
\emergencystretch=1.5em
\hbadness=10000
\hfuzz=8pt

% Page-breaking discipline. Without these a paragraph can leave a single
% line stranded at the foot of one page or the head of the next, which is
% the most visible sign of unattended typesetting.
\widowpenalty=10000          % no first line of a paragraph left alone at a page foot
\clubpenalty=10000           % no last line of a paragraph left alone at a page head
\displaywidowpenalty=10000   % the same around displayed equations
\brokenpenalty=10000         % do not break a page after a hyphenated line
\predisplaypenalty=10000     % keep the line before a display with it
% Pages are built from boxes of many heights, so let the last line float
% up rather than stretching the gaps between boxes to force a flush foot.
\raggedbottom
% Prefer a slightly loose line to an overfull one, but do not hyphenate
% freely -- short sans lines look ragged when over-hyphenated.
\tolerance=1200
\hyphenpenalty=200

% Character protrusion and font expansion. Both are what stop a justified
% sans column from showing visible rivers down the page.
\microtypesetup{protrusion=true,expansion=true,tracking=true}

\pgfplotsset{
  compat=1.18,
  every axis/.append style={
    axis line style={draw=UTPMuted!85,line width=0.55pt},
    tick style={draw=UTPMuted!85},
    tick label style={font=\scriptsize\color{UTPMuted}},
    label style={font=\small\color{UTPNavy}},
    title style={font=\small\bfseries\color{UTPNavy}},
    legend style={font=\scriptsize,draw=UTPBlue!22,fill=white,rounded corners=1pt}
  },
  every axis plot/.append style={line width=1.15pt}
}
\tikzset{every picture/.append style={line cap=round,line join=round}}

% --------------------------------------------------------------------
%  5.  TABLE OF CONTENTS
% --------------------------------------------------------------------
%  Formatting for part, chapter and lab entries, and for a list of
%  figures if a book prints one. The lab level is a custom tocloft list
%  so lab banners appear in the contents.
%
%  Part styling is skipped when \utpstyleparts is 0, which is what a
%  book without \part divisions should set.
%
%  The contents page is meant to fill exactly one page, and how loosely
%  it is set depends on how many entries a book has. One line after the
%  input tunes it:
%
%      \utpcontentspage{<base size>}{<part size>}%
%                      {<space before a part>}{<before a chapter>}{<before a lab>}
%      \utpcontentspage{\normalsize}{\large}{1.15em}{0.38em}{0.30em}   % the default

\addto\captionsamerican{\renewcommand{\contentsname}{Learning Path}}
\newlistentry{lab}{toc}{-1}
\providecommand{\utpstyleparts}{1}
\ifnum\utpstyleparts=1
\renewcommand{\cftpartpresnum}{}
\renewcommand{\cftpartaftersnum}{}
\setlength{\cftpartnumwidth}{5.2em}
\setlength{\cftbeforepartskip}{0.38em}
\renewcommand{\cftpartfont}{\normalsize\bfseries\color{UTPNavy}}
\renewcommand{\cftpartpagefont}{\normalsize\bfseries\color{UTPNavy}}
\renewcommand{\cftpartleader}{\hfill}
\fi
\renewcommand{\cftchappresnum}{\utpchaptertocword}
\renewcommand{\cftchapaftersnum}{\quad}
\setlength{\cftchapindent}{\utpchaptertocindent}
\setlength{\cftchapnumwidth}{\utpchaptertocnumwidth}
\setlength{\cftbeforechapskip}{0pt}
\renewcommand{\cftchapfont}{\normalfont}
\renewcommand{\cftchappagefont}{\normalfont}
\renewcommand{\cftchapleader}{\cftdotfill{\cftdotsep}}
\setlength{\cftlabindent}{3em}
\setlength{\cftlabnumwidth}{0em}
\setlength{\cftbeforelabskip}{0.10em}
\renewcommand{\cftlabfont}{\bfseries\color{UTPOrange!90!black}}
\renewcommand{\cftlabpagefont}{\bfseries\color{UTPOrange!90!black}}
\renewcommand{\cftlableader}{\color{UTPOrange!65!UTPMuted}\cftdotfill{\cftlabdotsep}}

% A list of figures, for a book that prints one (\listoffigures).
\renewcommand{\cftfigfont}{\normalfont\small}
\renewcommand{\cftfigpagefont}{\normalfont\small}
\renewcommand{\cftfigpresnum}{Figure~}
\renewcommand{\cftfigaftersnum}{\quad}
\setlength{\cftfigindent}{1.5em}
\setlength{\cftfignumwidth}{5.6em}
\setlength{\cftbeforefigskip}{0pt}
\renewcommand{\cftfigleader}{\cftdotfill{\cftdotsep}}

% The contents page itself, printed by \utpfrontmatter (section 14).
\newcommand{\utpcontentspage}[5]{%
  \def\utptocsize{#1}\def\utptocpartsize{#2}%
  \def\utptocpartskip{#3}\def\utptocchapterskip{#4}\def\utptoclabskip{#5}}
\utpcontentspage{\normalsize}{\large}{1.15em}{0.38em}{0.30em}
\newcommand{\utptableofcontents}{%
  \begingroup
  \utptocsize
  \setlength{\parskip}{0pt}%
  \setlength{\cftbeforepartskip}{\utptocpartskip}%
  \setlength{\cftbeforechapskip}{\utptocchapterskip}%
  \setlength{\cftbeforelabskip}{\utptoclabskip}%
  \renewcommand{\cftpartfont}{\utptocpartsize\bfseries\color{UTPNavy}}%
  \renewcommand{\cftpartpagefont}{\utptocpartsize\bfseries\color{UTPNavy}}%
  \tableofcontents
  \endgroup
}

% --------------------------------------------------------------------
%  6.  HEADINGS
% --------------------------------------------------------------------
%  The chapter badge (a coloured block with a word and the number),
%  part pages, and section spacing.
%
%  The word in the badge is \utpchapterlabel; Data Communication sets
%  it to LAB because each of its chapters is a lab.

\titleformat{\chapter}[display]
  {\bfseries\color{UTPNavy}}
  {\filleft\colorbox{UTPBlue}{\parbox{2.55cm}{\centering\color{white}\Large
    \utpchapterlabel\\[-0.15em]\Huge\thechapter}}}
  {0.65em}{\Huge\raggedright}
  [\vspace{0.25em}{\color{UTPBlue}\titlerule[1.2pt]}]
\renewcommand{\thepart}{Stage~\arabic{part}}
\titleformat{\part}[display]
  {\centering\color{UTPNavy}\bfseries}
  {\Large\color{UTPTeal}\thepart}{0.45em}{\Huge}
  [\vspace{0.55em}\color{UTPTeal}\rule{0.45\linewidth}{1.5pt}]
\titleformat{\section}{\Large\bfseries\color{UTPNavy}}{\thesection}{0.55em}{}
\titleformat{\subsection}{\large\bfseries\color{UTPBlue}}{\thesubsection}{0.5em}{}
\titlespacing*{\chapter}{0pt}{-1.8em}{1.0em}
\titlespacing*{\part}{0pt}{0pt}{1em}
\titlespacing*{\section}{0pt}{1.2em}{0.35em}
\titlespacing*{\subsection}{0pt}{0.95em}{0.2em}

% --------------------------------------------------------------------
%  7.  PAGE FURNITURE
% --------------------------------------------------------------------
%  Running heads and footers. The left header is \utprunninghead;
%  the right one follows the current chapter.

\pagestyle{fancy}
\fancyhf{}
\lhead{\small\color{UTPMuted}\utprunninghead}
\rhead{\small\color{UTPMuted}\nouppercase{\leftmark}}
\cfoot{\small\color{UTPMuted}\thepage}
\renewcommand{\headrulewidth}{0.5pt}
\renewcommand{\headrule}{\hbox to\headwidth{\color{UTPPaleBlue}\leaders\hrule height \headrulewidth\hfill}}
\renewcommand{\chaptermark}[1]{\markboth{#1}{}}
\fancypagestyle{plain}{%
  \fancyhf{}%
  \cfoot{\small\color{UTPMuted}\thepage}%
  \renewcommand{\headrulewidth}{0pt}%
}

% --------------------------------------------------------------------
%  8.  LEARNING BOXES
% --------------------------------------------------------------------
%  The coloured callouts the chapters use. Every title is a macro, so
%  a book can rename a box without restating its geometry:
%
%      \def\utptitlepitfall{Technical note}   % before the input
%
%  To change how a box LOOKS in one book, use one of the named
%  variations at the end of this section, or \renewtcolorbox after the
%  input. To add a box only one book needs, declare it in that book
%  with \utpnewcallout or \utpnewquietbox.

% Box titles -- override any of these before the input line.
\providecommand{\utptitlekeyidea}{Key idea}
\providecommand{\utptitleworked}{Worked example}
\providecommand{\utptitlepitfall}{Common mistake}
\providecommand{\utptitletryit}{Try it yourself}
\providecommand{\utptitlequickcheck}{Check your understanding}
\providecommand{\utptitlerecap}{Chapter recap}
\providecommand{\utptitledeliverable}{Lab deliverable}
\providecommand{\utptitlelabmode}{How much is scaffolded?}
\providecommand{\utptitleintuition}{Plain idea}
\providecommand{\utptitlefirstread}{Core path}
\providecommand{\utptitleglance}{Chapter at a glance}
\providecommand{\utptitledebug}{Debugging task}
\providecommand{\utptitleoptional}{Beyond the core}
\providecommand{\utptitleverification}{Verification checklist}
\providecommand{\utptitletroubleshoot}{Troubleshooting}
\providecommand{\utptitlesourcebasis}{Source basis}

% Explaining an idea
\newtcolorbox{keyidea}{enhanced,breakable,colback=UTPPaleTeal,colframe=UTPTeal, boxrule=0.8pt,arc=2mm,left=2.5mm,right=2.5mm,top=1.8mm,bottom=1.8mm, fonttitle=\bfseries\color{UTPNavy},title=\utptitlekeyidea}
\newtcolorbox{workedexample}{enhanced,breakable,colback=UTPPaleBlue,colframe=UTPBlue, boxrule=0.8pt,arc=2mm,left=2.5mm,right=2.5mm,top=1.8mm,bottom=1.8mm, fonttitle=\bfseries\color{UTPNavy},title=\utptitleworked}
\newtcolorbox{pitfall}{enhanced,breakable,colback=UTPPaleRed,colframe=UTPRed, boxrule=0.8pt,arc=2mm,left=2.5mm,right=2.5mm,top=1.8mm,bottom=1.8mm, fonttitle=\bfseries\color{UTPNavy},title=\utptitlepitfall}
% Asking the reader to do something
\newtcolorbox{tryit}{enhanced,breakable,colback=white,colframe=UTPTeal, boxrule=0.9pt,arc=2mm,left=2.5mm,right=2.5mm,top=1.8mm,bottom=1.8mm, fonttitle=\bfseries\color{white},colbacktitle=UTPTeal,title=\utptitletryit}
\newtcolorbox{quickcheck}{enhanced,breakable,colback=white,colframe=UTPOrange, boxrule=0.9pt,arc=2mm,left=2.5mm,right=2.5mm,top=1.8mm,bottom=1.8mm, fonttitle=\bfseries\color{white},colbacktitle=UTPOrange,title=\utptitlequickcheck}
\newtcolorbox{chapterrecap}{enhanced,breakable,colback=UTPPaleTeal,colframe=UTPGreen, boxrule=0.9pt,arc=2mm,left=2.5mm,right=2.5mm,top=1.8mm,bottom=1.8mm, fonttitle=\bfseries\color{white},colbacktitle=UTPGreen,title=\utptitlerecap}
% Lab structure
\newtcolorbox{labdeliverable}{enhanced,breakable,colback=white,colframe=UTPGreen, boxrule=0.9pt,arc=2mm,left=2.6mm,right=2.6mm,top=1.9mm,bottom=1.9mm, fonttitle=\bfseries\color{white},colbacktitle=UTPGreen,title=\utptitledeliverable}
\newtcolorbox{labmode}{enhanced,breakable,colback=UTPPaleBlue!52,colframe=UTPBlue, boxrule=0.75pt,arc=2mm,left=2.6mm,right=2.6mm,top=1.7mm,bottom=1.7mm, fonttitle=\bfseries\color{white},colbacktitle=UTPBlue,title=\utptitlelabmode}
% Softer, lighter-weight asides
\newtcolorbox{intuition}{enhanced,breakable,colback=white,colframe=UTPOrange!42, boxrule=0.5pt,arc=2mm,left=2.5mm,right=2.5mm,top=1.5mm,bottom=1.5mm, fonttitle=\bfseries\color{UTPOrange!85!black},title=\utptitleintuition}
\newtcolorbox{firstread}{enhanced,breakable,colback=white,colframe=UTPBlue!42, boxrule=0.5pt,arc=2mm,left=2.5mm,right=2.5mm,top=1.5mm,bottom=1.5mm, fonttitle=\bfseries\color{UTPBlue}, title=\utptitlefirstread}
\newtcolorbox{debugtask}{enhanced,breakable,colback=UTPPaleRed!58,colframe=UTPRed, boxrule=0.85pt,arc=2mm,left=2.6mm,right=2.6mm,top=1.9mm,bottom=1.9mm, fonttitle=\bfseries\color{white},colbacktitle=UTPRed,title=\utptitledebug}
\newtcolorbox{analogy}{enhanced,breakable,colback=white,colframe=UTPOrange!38, boxrule=0.5pt,arc=2mm,left=2.5mm,right=2.5mm,top=1.5mm,bottom=1.5mm, fonttitle=\bfseries\color{UTPOrange!85!black},title=Analogy}
\newtcolorbox{coursechoice}{enhanced,breakable,colback=white,colframe=UTPNavy, boxrule=0.8pt,arc=2mm,left=2.8mm,right=2.8mm,top=2mm,bottom=2mm, fonttitle=\bfseries\color{white},colbacktitle=UTPNavy,title=Course choice}
\newtcolorbox{decisionmemo}{enhanced,breakable,colback=UTPPaleOrange!55,colframe=UTPOrange, boxrule=0.85pt,arc=2mm,left=2.6mm,right=2.6mm,top=1.9mm,bottom=1.9mm, fonttitle=\bfseries\color{white},colbacktitle=UTPOrange,title=Decision memo}
\newtcolorbox{libnote}{enhanced,breakable,colback=white,colframe=UTPMuted!48, boxrule=0.5pt,arc=2mm,left=2.5mm,right=2.5mm,top=1.5mm,bottom=1.5mm, fonttitle=\bfseries\color{UTPNavy},title=Reference note}
\newtcolorbox{miniproject}{enhanced,breakable,colback=UTPPaleTeal!52,colframe=UTPGreen, boxrule=0.9pt,arc=2mm,left=2.6mm,right=2.6mm,top=1.9mm,bottom=1.9mm, fonttitle=\bfseries\color{white},colbacktitle=UTPGreen,title=Mini-project}
% Where the material came from: the lecturer's slides, manual or tutorial
% that a chapter is built on.
\newtcolorbox{sourcebasis}{enhanced,breakable,colback=white,colframe=UTPMuted!55, boxrule=0.5pt,arc=2mm,left=2.5mm,right=2.5mm,top=1.5mm,bottom=1.5mm, fonttitle=\bfseries\color{UTPMuted},title=\utptitlesourcebasis}
% Boxes that take a title argument
\newtcolorbox{modelcard}[1]{enhanced,breakable,colback=white,colframe=UTPBlue, boxrule=0.9pt,arc=2mm,left=2.5mm,right=2.5mm,top=1.8mm,bottom=1.8mm, fonttitle=\bfseries\color{white},colbacktitle=UTPBlue,title=#1}
\newtcolorbox{optionaltopic}{enhanced,breakable,colback=white,colframe=UTPMuted!58, boxrule=0.5pt,arc=2mm,left=2.5mm,right=2.5mm,top=1.5mm,bottom=1.5mm, fonttitle=\bfseries\color{UTPMuted},title=\utptitleoptional}
\newtcolorbox{secondpass}{enhanced,breakable,colback=UTPPaleGray!82,colframe=UTPMuted, boxrule=0.75pt,arc=2mm,left=2.6mm,right=2.6mm,top=1.9mm,bottom=1.9mm, fonttitle=\bfseries\color{white},colbacktitle=UTPMuted,title=Second pass}
\newtcolorbox{labproject}[1]{enhanced,breakable,colback=white,colframe=UTPBlue, boxrule=0.9pt,arc=2mm,left=2.6mm,right=2.6mm,top=1.9mm,bottom=1.9mm, fonttitle=\bfseries\color{white},colbacktitle=UTPBlue,title={#1}}
% Evidence that a lab worked, and what to do when it did not
\newtcolorbox{verification}{enhanced,breakable,colback=UTPPaleTeal!65,colframe=UTPGreen, boxrule=0.8pt,arc=2mm,left=2.6mm,right=2.6mm,top=1.8mm,bottom=1.8mm, fonttitle=\bfseries\color{white},colbacktitle=UTPGreen,title=\utptitleverification}
\newtcolorbox{troubleshoot}{enhanced,breakable,colback=UTPPaleOrange!70,colframe=UTPOrange, boxrule=0.8pt,arc=2mm,left=2.6mm,right=2.6mm,top=1.8mm,bottom=1.8mm, fonttitle=\bfseries\color{white},colbacktitle=UTPOrange,title=\utptitletroubleshoot}

% Named variations a book asks for after the input. Each restyles boxes
% that already exist, so the chapters do not change.
%
%   \utpheavycallouts          a solid colour bar behind the titles of
%                              keyidea, workedexample, intuition and
%                              firstread, with a firmer frame.
%   \utpbreakableframes        let utpfigure and utptable split over a
%                              page break, in a lighter blue (section 11).
%   \utplabprojectcolour{...}  the labproject box in another colour, with
%                              a little more room inside.
\newcommand{\utplabprojectcolour}[1]{%
  \renewtcolorbox{labproject}[1]{enhanced,breakable,colback=white,colframe=#1, boxrule=0.9pt,arc=2mm,left=2.7mm,right=2.7mm,top=2mm,bottom=2mm, fonttitle=\bfseries\color{white},colbacktitle=#1,title={##1}}%
}
\newcommand{\utpheavycallouts}{%
  \renewtcolorbox{keyidea}{enhanced,breakable,colback=UTPPaleTeal,colframe=UTPTeal, boxrule=0.85pt,arc=2mm,left=2.5mm,right=2.5mm,top=1.8mm,bottom=1.8mm, fonttitle=\bfseries\color{white},colbacktitle=UTPTeal,title=\utptitlekeyidea}%
  \renewtcolorbox{workedexample}{enhanced,breakable,colback=UTPPaleBlue,colframe=UTPBlue, boxrule=0.8pt,arc=2mm,left=2.5mm,right=2.5mm,top=1.8mm,bottom=1.8mm, fonttitle=\bfseries\color{white},colbacktitle=UTPBlue,title=\utptitleworked}%
  \renewtcolorbox{intuition}{enhanced,breakable,colback=UTPPaleOrange,colframe=UTPOrange, boxrule=0.8pt,arc=2mm,left=2.5mm,right=2.5mm,top=1.8mm,bottom=1.8mm, fonttitle=\bfseries\color{white},colbacktitle=UTPOrange,title=\utptitleintuition}%
  \renewtcolorbox{firstread}{enhanced,breakable,colback=white,colframe=UTPBlue, boxrule=0.9pt,arc=2mm,left=2.5mm,right=2.5mm,top=1.8mm,bottom=1.8mm, fonttitle=\bfseries\color{white},colbacktitle=UTPBlue,title=\utptitlefirstread}%
}

% Declaring a box a single book needs. Both constructors take the house
% geometry and let a trailing key list override any of it, so a new box
% is one line and still matches the others:
%
%   \utpnewcallout{odoobox}{white}{UTPBlue}{Odoo 19 Community}
%   \utpnewcallout[boxrule=0.9pt]{name}{back}{frame}{Title}
%   \utpnewquietbox{referencebox}{UTPMuted!58}{UTPMuted}{Reference note}
%
% A callout is a tinted box under a solid title bar, for something the
% reader has to act on. A quiet box is a light frame with coloured
% title text, for an aside they can read past. If a second book ever
% wants the same box, move its line into this section.
\newcommand{\utpnewcallout}[5][]{%
  \newtcolorbox{#2}{enhanced,breakable,colback=#3,colframe=#4, boxrule=0.8pt,arc=2mm,left=2.6mm,right=2.6mm,top=1.8mm,bottom=1.8mm, fonttitle=\bfseries\color{white},colbacktitle=#4,title={#5},#1}%
}
\newcommand{\utpnewquietbox}[5][]{%
  \newtcolorbox{#2}{enhanced,breakable,colback=white,colframe=#3, boxrule=0.5pt,arc=2mm,left=2.5mm,right=2.5mm,top=1.5mm,bottom=1.5mm, fonttitle=\bfseries\color{#4},title={#5},#1}%
}

% --------------------------------------------------------------------
%  9.  SMALL COMMANDS
% --------------------------------------------------------------------
%  \term for a defined term, \cmd for inline code, \chapref and
%  \secref for cross-references, and the utpsteps numbered-step list.
%
%      \term{overfitting}          a term being defined
%      \cmd{git status}            inline code, in any language
%      \chapref{ch:trees}          "Chapter 12"
%      \secref{sec:pruning}        "Section 12.3"
%      \chapterguide{prior}{outcomes}   the box that opens a chapter
%      \begin{utpsteps} ... \end{utpsteps}  a numbered Step 1/Step 2 list
%
%  The second group is the shorthand the chapters use in running text:
%
%      \figref{fig:tree}           "Figure 4.2"
%      \chapsref{ch:a}{ch:b}       "Chapters 3--5"
%      \lib{pandas}                a library name
%      \cxg{x}, \cxo{x}, \cxr{x}   maths coloured green/orange/red
%      \given                      the bar in P(a \given b)
%      \coretag, \latertag         read-now / come-back-to markers
%      \kagglesource{url}{name}    the dataset a lab is built on

\newcommand{\term}[1]{\textbf{\color{UTPNavy}#1}}
\newcommand{\cmd}[1]{\texttt{#1}}
\newcommand{\chapref}[1]{Chapter~\ref{#1}}
\newcommand{\secref}[1]{Section~\ref{#1}}
\newcommand{\chapterguide}[2]{%
  \begin{tcolorbox}[enhanced,breakable,colback=white,colframe=UTPNavy,boxrule=0.8pt,
    arc=2mm,left=2.5mm,right=2.5mm,top=2mm,bottom=2mm,
    title=\utptitleglance,fonttitle=\bfseries\color{white},colbacktitle=UTPNavy]
    \textbf{Before you start:} #1\par
    \textbf{By the end, you can:} #2
  \end{tcolorbox}}
\newenvironment{utpsteps}
  {\begin{enumerate}[label=\textcolor{UTPBlue}{\bfseries Step~\arabic*:},leftmargin=4.5em,itemsep=0.38em]}
  {\end{enumerate}}

\newcommand{\figref}[1]{Figure~\ref{#1}}
\newcommand{\chapsref}[2]{Chapters~\ref{#1}--\ref{#2}}
\newcommand{\lib}[1]{\textbf{\color{UTPBlue}\texttt{#1}}}
\newcommand{\cxg}[1]{{\color{UTPGreen}\bfseries $#1$}}
\newcommand{\cxo}[1]{{\color{UTPOrange}\bfseries $#1$}}
\newcommand{\cxr}[1]{{\color{UTPRed}\bfseries $#1$}}
\newcommand{\given}{\mathrel{|}}
\DeclareMathOperator*{\argmin}{arg\,min}
\DeclareMathOperator*{\argmax}{arg\,max}
\newcommand{\coretag}{\par\smallskip\noindent\colorbox{UTPPaleBlue}{\strut\small\bfseries\color{UTPBlue}CORE NOW}\par\smallskip}
\newcommand{\latertag}{\par\smallskip\noindent\colorbox{UTPPaleGray}{\strut\small\bfseries\color{UTPMuted}RETURN LATER}\par\smallskip}
\newcommand{\kagglesource}[2]{%
  \par\smallskip\noindent\textbf{Kaggle source:} \href{#1}{#2}\par
}

% --------------------------------------------------------------------
%  10.  DIAGRAM NODE STYLES
% --------------------------------------------------------------------
%  The vocabulary the hand-drawn TikZ diagrams are built from: a cell
%  in an array, the one cell under the cursor, a circular graph node,
%  the arrow between two of them. Every book has the whole set:
%
%      utpcell, utphot, utpdone, utpdim, utpwall        array cells
%      utpnode, utpnodehot, utpnodedone, utpnodedim     graph nodes
%      utparrow, utpptr, utpdash, utpedge               connectors
%      utplabel, utptag, utpbrace                       annotation
%      utpstep, utpstage, utppanel, utpbox              boxes
%
%  A cell is 1.0cm by 0.8cm. A book whose diagrams were drawn to a
%  different grid resizes the cells once, after the input:
%
%      \utpdiagramsize{1.05cm}{0.85cm}
%
%  Larger vocabularies for one kind of drawing live in .shared/diagrams/
%  and are loaded by name, after the input:
%
%      \utpusediagrams{network}         devices, links, packet headers,
%                                       sequence diagrams, state machines
%      \utpusediagrams{process}         business-process flows, documents,
%                                       swimlanes, status chips
%      \utpusediagrams{data-objects}    vectors, lists and data frames
%                                       drawn as cells; flowcharts
%      \utpusediagrams{network,process} several at once

\newcommand{\utpdiagramsize}[2]{%
  \tikzset{
    utpcell/.style={draw=UTPBlue,fill=UTPPaleBlue,minimum width=#1,minimum height=#2,
                    font=\small\bfseries\color{UTPNavy},inner sep=1pt},
    utphot/.style={draw=UTPOrange,fill=UTPPaleOrange,minimum width=#1,minimum height=#2,
                   font=\small\bfseries\color{UTPNavy},line width=1pt,inner sep=1pt},
    utpdone/.style={draw=UTPGreen,fill=UTPPaleTeal,minimum width=#1,minimum height=#2,
                    font=\small\bfseries\color{UTPNavy},inner sep=1pt},
    utpdim/.style={draw=UTPMuted!45,fill=UTPPaleGray,minimum width=#1,minimum height=#2,
                   font=\small\color{UTPMuted},inner sep=1pt},
    utpwall/.style={draw=UTPNavy,fill=UTPNavy!82,minimum width=#1,minimum height=#2,
                    inner sep=1pt},
    utpnode/.style={circle,draw=UTPBlue,fill=UTPPaleBlue,minimum size=8mm,
                    font=\small\bfseries\color{UTPNavy},inner sep=0.5pt},
    utpnodehot/.style={circle,draw=UTPOrange,fill=UTPPaleOrange,minimum size=8mm,line width=1pt,
                       font=\small\bfseries\color{UTPNavy},inner sep=0.5pt},
    utpnodedone/.style={circle,draw=UTPGreen,fill=UTPPaleTeal,minimum size=8mm,
                        font=\small\bfseries\color{UTPNavy},inner sep=0.5pt},
    utpnodedim/.style={circle,draw=UTPMuted!45,fill=UTPPaleGray,minimum size=8mm,
                       font=\small\color{UTPMuted},inner sep=0.5pt},
    utparrow/.style={-{Stealth[length=2.2mm]},thick,draw=UTPTeal},
    utpptr/.style={-{Stealth[length=2mm]},semithick,draw=UTPBlue},
    utpdash/.style={-{Stealth[length=2mm]},semithick,draw=UTPRed,dashed},
    utpedge/.style={thick,draw=UTPTeal},
    utplabel/.style={font=\scriptsize\color{UTPMuted},inner sep=1pt},
    utptag/.style={font=\scriptsize\bfseries\color{UTPOrange},inner sep=1pt},
    utpbrace/.style={decorate,decoration={brace,amplitude=4pt,raise=2pt},draw=UTPMuted,semithick},
    utpstep/.style={draw=UTPNavy,fill=white,rounded corners=1.5mm,minimum height=0.75cm,
                    font=\scriptsize\color{UTPNavy},align=center,inner sep=2pt},
    utpstage/.style={draw=UTPNavy,fill=white,rounded corners=1.5mm,minimum height=0.75cm,
                     font=\scriptsize\color{UTPNavy},align=center,inner sep=2pt},
    utppanel/.style={draw=UTPBlue!30,fill=white,rounded corners=1.5mm,inner sep=2.5mm},
    utpbox/.style={draw=UTPBlue,fill=UTPPaleBlue,rounded corners=2mm,align=center,
                   font=\small\color{UTPNavy},inner sep=2.2mm}
  }%
}
\utpdiagramsize{1.0cm}{0.8cm}

% \@for, not \foreach: \foreach runs each pass in a group, which would
% throw away the styles a library declares.
\makeatletter
\newcommand{\utpusediagrams}[1]{%
  \@for\utp@library:=#1\do{\input{../.shared/diagrams/\utp@library}}%
}
\makeatother

% --------------------------------------------------------------------
%  11.  FIGURE FRAMES
% --------------------------------------------------------------------
%  Three ways to put a drawing, chart or screenshot on the page:
%
%      \begin{utpfigure} ... \end{utpfigure}
%          a centred frame with no caption, for a diagram the text
%          talks about directly. An optional title puts a bar on top:
%          \begin{utpfigure}[Where the VLAN tag is added] ...
%
%      \begin{utptable} ... \end{utptable}
%          the same frame in teal, around a tabular or tabularx.
%
%      \begin{utpfig}[short entry]{The caption.}{fig:label}
%        \begin{tikzpicture} ... \end{tikzpicture}
%      \end{utpfig}
%          a frame with a numbered caption that \figref can point to.
%          The optional first argument is the short form for a list of
%          figures. A page never breaks between the drawing and its
%          caption.
%
%  A screenshot or slide lifted from the lecturer's material goes in a
%  titled frame, so the reader can tell it from a drawing made for the
%  notes; the optional width defaults to 0.9 and 0.78 of the line:
%
%      \sourceimage[0.8\linewidth]{Static-routing topology}{static-topology.png}
%      \compactsourceimage{The same, kept to a smaller height}{static-topology.png}
%
%  \fignote{...} adds a short muted line under a drawing, inside its frame.

\newtcolorbox{utpfigure}[1][]{enhanced,width=\linewidth,colback=white,colframe=UTPBlue!28, boxrule=0.7pt,arc=2mm, left=2.2mm,right=2.2mm,top=2.2mm,bottom=1.8mm, before skip=0.75em,after skip=0.85em, before upper={\centering}, halign=center, fonttitle=\bfseries\color{white},colbacktitle=UTPBlue,title={#1}}
\newtcolorbox{utptable}{enhanced,width=\linewidth,colback=white,colframe=UTPTeal!32, boxrule=0.7pt,arc=2mm, left=2.2mm,right=2.2mm,top=1.8mm,bottom=1.8mm, before skip=0.75em,after skip=0.85em, halign=center}
\newcommand{\utpbreakableframes}{%
  \renewtcolorbox{utpfigure}[1][]{enhanced,breakable,colback=white,colframe=UTPBlue!35, boxrule=0.5pt,arc=2mm,left=2.5mm,right=2.5mm,top=2.2mm,bottom=2.2mm, halign=center, fonttitle=\bfseries\color{white},colbacktitle=UTPBlue,title={##1}}%
  \renewtcolorbox{utptable}{enhanced,breakable,colback=white,colframe=UTPBlue!28, boxrule=0.45pt,arc=2mm,left=2.2mm,right=2.2mm,top=1.8mm,bottom=1.8mm}%
}

\captionsetup{
  hypcap=false,
  font={small},
  labelfont={bf,color=UTPNavy},
  textfont={color=UTPInk},
  labelsep=period,
  justification=centering,
  skip=4pt
}

\newenvironment{utpfig}[3][\relax]
  {\def\utpfigshort{#1}\def\utpfigcap{#2}\def\utpfiglab{#3}%
   \par\addvspace{0.55em}%
   \begin{tcolorbox}[enhanced,breakable,colback=white,colframe=UTPBlue!32,
     boxrule=0.5pt,arc=2mm,left=2.5mm,right=2.5mm,top=2.6mm,bottom=2.2mm]
   \centering}
  {\par\nopagebreak\vspace{0.45em}%
   % The caption is boxed: \captionof starts with a write and a zero
   % rule, and the glue after them would otherwise be a place where the
   % page could break, leaving the caption alone at the top of the next.
   % Inside the box the caption is told it sits below the drawing, which
   % it can no longer see, so the gap stays above it.
   \nointerlineskip
   \vbox{\setlength{\parskip}{0pt}\captionsetup{position=below}%
   \ifx\utpfigshort\relax
     \captionof{figure}{\utpfigcap}%
   \else
     \captionof{figure}[\utpfigshort]{\utpfigcap}%
   \fi
   \label{\utpfiglab}}%
   \end{tcolorbox}%
   \par\addvspace{0.55em}}

\newcommand{\sourceimage}[3][0.9\linewidth]{%
  \begin{utpfigure}[Source visual --- #2]%
    \includegraphics[width=#1,height=0.56\textheight,keepaspectratio]{#3}%
  \end{utpfigure}}
\newcommand{\compactsourceimage}[3][0.78\linewidth]{%
  \begin{utpfigure}[Source visual --- #2]%
    \includegraphics[width=#1,height=0.42\textheight,keepaspectratio]{#3}%
  \end{utpfigure}}

\newcommand{\fignote}[1]{\par\smallskip{\scriptsize\color{UTPMuted}#1\par}}

% --------------------------------------------------------------------
%  12.  LAB SHEETS
% --------------------------------------------------------------------
%  Every lab sheet opens with the banner, a full-width page break that
%  also puts the lab in the contents:
%
%      \labsection{3}{Scheduling Policies}{sec:lab03}
%
%  The "of N" comes from \utplabtotal, so a book states its lab count
%  once in its configuration rather than in every lab file.
%
%  The body of the sheet is a labproject box (section 8), and inside it
%  the same toolkit serves every course:
%
%      \labflow{chapter link}{goal}{what to explain}
%          the three-line brief under the lab title.
%      \labsourcecard{file}{what it does}{what to inspect}
%          the heading over a listing of one supplied source file.
%      \begin{labquestion}{3}{Short title} ... \end{labquestion}
%      \begin{labanswer} ... \expectedoutput ... \end{labanswer}
%          a question as the student meets it, then the worked code, the
%          output it should print, and what that output means.
%      \begin{checkanswers} ... \end{checkanswers}
%          answers to a quickcheck, set apart so they can be skipped.
%      labmode, verification, troubleshoot, labdeliverable (section 8)
%          how much is scaffolded, the evidence it worked, what to do if
%          it did not, and what to hand in.

\newcommand{\labsection}[3]{%
  \clearpage
  \phantomsection
  \addcontentsline{toc}{lab}{Lab #1 of \utplabtotal\space--- #2}%
  \markboth{Lab #1 of \utplabtotal\space--- #2}{}%
  \begin{tcolorbox}[enhanced,width=\linewidth,colback=UTPNavy,colframe=UTPNavy,
    boxrule=0pt,arc=2mm,left=5mm,right=5mm,top=5mm,bottom=5mm,
    before skip=0pt,after skip=1.2em]
    {\large\bfseries\color{UTPTeal!35!white} PRACTICAL MILESTONE\par}
    \vspace{0.25em}
    {\Huge\bfseries\color{white} Lab #1 of \utplabtotal\par}
    \vspace{0.20em}
    {\Large\bfseries\color{white} #2\par}
  \end{tcolorbox}
  \label{#3}%
}

\newcommand{\labflow}[3]{%
  \par\smallskip\noindent
  \textbf{Chapter connection.} #1\par
  \textbf{Goal.} #2\par
  \textbf{Explain.} #3\par\medskip}
\newcommand{\labsourcecard}[3]{%
  \par\medskip\noindent\textbf{#1 --- #2}\par
  \textbf{Inspect.} #3\par\smallskip}

\newtcolorbox{labquestion}[2]{enhanced,breakable,colback=UTPPaleOrange!55,colframe=UTPOrange,
  boxrule=0.8pt,arc=2mm,left=2.6mm,right=2.6mm,top=1.8mm,bottom=1.8mm,
  fonttitle=\bfseries\color{white},colbacktitle=UTPOrange,title={Question #1 --- #2}}
\newtcolorbox{labanswer}{enhanced,breakable,colback=white,colframe=UTPTeal,
  boxrule=0.8pt,arc=2mm,left=2.6mm,right=2.6mm,top=1.8mm,bottom=1.8mm,
  fonttitle=\bfseries\color{white},colbacktitle=UTPTeal,title=Worked answer}
\newcommand{\expectedoutput}{\par\smallskip\noindent\textbf{\color{UTPNavy}Reference output.}\nopagebreak\par\nopagebreak}
\utpnewcallout[boxrule=0.65pt,left=2.5mm,right=2.5mm,top=1.6mm,bottom=1.6mm,colbacktitle=UTPMuted]{checkanswers}{UTPPaleGray!72}{UTPMuted!62}{Quick-check answers}

% --------------------------------------------------------------------
%  13.  STAGE DIVIDER
% --------------------------------------------------------------------
%  A full-page divider used by the longer books to open a part:
%  \substage{3A}{Supervised Workflow}{one-line summary}.
%  Harmless to leave unused.

\newcommand{\substage}[3]{%
  \cleardoublepage
  \thispagestyle{empty}
  \phantomsection
  \addcontentsline{toc}{chapter}{Stage #1 --- #2}
  \begin{center}
    \vspace*{0.20\textheight}
    {\Large\bfseries\color{UTPTeal} Stage #1\par}
    \vspace{0.45em}
    {\Huge\bfseries\color{UTPNavy} #2\par}
    \vspace{0.65em}
    {\color{UTPTeal}\rule{0.46\linewidth}{1.5pt}\par}
    \vspace{1.1em}
    \begin{minipage}{0.76\linewidth}
      \centering\large\color{UTPMuted} #3
    \end{minipage}
  \end{center}
  \cleardoublepage
}

% --------------------------------------------------------------------
%  14.  FRONT MATTER
% --------------------------------------------------------------------
%  The opening every book shares: the cover, the contents in roman
%  numerals, the course map, then chapter one on an arabic page 1.
%
%      \begin{document}
%      \utpfrontmatter
%
%  What it prints comes from the configuration in section 1 --- the
%  cover text from \utpcoursetitle and \utpcourseoutline, the overview
%  chapter from \utpcoursemap --- and the contents page is set as
%  \utpcontentspage (section 5) describes. A book with no overview
%  chapter sets \def\utpcoursemap{} and the step is skipped.
%
%  Page anchors are off across the cover so hyperref does not warn
%  about the unnumbered pages, and back on before the contents links
%  are used.
%
%  The number of front-matter pages is recorded in the .aux file as
%  \utpfrontpages, which is how .shared/build.py turns the page numbers
%  in the contents into physical pages when it splits the book.

\newcommand{\utpfrontmatter}{%
  \hypersetup{pageanchor=false}%
  % ---------- Cover ----------
%  The title page, printed by \utpfrontmatter. Everything that changes
%  between books is a configuration macro from style/configuration.tex:
%  \utpcoursetitle, \utpcourseoutline, \utpoutlineheading, \utpauthor
%  and \utpauthorrole. The colours are the shared palette.

\begin{titlepage}
\thispagestyle{empty}
\begin{tikzpicture}[remember picture,overlay]
  \fill[UTPGold]
    (current page.north west) rectangle
    ([yshift=-0.12cm]current page.north east);
  \fill[UTPGold]
    (current page.south west) rectangle
    ([yshift=0.12cm]current page.south east);
\end{tikzpicture}

\centering
\includegraphics[
  height=2.85cm,
  clip
]{logo-UTP.png}\par

\vspace{0.50cm}
{\fontsize{29}{33}\selectfont\bfseries\color{UTPNavy}\utpcoursetitle\par}
\vspace{0.24cm}
{\Large\color{UTPMuted} \utpdocumentsubtitle\par}
\vspace{0.42cm}
{\color{UTPGold}\rule{0.26\linewidth}{1.3pt}\par}
\vspace{0.72cm}
\begin{tikzpicture}[
  core/.style={circle,draw=UTPGold,fill=UTPNavy,minimum size=1.95cm,
    align=center,font=\small\bfseries\color{white},line width=1pt},
  stage/.style={rounded corners=1.5mm,draw=UTPBlue!65,fill=UTPPaleBlue!42,
    text width=3.15cm,minimum height=1.42cm,align=center,
    font=\small\bfseries\color{UTPNavy},inner xsep=3mm,
    inner ysep=3.5mm,line width=0.65pt},
  cycle/.style={-{Stealth[length=2.1mm]},semithick,draw=UTPMuted,
    rounded corners=4mm,shorten <=7pt,shorten >=7pt}
]
\node[core] (guide) at (0,-2.55) {TA\\WORKFLOW};
\node[stage] (prepare) at (0,0) {SESSION\\PLANNING};
\node[stage] (explain) at (4.75,-2.55) {INSTRUCTIONAL\\SUPPORT};
\node[stage] (facilitate) at (0,-5.10) {STUDENT\\ENGAGEMENT};
\node[stage] (assess) at (-4.75,-2.55) {LEARNING\\ASSESSMENT};

\draw[cycle] (prepare.east) -- (explain.north);
\draw[cycle] (explain.south) -- (facilitate.east);
\draw[cycle] (facilitate.west) -- (assess.south);
\draw[cycle] (assess.north) -- (prepare.west);
\end{tikzpicture}

\vspace{0.72cm}
\begin{tcolorbox}[
  enhanced,width=0.82\linewidth,colback=white,colframe=UTPBlue!18,
  arc=1.5mm,boxrule=0.45pt,left=6mm,right=6mm,top=3.2mm,bottom=3.2mm,
  borderline west={1.6pt}{0pt}{UTPGold}]
\centering\normalsize\color{UTPInk}
{\large\bfseries\color{UTPBlue}\utpoutlineheading\par}
\vspace{0.3em}
\utpcourseoutline
\end{tcolorbox}

\vfill
\begin{minipage}{0.82\textwidth}
\begin{tabular}{@{}>{\RaggedRight\arraybackslash}m{0.42\linewidth}@{\hspace{0.10\linewidth}}>{\RaggedLeft\arraybackslash}m{0.48\linewidth}@{}}
{\normalsize\bfseries\color{UTPBlue}PREPARED BY\par\vspace{0.25em}
 \Large\bfseries\color{UTPNavy}\utpauthor\par\vspace{0.12em}
 \normalfont\small\color{UTPMuted}\utpauthorrole}
&
\includegraphics[height=2.10cm,keepaspectratio]{logo-DC.png}
\end{tabular}
\end{minipage}
\vspace{0.5cm}
\end{titlepage}%
  \cleardoublepage
  \pagenumbering{roman}%
  \setcounter{tocdepth}{0}%
  \utptableofcontents
  \newpage
  \hypersetup{pageanchor=true}%
  \ifx\utpcoursemap\empty\else\input{\utpcoursemap}\fi
  \cleardoublepage
  \ifdefined\ReadonlyShipoutCounter
    \immediate\write\csname @auxout\endcsname
      {\string\gdef\string\utpfrontpages{\the\ReadonlyShipoutCounter}}%
  \fi
  \pagenumbering{arabic}%
}

% --------------------------------------------------------------------
%  15.  LISTING STYLES
% --------------------------------------------------------------------
%  Two base looks, and one named style per language built on them, so
%  the colours and frame are the same in every book:
%
%      utpcode          the base for source code
%      utpterminal      the base for commands typed and output printed
%      utppython        Python
%      utpr             R
%      utpc             C
%      utpcpp           C++
%      utpbash          shell scripts and commands
%      utpcnumbered     C with line numbers, for a complete source file
%      utpbashnumbered  shell with line numbers, likewise
%
%      \begin{lstlisting}[style=utppython] ... \end{lstlisting}
%      \lstinputlisting[style=utpcnumbered]{all_experiments/Lab04/fcfs.c}
%
%  A book that sets most of its listings in one style makes it the
%  default once, after the input:  \lstset{style=utpcpp}

\lstdefinestyle{utpcode}{
  basicstyle=\ttfamily\scriptsize,
  keywordstyle=\bfseries\color{UTPBlue},
  commentstyle=\itshape\color{UTPMuted},
  stringstyle=\color{UTPGreen!70!black},
  backgroundcolor=\color{UTPPaleGray!72},
  frame=single,
  rulecolor=\color{UTPBlue!28},
  framesep=5pt,
  breaklines=true,
  breakatwhitespace=false,
  showstringspaces=false,
  keepspaces=true,
  columns=fullflexible,
  upquote=true,
  tabsize=4,
  aboveskip=0.7em,
  belowskip=0.8em
}
\lstdefinestyle{utpterminal}{
  basicstyle=\ttfamily\scriptsize,
  backgroundcolor=\color{UTPPaleGray!72},
  frame=single,
  rulecolor=\color{UTPMuted!35},
  framesep=5pt,
  breaklines=true,
  showstringspaces=false,
  keepspaces=true,
  columns=fullflexible,
  upquote=true,
  aboveskip=0.6em,
  belowskip=0.7em
}

\lstdefinestyle{utppython}{style=utpcode,language=Python}
\lstdefinestyle{utpr}{style=utpcode,language=R,tabsize=2}
\lstdefinestyle{utpc}{style=utpcode,language=C}
\lstdefinestyle{utpcpp}{style=utpcode,language=C++}
\lstdefinestyle{utpbash}{style=utpcode,language=bash,
  rulecolor=\color{UTPMuted!35},aboveskip=0.6em,belowskip=0.7em}
\lstdefinestyle{utpcnumbered}{style=utpc,numbers=left,numberstyle=\tiny\color{UTPMuted},
  numbersep=7pt,xleftmargin=1.7em,framexleftmargin=1.4em}
\lstdefinestyle{utpbashnumbered}{style=utpbash,numbers=left,numberstyle=\tiny\color{UTPMuted},
  numbersep=7pt,xleftmargin=1.7em,framexleftmargin=1.4em}

% Pseudocode. A book that sets algorithms asks for them once, after the
% input, and gets algorithm2e loaded and coloured to match:
%
%     \utpusealgorithms
%
% \algostyle then sets one algorithm environment a size smaller than the
% text around it.
\newcommand{\utpalgorithmstyle}{%
  \SetAlCapNameFnt{\bfseries\color{UTPNavy}}%
  \SetAlCapFnt{\bfseries\color{UTPNavy}}%
  \SetAlgoCaptionSeparator{.}%
  \SetKwInput{KwIn}{Input}%
  \SetKwInput{KwOut}{Output}%
  \SetNlSty{color}{UTPMuted}{}%
  \SetCommentSty{\itshape\color{UTPMuted}}%
  \SetKwSty{bfseries\color{UTPBlue}}%
  \SetArgSty{}%
  \providecommand{\algostyle}{\small\SetAlgoNlRelativeSize{-1}\DontPrintSemicolon\SetAlgoVlined}%
}
\newcommand{\utpusealgorithms}{%
  \usepackage[linesnumbered,ruled,vlined]{algorithm2e}%
  \utpalgorithmstyle
}

%
%      \utpcontentspage{\small}{\normalsize}{0.65em}{0.07em}{0.16em}
%
%      \begin{document}
%      \utpfrontmatter                     % cover, contents, map
%      \part{Orientation}
%      \input{chapters/01-where-to-begin}
%      \end{document}
%
%  Configuration goes BEFORE the \input; overrides go AFTER it.
%  Paths are relative, so compile from inside the course folder.
%
%  ------------------------------------------------------------------
%  NAMING
%  ------------------------------------------------------------------
%  One name per idea, used by every book. Anything a book can use is
%  either an unprefixed learning box (keyidea, pitfall, labproject,
%  ...) or carries the utp prefix (utpfigure, utptable, utpfig,
%  utpsteps, utppython, UTPNavy, utpcell, ...). Course prefixes such
%  as ml-, ds- or erp- are not used: a chapter moved from one book to
%  another compiles unchanged. .shared/README.md lists every name.
%
%  Adding a new course? Copy .shared/TEMPLATE-main.tex and follow
%  .shared/ADDING-A-COURSE.md. Every step is a numbered checklist.
%
%  ------------------------------------------------------------------
%  CONTENTS
%  ------------------------------------------------------------------
%    1  Course configuration     the knobs a book may set
%    2  Packages                 the shared package set and fonts
%    3  Colour palette           the one UTP palette
%    4  Page layout and tables   margins, lists, table styling
%    5  Table of contents        parts, chapters, labs, figures
%    6  Headings                 chapter badge, part pages, sections
%    7  Page furniture           running heads and footers
%    8  Learning boxes           key idea, pitfall, try it, and friends
%    9  Small commands           \term, \cmd, \chapterguide, utpsteps
%   10  Diagram node styles      the TikZ vocabulary every book shares
%   11  Figure frames            utpfigure, utpfig, source images
%   12  Lab sheets               the lab banner and the lab toolkit
%   13  Stage divider            the part-opening page in longer books
%   14  Front matter             \utpfrontmatter: cover, contents, map
%   15  Listing styles           code, terminal and pseudocode
% =====================================================================

% --------------------------------------------------------------------
%  1.  COURSE CONFIGURATION
% --------------------------------------------------------------------
%  The knobs below, plus \utpstyleparts in section 5 and the box titles
%  in section 8, are the whole configuration surface. A book sets what
%  it needs BEFORE the \input and inherits the rest.
%
%  \providecommand means "define unless the book already did", which is
%  what makes the \def-before-\input pattern work.
%
%  Running text
%    \utprunninghead   the left running head on every page.
%    \utpchapterlabel  the word in the chapter badge. CHAPTER by
%                      default; a book whose chapters are labs sets LAB.
%    \utplabtotal      the N in the "Lab 3 of N" banners.

\providecommand{\utprunninghead}{Course Notes}
\providecommand{\utpchapterlabel}{CHAPTER}
\providecommand{\utplabtotal}{10}

%  Cover, printed by section 14
%    \utpcoursetitle       the title on the title page.
%    \utpdocumentsubtitle  the line under the title; Teaching Assistant Guide by default.
%    \utpcourseoutline     the sentence or two below it; \\ breaks a line.
%    \utpoutlineheading    the small heading above that sentence.
%    \utpauthor            the name in the PREPARED BY block.
%    \utpauthorrole        the line under the name.

\providecommand{\utpcoursetitle}{\utprunninghead}
\providecommand{\utpdocumentsubtitle}{Teaching Assistant Guide}
\providecommand{\utpcourseoutline}{}
\providecommand{\utpoutlineheading}{COURSE OUTLINE}
\providecommand{\utpauthor}{Reeyan Afzal Khan}
\providecommand{\utpauthorrole}{Graduate Research Assistant}

%  Files and folders
%    \utpextragraphicspath  image folders to search before the standard
%                           ones, as brace groups: {{screenshots/}}.
%                           A course's own figures/ folder and the
%                           shared logos are always searched.
%    \utpcoursemap          the overview chapter \utpfrontmatter inputs
%                           after the contents. Empty means none.

\providecommand{\utpextragraphicspath}{}
\providecommand{\utpcoursemap}{chapters/course-map}

%  Contents entries for chapters (section 5 applies them). A book whose
%  chapters are labs says so once here rather than restating the three
%  tocloft settings after the input.
%    \utpchaptertocword      Chapter~ by default, Lab~ in one book.
%    \utpchaptertocindent    left indent of the entry.
%    \utpchaptertocnumwidth  room reserved for the word and number.

\providecommand{\utpchaptertocword}{Chapter~}
\providecommand{\utpchaptertocindent}{1.5em}
\providecommand{\utpchaptertocnumwidth}{6em}

%  Monospace face (section 2)
%    \utpmonofont   sourcecodepro by default. inconsolata is the
%                   alternative, for a book whose pages are mostly
%                   device transcripts.

\providecommand{\utpmonofont}{sourcecodepro}

% --------------------------------------------------------------------
%  2.  PACKAGES
% --------------------------------------------------------------------
%  The union of what the books need. A course that needs something
%  extra loads it in its own file, under the heading "packages only
%  this book needs". Pseudocode is the one optional package the shared
%  file knows about: \utpusealgorithms in section 15 loads and styles it.
%
%  The tikz and pgfplots library lists are the union across all books,
%  so a diagram or chart copied from one course compiles in another.

\usepackage[T1]{fontenc}
\IfFileExists{glyphtounicode.tex}{\input{glyphtounicode}\pdfgentounicode=1}{}
\usepackage[american]{babel}
\IfFileExists{sourcesanspro.sty}
  {\usepackage[default]{sourcesanspro}}
  {\usepackage[scaled]{helvet}\renewcommand{\familydefault}{\sfdefault}}
\DeclareFontShape{T1}{cmss}{b}{n}{<->ssub*cmss/bx/n}{}
\DeclareFontShape{TS1}{cmss}{b}{n}{<->ssub*cmss/bx/n}{}
% Monospace. Without this the tt default stays Computer Modern Typewriter,
% a serif face that clashes with the sans body text in every listing and
% every \cmd. Source Code Pro is the companion to Source Sans Pro; it is
% scaled down because mono digits otherwise sit large next to the sans.
%
% Inconsolata, chosen with \def\utpmonofont{inconsolata}, distinguishes
% 0/O and 1/l/I at a glance, which matters when a reader copies an
% interface name or an address out of a transcript. Only one mono package
% is ever loaded: a second would win the family but inherit the first's
% series names, and every \ttfamily would then resolve by substitution.
\def\utpmonoinconsolata{inconsolata}
\ifx\utpmonofont\utpmonoinconsolata
  \IfFileExists{inconsolata.sty}{\usepackage[varqu,varl,scaled=1.02]{inconsolata}}{}
\else
  \IfFileExists{sourcecodepro.sty}
    {\usepackage[scaled=0.88]{sourcecodepro}}
    {\usepackage[scaled=0.86]{beramono}}
  % Source Sans names its normal weight "regular", but the Computer Modern
  % fallbacks call it "m". Map one onto the other so a \ttfamily reached
  % while the regular series is active resolves instead of warning.
  \DeclareFontShape{T1}{cmtt}{regular}{n}{<->ssub*cmtt/m/n}{}
  \DeclareFontShape{TS1}{cmtt}{regular}{n}{<->ssub*cmtt/m/n}{}
  \DeclareFontShape{T1}{cmtt}{regular}{it}{<->ssub*cmtt/m/it}{}
\fi
\usepackage[margin=21mm,headheight=22pt]{geometry}
\usepackage{amsmath,amssymb,mathtools,bm}
\usepackage{microtype}
\usepackage{enumitem}
\usepackage{booktabs}
\usepackage{tabularx}
\usepackage{array}
\usepackage{ragged2e}
\usepackage{etoolbox}
\usepackage{titlesec}
\usepackage{tocloft}
\usepackage{fancyhdr}
\usepackage[table]{xcolor}
\usepackage{tcolorbox}
\tcbuselibrary{breakable,skins}
\usepackage[hidelinks]{hyperref}
\usepackage{bookmark}
% Document properties, so the PDF identifies itself in a reader's
% title bar, in a library catalogue, and in search results. Set from
% the same configuration macros the cover uses.
\hypersetup{
  pdftitle={\utpcoursetitle},
  pdfauthor={\utpauthor},
  pdfsubject={Universiti Teknologi PETRONAS course notes and laboratory guide},
  pdfkeywords={UTP, course notes, laboratory manual, \utpcoursetitle},
  pdfdisplaydoctitle=true,
  bookmarksnumbered=true,
  bookmarksopenlevel=1
}
\usepackage{listings}
\usepackage{graphicx}
\usepackage{caption}
\usepackage{tikz}
\usetikzlibrary{angles,arrows.meta,backgrounds,calc,chains,decorations.pathmorphing,decorations.pathreplacing,fit,matrix,patterns,positioning,quotes,shadows.blur,shapes.geometric,shapes.misc,shapes.symbols}
\usepackage{pgfplots}
% statistics gives boxplot prepared; groupplots gives multi-panel charts.
\usepgfplotslibrary{statistics}
\usepgfplotslibrary{groupplots}

% Images: any folders the book names first, then its own figures/
% folder, then the shared logos used by the cover.
\expandafter\graphicspath\expandafter{\utpextragraphicspath{figures/}{../.shared/logos/}}

% --------------------------------------------------------------------
%  3.  COLOUR PALETTE
% --------------------------------------------------------------------
%  One palette, defined once under UTP* names, and used by every book,
%  box, diagram and chart. Five accents (blue, teal, green, orange, red)
%  with a pale tint of each for box backgrounds, plus navy for headings,
%  gold for the cover, ink and muted for text, and a sixth accent,
%  purple, for charts and layered diagrams that need one more series.

\definecolor{UTPNavy}{HTML}{18324A}
\definecolor{UTPBlue}{HTML}{246B9E}
\definecolor{UTPTeal}{HTML}{138A80}
\definecolor{UTPGreen}{HTML}{2D7D46}
\definecolor{UTPOrange}{HTML}{C86B24}
\definecolor{UTPRed}{HTML}{B34843}
\definecolor{UTPPurple}{HTML}{6B4FA0}
\definecolor{UTPGold}{HTML}{D4A13E}
\definecolor{UTPInk}{HTML}{263238}
\definecolor{UTPMuted}{HTML}{60717C}
\definecolor{UTPPaleBlue}{HTML}{EEF6FB}
\definecolor{UTPPaleTeal}{HTML}{EEF8F6}
\definecolor{UTPPaleOrange}{HTML}{FFF5EA}
\definecolor{UTPPaleRed}{HTML}{FFF0EF}
\definecolor{UTPPalePurple}{HTML}{F0ECF9}
\definecolor{UTPPaleGray}{HTML}{F4F6F7}
\color{UTPInk}

% --------------------------------------------------------------------
%  4.  PAGE LAYOUT AND TABLES
% --------------------------------------------------------------------
%  Paragraph spacing, list indents, and the table look: alternating
%  row tint, top-aligned paragraph columns, ragged-right cells.
%
%  Two extra column types for tabularx tables:
%      D{width}   a fixed-width, ragged-right, vertically centred column
%      Y          a ragged-right column that fills the remaining width
%
%  Every table in every book has the same shape, the one Algorithms &
%  Data Structures set: a utptable frame around a full-width tabularx,
%  booktabs rules, a tinted header row of \utpth cells (centred, bold,
%  navy), and a first column in bold navy when it names the rows.
%  Cells are vertically centred, so use m{} rather than p{} columns.
%
%      \begin{utptable}
%      \begin{tabularx}{\linewidth}{@{}>{\bfseries\leavevmode\color{UTPNavy}}m{0.22\linewidth}X@{}}
%      \toprule
%      \rowcolor{UTPPaleBlue!92}
%      \utpth{Operation} & \utpth{What it does} \\
%      \midrule
%      Insertion & Add a new element. \\
%      \bottomrule
%      \end{tabularx}
%      \end{utptable}
%
%  A book whose tables carry dense text asks for \utpcompacttables after
%  the input: tabularx tables then set in \footnotesize with tighter
%  padding and vertically centred cells.

\setlength{\parindent}{0pt}
\setlength{\parskip}{0.55em}
\setlist[itemize]{leftmargin=1.45em,itemsep=0.22em,topsep=0.25em}
\setlist[enumerate]{leftmargin=1.75em,itemsep=0.28em,topsep=0.3em}
\setlist[description]{font=\bfseries\color{UTPNavy},itemsep=0.25em,topsep=0.3em}
\renewcommand{\arraystretch}{1.34}
\setlength{\tabcolsep}{6.5pt}
\arrayrulecolor{UTPBlue!45}
% Keep paragraph-style columns ragged and vertically centred. Tables use
% fixed m columns beside flexible X columns; giving X the same centring
% keeps a short cell level with a tall neighbour in a multi-line row.
\renewcommand{\tabularxcolumn}[1]{m{#1}}
\makeatletter
\g@addto@macro\@arrayparboxrestore{\RaggedRight\arraybackslash}
\makeatother
% A header cell: centred, bold and navy, on the tinted header row.
\newcommand{\utpth}[1]{\multicolumn{1}{c}{\bfseries\color{UTPNavy}#1}}
\AtBeginEnvironment{tabular}{\small\setlength{\extrarowheight}{0.6pt}\rowcolors{2}{UTPPaleBlue!72}{white}}
\AtBeginEnvironment{tabularx}{\small\setlength{\extrarowheight}{0.6pt}\rowcolors{2}{UTPPaleBlue!72}{white}}
\AtBeginEnvironment{figure}{\centering}
\newcolumntype{D}[1]{>{\raggedright\arraybackslash}m{#1}}
\newcolumntype{Y}{>{\raggedright\arraybackslash}X}
\newcommand{\utpcompacttables}{%
  \renewcommand{\tabularxcolumn}[1]{m{##1}}%
  \AtBeginEnvironment{tabularx}{%
    \footnotesize
    \setlength{\tabcolsep}{4pt}%
    \renewcommand{\arraystretch}{1.25}%
    \setlength{\extrarowheight}{0.6pt}%
  }%
}
\emergencystretch=1.5em
\hbadness=10000
\hfuzz=8pt

% Page-breaking discipline. Without these a paragraph can leave a single
% line stranded at the foot of one page or the head of the next, which is
% the most visible sign of unattended typesetting.
\widowpenalty=10000          % no first line of a paragraph left alone at a page foot
\clubpenalty=10000           % no last line of a paragraph left alone at a page head
\displaywidowpenalty=10000   % the same around displayed equations
\brokenpenalty=10000         % do not break a page after a hyphenated line
\predisplaypenalty=10000     % keep the line before a display with it
% Pages are built from boxes of many heights, so let the last line float
% up rather than stretching the gaps between boxes to force a flush foot.
\raggedbottom
% Prefer a slightly loose line to an overfull one, but do not hyphenate
% freely -- short sans lines look ragged when over-hyphenated.
\tolerance=1200
\hyphenpenalty=200

% Character protrusion and font expansion. Both are what stop a justified
% sans column from showing visible rivers down the page.
\microtypesetup{protrusion=true,expansion=true,tracking=true}

\pgfplotsset{
  compat=1.18,
  every axis/.append style={
    axis line style={draw=UTPMuted!85,line width=0.55pt},
    tick style={draw=UTPMuted!85},
    tick label style={font=\scriptsize\color{UTPMuted}},
    label style={font=\small\color{UTPNavy}},
    title style={font=\small\bfseries\color{UTPNavy}},
    legend style={font=\scriptsize,draw=UTPBlue!22,fill=white,rounded corners=1pt}
  },
  every axis plot/.append style={line width=1.15pt}
}
\tikzset{every picture/.append style={line cap=round,line join=round}}

% --------------------------------------------------------------------
%  5.  TABLE OF CONTENTS
% --------------------------------------------------------------------
%  Formatting for part, chapter and lab entries, and for a list of
%  figures if a book prints one. The lab level is a custom tocloft list
%  so lab banners appear in the contents.
%
%  Part styling is skipped when \utpstyleparts is 0, which is what a
%  book without \part divisions should set.
%
%  The contents page is meant to fill exactly one page, and how loosely
%  it is set depends on how many entries a book has. One line after the
%  input tunes it:
%
%      \utpcontentspage{<base size>}{<part size>}%
%                      {<space before a part>}{<before a chapter>}{<before a lab>}
%      \utpcontentspage{\normalsize}{\large}{1.15em}{0.38em}{0.30em}   % the default

\addto\captionsamerican{\renewcommand{\contentsname}{Learning Path}}
\newlistentry{lab}{toc}{-1}
\providecommand{\utpstyleparts}{1}
\ifnum\utpstyleparts=1
\renewcommand{\cftpartpresnum}{}
\renewcommand{\cftpartaftersnum}{}
\setlength{\cftpartnumwidth}{5.2em}
\setlength{\cftbeforepartskip}{0.38em}
\renewcommand{\cftpartfont}{\normalsize\bfseries\color{UTPNavy}}
\renewcommand{\cftpartpagefont}{\normalsize\bfseries\color{UTPNavy}}
\renewcommand{\cftpartleader}{\hfill}
\fi
\renewcommand{\cftchappresnum}{\utpchaptertocword}
\renewcommand{\cftchapaftersnum}{\quad}
\setlength{\cftchapindent}{\utpchaptertocindent}
\setlength{\cftchapnumwidth}{\utpchaptertocnumwidth}
\setlength{\cftbeforechapskip}{0pt}
\renewcommand{\cftchapfont}{\normalfont}
\renewcommand{\cftchappagefont}{\normalfont}
\renewcommand{\cftchapleader}{\cftdotfill{\cftdotsep}}
\setlength{\cftlabindent}{3em}
\setlength{\cftlabnumwidth}{0em}
\setlength{\cftbeforelabskip}{0.10em}
\renewcommand{\cftlabfont}{\bfseries\color{UTPOrange!90!black}}
\renewcommand{\cftlabpagefont}{\bfseries\color{UTPOrange!90!black}}
\renewcommand{\cftlableader}{\color{UTPOrange!65!UTPMuted}\cftdotfill{\cftlabdotsep}}

% A list of figures, for a book that prints one (\listoffigures).
\renewcommand{\cftfigfont}{\normalfont\small}
\renewcommand{\cftfigpagefont}{\normalfont\small}
\renewcommand{\cftfigpresnum}{Figure~}
\renewcommand{\cftfigaftersnum}{\quad}
\setlength{\cftfigindent}{1.5em}
\setlength{\cftfignumwidth}{5.6em}
\setlength{\cftbeforefigskip}{0pt}
\renewcommand{\cftfigleader}{\cftdotfill{\cftdotsep}}

% The contents page itself, printed by \utpfrontmatter (section 14).
\newcommand{\utpcontentspage}[5]{%
  \def\utptocsize{#1}\def\utptocpartsize{#2}%
  \def\utptocpartskip{#3}\def\utptocchapterskip{#4}\def\utptoclabskip{#5}}
\utpcontentspage{\normalsize}{\large}{1.15em}{0.38em}{0.30em}
\newcommand{\utptableofcontents}{%
  \begingroup
  \utptocsize
  \setlength{\parskip}{0pt}%
  \setlength{\cftbeforepartskip}{\utptocpartskip}%
  \setlength{\cftbeforechapskip}{\utptocchapterskip}%
  \setlength{\cftbeforelabskip}{\utptoclabskip}%
  \renewcommand{\cftpartfont}{\utptocpartsize\bfseries\color{UTPNavy}}%
  \renewcommand{\cftpartpagefont}{\utptocpartsize\bfseries\color{UTPNavy}}%
  \tableofcontents
  \endgroup
}

% --------------------------------------------------------------------
%  6.  HEADINGS
% --------------------------------------------------------------------
%  The chapter badge (a coloured block with a word and the number),
%  part pages, and section spacing.
%
%  The word in the badge is \utpchapterlabel; Data Communication sets
%  it to LAB because each of its chapters is a lab.

\titleformat{\chapter}[display]
  {\bfseries\color{UTPNavy}}
  {\filleft\colorbox{UTPBlue}{\parbox{2.55cm}{\centering\color{white}\Large
    \utpchapterlabel\\[-0.15em]\Huge\thechapter}}}
  {0.65em}{\Huge\raggedright}
  [\vspace{0.25em}{\color{UTPBlue}\titlerule[1.2pt]}]
\renewcommand{\thepart}{Stage~\arabic{part}}
\titleformat{\part}[display]
  {\centering\color{UTPNavy}\bfseries}
  {\Large\color{UTPTeal}\thepart}{0.45em}{\Huge}
  [\vspace{0.55em}\color{UTPTeal}\rule{0.45\linewidth}{1.5pt}]
\titleformat{\section}{\Large\bfseries\color{UTPNavy}}{\thesection}{0.55em}{}
\titleformat{\subsection}{\large\bfseries\color{UTPBlue}}{\thesubsection}{0.5em}{}
\titlespacing*{\chapter}{0pt}{-1.8em}{1.0em}
\titlespacing*{\part}{0pt}{0pt}{1em}
\titlespacing*{\section}{0pt}{1.2em}{0.35em}
\titlespacing*{\subsection}{0pt}{0.95em}{0.2em}

% --------------------------------------------------------------------
%  7.  PAGE FURNITURE
% --------------------------------------------------------------------
%  Running heads and footers. The left header is \utprunninghead;
%  the right one follows the current chapter.

\pagestyle{fancy}
\fancyhf{}
\lhead{\small\color{UTPMuted}\utprunninghead}
\rhead{\small\color{UTPMuted}\nouppercase{\leftmark}}
\cfoot{\small\color{UTPMuted}\thepage}
\renewcommand{\headrulewidth}{0.5pt}
\renewcommand{\headrule}{\hbox to\headwidth{\color{UTPPaleBlue}\leaders\hrule height \headrulewidth\hfill}}
\renewcommand{\chaptermark}[1]{\markboth{#1}{}}
\fancypagestyle{plain}{%
  \fancyhf{}%
  \cfoot{\small\color{UTPMuted}\thepage}%
  \renewcommand{\headrulewidth}{0pt}%
}

% --------------------------------------------------------------------
%  8.  LEARNING BOXES
% --------------------------------------------------------------------
%  The coloured callouts the chapters use. Every title is a macro, so
%  a book can rename a box without restating its geometry:
%
%      \def\utptitlepitfall{Technical note}   % before the input
%
%  To change how a box LOOKS in one book, use one of the named
%  variations at the end of this section, or \renewtcolorbox after the
%  input. To add a box only one book needs, declare it in that book
%  with \utpnewcallout or \utpnewquietbox.

% Box titles -- override any of these before the input line.
\providecommand{\utptitlekeyidea}{Key idea}
\providecommand{\utptitleworked}{Worked example}
\providecommand{\utptitlepitfall}{Common mistake}
\providecommand{\utptitletryit}{Try it yourself}
\providecommand{\utptitlequickcheck}{Check your understanding}
\providecommand{\utptitlerecap}{Chapter recap}
\providecommand{\utptitledeliverable}{Lab deliverable}
\providecommand{\utptitlelabmode}{How much is scaffolded?}
\providecommand{\utptitleintuition}{Plain idea}
\providecommand{\utptitlefirstread}{Core path}
\providecommand{\utptitleglance}{Chapter at a glance}
\providecommand{\utptitledebug}{Debugging task}
\providecommand{\utptitleoptional}{Beyond the core}
\providecommand{\utptitleverification}{Verification checklist}
\providecommand{\utptitletroubleshoot}{Troubleshooting}
\providecommand{\utptitlesourcebasis}{Source basis}

% Explaining an idea
\newtcolorbox{keyidea}{enhanced,breakable,colback=UTPPaleTeal,colframe=UTPTeal, boxrule=0.8pt,arc=2mm,left=2.5mm,right=2.5mm,top=1.8mm,bottom=1.8mm, fonttitle=\bfseries\color{UTPNavy},title=\utptitlekeyidea}
\newtcolorbox{workedexample}{enhanced,breakable,colback=UTPPaleBlue,colframe=UTPBlue, boxrule=0.8pt,arc=2mm,left=2.5mm,right=2.5mm,top=1.8mm,bottom=1.8mm, fonttitle=\bfseries\color{UTPNavy},title=\utptitleworked}
\newtcolorbox{pitfall}{enhanced,breakable,colback=UTPPaleRed,colframe=UTPRed, boxrule=0.8pt,arc=2mm,left=2.5mm,right=2.5mm,top=1.8mm,bottom=1.8mm, fonttitle=\bfseries\color{UTPNavy},title=\utptitlepitfall}
% Asking the reader to do something
\newtcolorbox{tryit}{enhanced,breakable,colback=white,colframe=UTPTeal, boxrule=0.9pt,arc=2mm,left=2.5mm,right=2.5mm,top=1.8mm,bottom=1.8mm, fonttitle=\bfseries\color{white},colbacktitle=UTPTeal,title=\utptitletryit}
\newtcolorbox{quickcheck}{enhanced,breakable,colback=white,colframe=UTPOrange, boxrule=0.9pt,arc=2mm,left=2.5mm,right=2.5mm,top=1.8mm,bottom=1.8mm, fonttitle=\bfseries\color{white},colbacktitle=UTPOrange,title=\utptitlequickcheck}
\newtcolorbox{chapterrecap}{enhanced,breakable,colback=UTPPaleTeal,colframe=UTPGreen, boxrule=0.9pt,arc=2mm,left=2.5mm,right=2.5mm,top=1.8mm,bottom=1.8mm, fonttitle=\bfseries\color{white},colbacktitle=UTPGreen,title=\utptitlerecap}
% Lab structure
\newtcolorbox{labdeliverable}{enhanced,breakable,colback=white,colframe=UTPGreen, boxrule=0.9pt,arc=2mm,left=2.6mm,right=2.6mm,top=1.9mm,bottom=1.9mm, fonttitle=\bfseries\color{white},colbacktitle=UTPGreen,title=\utptitledeliverable}
\newtcolorbox{labmode}{enhanced,breakable,colback=UTPPaleBlue!52,colframe=UTPBlue, boxrule=0.75pt,arc=2mm,left=2.6mm,right=2.6mm,top=1.7mm,bottom=1.7mm, fonttitle=\bfseries\color{white},colbacktitle=UTPBlue,title=\utptitlelabmode}
% Softer, lighter-weight asides
\newtcolorbox{intuition}{enhanced,breakable,colback=white,colframe=UTPOrange!42, boxrule=0.5pt,arc=2mm,left=2.5mm,right=2.5mm,top=1.5mm,bottom=1.5mm, fonttitle=\bfseries\color{UTPOrange!85!black},title=\utptitleintuition}
\newtcolorbox{firstread}{enhanced,breakable,colback=white,colframe=UTPBlue!42, boxrule=0.5pt,arc=2mm,left=2.5mm,right=2.5mm,top=1.5mm,bottom=1.5mm, fonttitle=\bfseries\color{UTPBlue}, title=\utptitlefirstread}
\newtcolorbox{debugtask}{enhanced,breakable,colback=UTPPaleRed!58,colframe=UTPRed, boxrule=0.85pt,arc=2mm,left=2.6mm,right=2.6mm,top=1.9mm,bottom=1.9mm, fonttitle=\bfseries\color{white},colbacktitle=UTPRed,title=\utptitledebug}
\newtcolorbox{analogy}{enhanced,breakable,colback=white,colframe=UTPOrange!38, boxrule=0.5pt,arc=2mm,left=2.5mm,right=2.5mm,top=1.5mm,bottom=1.5mm, fonttitle=\bfseries\color{UTPOrange!85!black},title=Analogy}
\newtcolorbox{coursechoice}{enhanced,breakable,colback=white,colframe=UTPNavy, boxrule=0.8pt,arc=2mm,left=2.8mm,right=2.8mm,top=2mm,bottom=2mm, fonttitle=\bfseries\color{white},colbacktitle=UTPNavy,title=Course choice}
\newtcolorbox{decisionmemo}{enhanced,breakable,colback=UTPPaleOrange!55,colframe=UTPOrange, boxrule=0.85pt,arc=2mm,left=2.6mm,right=2.6mm,top=1.9mm,bottom=1.9mm, fonttitle=\bfseries\color{white},colbacktitle=UTPOrange,title=Decision memo}
\newtcolorbox{libnote}{enhanced,breakable,colback=white,colframe=UTPMuted!48, boxrule=0.5pt,arc=2mm,left=2.5mm,right=2.5mm,top=1.5mm,bottom=1.5mm, fonttitle=\bfseries\color{UTPNavy},title=Reference note}
\newtcolorbox{miniproject}{enhanced,breakable,colback=UTPPaleTeal!52,colframe=UTPGreen, boxrule=0.9pt,arc=2mm,left=2.6mm,right=2.6mm,top=1.9mm,bottom=1.9mm, fonttitle=\bfseries\color{white},colbacktitle=UTPGreen,title=Mini-project}
% Where the material came from: the lecturer's slides, manual or tutorial
% that a chapter is built on.
\newtcolorbox{sourcebasis}{enhanced,breakable,colback=white,colframe=UTPMuted!55, boxrule=0.5pt,arc=2mm,left=2.5mm,right=2.5mm,top=1.5mm,bottom=1.5mm, fonttitle=\bfseries\color{UTPMuted},title=\utptitlesourcebasis}
% Boxes that take a title argument
\newtcolorbox{modelcard}[1]{enhanced,breakable,colback=white,colframe=UTPBlue, boxrule=0.9pt,arc=2mm,left=2.5mm,right=2.5mm,top=1.8mm,bottom=1.8mm, fonttitle=\bfseries\color{white},colbacktitle=UTPBlue,title=#1}
\newtcolorbox{optionaltopic}{enhanced,breakable,colback=white,colframe=UTPMuted!58, boxrule=0.5pt,arc=2mm,left=2.5mm,right=2.5mm,top=1.5mm,bottom=1.5mm, fonttitle=\bfseries\color{UTPMuted},title=\utptitleoptional}
\newtcolorbox{secondpass}{enhanced,breakable,colback=UTPPaleGray!82,colframe=UTPMuted, boxrule=0.75pt,arc=2mm,left=2.6mm,right=2.6mm,top=1.9mm,bottom=1.9mm, fonttitle=\bfseries\color{white},colbacktitle=UTPMuted,title=Second pass}
\newtcolorbox{labproject}[1]{enhanced,breakable,colback=white,colframe=UTPBlue, boxrule=0.9pt,arc=2mm,left=2.6mm,right=2.6mm,top=1.9mm,bottom=1.9mm, fonttitle=\bfseries\color{white},colbacktitle=UTPBlue,title={#1}}
% Evidence that a lab worked, and what to do when it did not
\newtcolorbox{verification}{enhanced,breakable,colback=UTPPaleTeal!65,colframe=UTPGreen, boxrule=0.8pt,arc=2mm,left=2.6mm,right=2.6mm,top=1.8mm,bottom=1.8mm, fonttitle=\bfseries\color{white},colbacktitle=UTPGreen,title=\utptitleverification}
\newtcolorbox{troubleshoot}{enhanced,breakable,colback=UTPPaleOrange!70,colframe=UTPOrange, boxrule=0.8pt,arc=2mm,left=2.6mm,right=2.6mm,top=1.8mm,bottom=1.8mm, fonttitle=\bfseries\color{white},colbacktitle=UTPOrange,title=\utptitletroubleshoot}

% Named variations a book asks for after the input. Each restyles boxes
% that already exist, so the chapters do not change.
%
%   \utpheavycallouts          a solid colour bar behind the titles of
%                              keyidea, workedexample, intuition and
%                              firstread, with a firmer frame.
%   \utpbreakableframes        let utpfigure and utptable split over a
%                              page break, in a lighter blue (section 11).
%   \utplabprojectcolour{...}  the labproject box in another colour, with
%                              a little more room inside.
\newcommand{\utplabprojectcolour}[1]{%
  \renewtcolorbox{labproject}[1]{enhanced,breakable,colback=white,colframe=#1, boxrule=0.9pt,arc=2mm,left=2.7mm,right=2.7mm,top=2mm,bottom=2mm, fonttitle=\bfseries\color{white},colbacktitle=#1,title={##1}}%
}
\newcommand{\utpheavycallouts}{%
  \renewtcolorbox{keyidea}{enhanced,breakable,colback=UTPPaleTeal,colframe=UTPTeal, boxrule=0.85pt,arc=2mm,left=2.5mm,right=2.5mm,top=1.8mm,bottom=1.8mm, fonttitle=\bfseries\color{white},colbacktitle=UTPTeal,title=\utptitlekeyidea}%
  \renewtcolorbox{workedexample}{enhanced,breakable,colback=UTPPaleBlue,colframe=UTPBlue, boxrule=0.8pt,arc=2mm,left=2.5mm,right=2.5mm,top=1.8mm,bottom=1.8mm, fonttitle=\bfseries\color{white},colbacktitle=UTPBlue,title=\utptitleworked}%
  \renewtcolorbox{intuition}{enhanced,breakable,colback=UTPPaleOrange,colframe=UTPOrange, boxrule=0.8pt,arc=2mm,left=2.5mm,right=2.5mm,top=1.8mm,bottom=1.8mm, fonttitle=\bfseries\color{white},colbacktitle=UTPOrange,title=\utptitleintuition}%
  \renewtcolorbox{firstread}{enhanced,breakable,colback=white,colframe=UTPBlue, boxrule=0.9pt,arc=2mm,left=2.5mm,right=2.5mm,top=1.8mm,bottom=1.8mm, fonttitle=\bfseries\color{white},colbacktitle=UTPBlue,title=\utptitlefirstread}%
}

% Declaring a box a single book needs. Both constructors take the house
% geometry and let a trailing key list override any of it, so a new box
% is one line and still matches the others:
%
%   \utpnewcallout{odoobox}{white}{UTPBlue}{Odoo 19 Community}
%   \utpnewcallout[boxrule=0.9pt]{name}{back}{frame}{Title}
%   \utpnewquietbox{referencebox}{UTPMuted!58}{UTPMuted}{Reference note}
%
% A callout is a tinted box under a solid title bar, for something the
% reader has to act on. A quiet box is a light frame with coloured
% title text, for an aside they can read past. If a second book ever
% wants the same box, move its line into this section.
\newcommand{\utpnewcallout}[5][]{%
  \newtcolorbox{#2}{enhanced,breakable,colback=#3,colframe=#4, boxrule=0.8pt,arc=2mm,left=2.6mm,right=2.6mm,top=1.8mm,bottom=1.8mm, fonttitle=\bfseries\color{white},colbacktitle=#4,title={#5},#1}%
}
\newcommand{\utpnewquietbox}[5][]{%
  \newtcolorbox{#2}{enhanced,breakable,colback=white,colframe=#3, boxrule=0.5pt,arc=2mm,left=2.5mm,right=2.5mm,top=1.5mm,bottom=1.5mm, fonttitle=\bfseries\color{#4},title={#5},#1}%
}

% --------------------------------------------------------------------
%  9.  SMALL COMMANDS
% --------------------------------------------------------------------
%  \term for a defined term, \cmd for inline code, \chapref and
%  \secref for cross-references, and the utpsteps numbered-step list.
%
%      \term{overfitting}          a term being defined
%      \cmd{git status}            inline code, in any language
%      \chapref{ch:trees}          "Chapter 12"
%      \secref{sec:pruning}        "Section 12.3"
%      \chapterguide{prior}{outcomes}   the box that opens a chapter
%      \begin{utpsteps} ... \end{utpsteps}  a numbered Step 1/Step 2 list
%
%  The second group is the shorthand the chapters use in running text:
%
%      \figref{fig:tree}           "Figure 4.2"
%      \chapsref{ch:a}{ch:b}       "Chapters 3--5"
%      \lib{pandas}                a library name
%      \cxg{x}, \cxo{x}, \cxr{x}   maths coloured green/orange/red
%      \given                      the bar in P(a \given b)
%      \coretag, \latertag         read-now / come-back-to markers
%      \kagglesource{url}{name}    the dataset a lab is built on

\newcommand{\term}[1]{\textbf{\color{UTPNavy}#1}}
\newcommand{\cmd}[1]{\texttt{#1}}
\newcommand{\chapref}[1]{Chapter~\ref{#1}}
\newcommand{\secref}[1]{Section~\ref{#1}}
\newcommand{\chapterguide}[2]{%
  \begin{tcolorbox}[enhanced,breakable,colback=white,colframe=UTPNavy,boxrule=0.8pt,
    arc=2mm,left=2.5mm,right=2.5mm,top=2mm,bottom=2mm,
    title=\utptitleglance,fonttitle=\bfseries\color{white},colbacktitle=UTPNavy]
    \textbf{Before you start:} #1\par
    \textbf{By the end, you can:} #2
  \end{tcolorbox}}
\newenvironment{utpsteps}
  {\begin{enumerate}[label=\textcolor{UTPBlue}{\bfseries Step~\arabic*:},leftmargin=4.5em,itemsep=0.38em]}
  {\end{enumerate}}

\newcommand{\figref}[1]{Figure~\ref{#1}}
\newcommand{\chapsref}[2]{Chapters~\ref{#1}--\ref{#2}}
\newcommand{\lib}[1]{\textbf{\color{UTPBlue}\texttt{#1}}}
\newcommand{\cxg}[1]{{\color{UTPGreen}\bfseries $#1$}}
\newcommand{\cxo}[1]{{\color{UTPOrange}\bfseries $#1$}}
\newcommand{\cxr}[1]{{\color{UTPRed}\bfseries $#1$}}
\newcommand{\given}{\mathrel{|}}
\DeclareMathOperator*{\argmin}{arg\,min}
\DeclareMathOperator*{\argmax}{arg\,max}
\newcommand{\coretag}{\par\smallskip\noindent\colorbox{UTPPaleBlue}{\strut\small\bfseries\color{UTPBlue}CORE NOW}\par\smallskip}
\newcommand{\latertag}{\par\smallskip\noindent\colorbox{UTPPaleGray}{\strut\small\bfseries\color{UTPMuted}RETURN LATER}\par\smallskip}
\newcommand{\kagglesource}[2]{%
  \par\smallskip\noindent\textbf{Kaggle source:} \href{#1}{#2}\par
}

% --------------------------------------------------------------------
%  10.  DIAGRAM NODE STYLES
% --------------------------------------------------------------------
%  The vocabulary the hand-drawn TikZ diagrams are built from: a cell
%  in an array, the one cell under the cursor, a circular graph node,
%  the arrow between two of them. Every book has the whole set:
%
%      utpcell, utphot, utpdone, utpdim, utpwall        array cells
%      utpnode, utpnodehot, utpnodedone, utpnodedim     graph nodes
%      utparrow, utpptr, utpdash, utpedge               connectors
%      utplabel, utptag, utpbrace                       annotation
%      utpstep, utpstage, utppanel, utpbox              boxes
%
%  A cell is 1.0cm by 0.8cm. A book whose diagrams were drawn to a
%  different grid resizes the cells once, after the input:
%
%      \utpdiagramsize{1.05cm}{0.85cm}
%
%  Larger vocabularies for one kind of drawing live in .shared/diagrams/
%  and are loaded by name, after the input:
%
%      \utpusediagrams{network}         devices, links, packet headers,
%                                       sequence diagrams, state machines
%      \utpusediagrams{process}         business-process flows, documents,
%                                       swimlanes, status chips
%      \utpusediagrams{data-objects}    vectors, lists and data frames
%                                       drawn as cells; flowcharts
%      \utpusediagrams{network,process} several at once

\newcommand{\utpdiagramsize}[2]{%
  \tikzset{
    utpcell/.style={draw=UTPBlue,fill=UTPPaleBlue,minimum width=#1,minimum height=#2,
                    font=\small\bfseries\color{UTPNavy},inner sep=1pt},
    utphot/.style={draw=UTPOrange,fill=UTPPaleOrange,minimum width=#1,minimum height=#2,
                   font=\small\bfseries\color{UTPNavy},line width=1pt,inner sep=1pt},
    utpdone/.style={draw=UTPGreen,fill=UTPPaleTeal,minimum width=#1,minimum height=#2,
                    font=\small\bfseries\color{UTPNavy},inner sep=1pt},
    utpdim/.style={draw=UTPMuted!45,fill=UTPPaleGray,minimum width=#1,minimum height=#2,
                   font=\small\color{UTPMuted},inner sep=1pt},
    utpwall/.style={draw=UTPNavy,fill=UTPNavy!82,minimum width=#1,minimum height=#2,
                    inner sep=1pt},
    utpnode/.style={circle,draw=UTPBlue,fill=UTPPaleBlue,minimum size=8mm,
                    font=\small\bfseries\color{UTPNavy},inner sep=0.5pt},
    utpnodehot/.style={circle,draw=UTPOrange,fill=UTPPaleOrange,minimum size=8mm,line width=1pt,
                       font=\small\bfseries\color{UTPNavy},inner sep=0.5pt},
    utpnodedone/.style={circle,draw=UTPGreen,fill=UTPPaleTeal,minimum size=8mm,
                        font=\small\bfseries\color{UTPNavy},inner sep=0.5pt},
    utpnodedim/.style={circle,draw=UTPMuted!45,fill=UTPPaleGray,minimum size=8mm,
                       font=\small\color{UTPMuted},inner sep=0.5pt},
    utparrow/.style={-{Stealth[length=2.2mm]},thick,draw=UTPTeal},
    utpptr/.style={-{Stealth[length=2mm]},semithick,draw=UTPBlue},
    utpdash/.style={-{Stealth[length=2mm]},semithick,draw=UTPRed,dashed},
    utpedge/.style={thick,draw=UTPTeal},
    utplabel/.style={font=\scriptsize\color{UTPMuted},inner sep=1pt},
    utptag/.style={font=\scriptsize\bfseries\color{UTPOrange},inner sep=1pt},
    utpbrace/.style={decorate,decoration={brace,amplitude=4pt,raise=2pt},draw=UTPMuted,semithick},
    utpstep/.style={draw=UTPNavy,fill=white,rounded corners=1.5mm,minimum height=0.75cm,
                    font=\scriptsize\color{UTPNavy},align=center,inner sep=2pt},
    utpstage/.style={draw=UTPNavy,fill=white,rounded corners=1.5mm,minimum height=0.75cm,
                     font=\scriptsize\color{UTPNavy},align=center,inner sep=2pt},
    utppanel/.style={draw=UTPBlue!30,fill=white,rounded corners=1.5mm,inner sep=2.5mm},
    utpbox/.style={draw=UTPBlue,fill=UTPPaleBlue,rounded corners=2mm,align=center,
                   font=\small\color{UTPNavy},inner sep=2.2mm}
  }%
}
\utpdiagramsize{1.0cm}{0.8cm}

% \@for, not \foreach: \foreach runs each pass in a group, which would
% throw away the styles a library declares.
\makeatletter
\newcommand{\utpusediagrams}[1]{%
  \@for\utp@library:=#1\do{\input{../.shared/diagrams/\utp@library}}%
}
\makeatother

% --------------------------------------------------------------------
%  11.  FIGURE FRAMES
% --------------------------------------------------------------------
%  Three ways to put a drawing, chart or screenshot on the page:
%
%      \begin{utpfigure} ... \end{utpfigure}
%          a centred frame with no caption, for a diagram the text
%          talks about directly. An optional title puts a bar on top:
%          \begin{utpfigure}[Where the VLAN tag is added] ...
%
%      \begin{utptable} ... \end{utptable}
%          the same frame in teal, around a tabular or tabularx.
%
%      \begin{utpfig}[short entry]{The caption.}{fig:label}
%        \begin{tikzpicture} ... \end{tikzpicture}
%      \end{utpfig}
%          a frame with a numbered caption that \figref can point to.
%          The optional first argument is the short form for a list of
%          figures. A page never breaks between the drawing and its
%          caption.
%
%  A screenshot or slide lifted from the lecturer's material goes in a
%  titled frame, so the reader can tell it from a drawing made for the
%  notes; the optional width defaults to 0.9 and 0.78 of the line:
%
%      \sourceimage[0.8\linewidth]{Static-routing topology}{static-topology.png}
%      \compactsourceimage{The same, kept to a smaller height}{static-topology.png}
%
%  \fignote{...} adds a short muted line under a drawing, inside its frame.

\newtcolorbox{utpfigure}[1][]{enhanced,width=\linewidth,colback=white,colframe=UTPBlue!28, boxrule=0.7pt,arc=2mm, left=2.2mm,right=2.2mm,top=2.2mm,bottom=1.8mm, before skip=0.75em,after skip=0.85em, before upper={\centering}, halign=center, fonttitle=\bfseries\color{white},colbacktitle=UTPBlue,title={#1}}
\newtcolorbox{utptable}{enhanced,width=\linewidth,colback=white,colframe=UTPTeal!32, boxrule=0.7pt,arc=2mm, left=2.2mm,right=2.2mm,top=1.8mm,bottom=1.8mm, before skip=0.75em,after skip=0.85em, halign=center}
\newcommand{\utpbreakableframes}{%
  \renewtcolorbox{utpfigure}[1][]{enhanced,breakable,colback=white,colframe=UTPBlue!35, boxrule=0.5pt,arc=2mm,left=2.5mm,right=2.5mm,top=2.2mm,bottom=2.2mm, halign=center, fonttitle=\bfseries\color{white},colbacktitle=UTPBlue,title={##1}}%
  \renewtcolorbox{utptable}{enhanced,breakable,colback=white,colframe=UTPBlue!28, boxrule=0.45pt,arc=2mm,left=2.2mm,right=2.2mm,top=1.8mm,bottom=1.8mm}%
}

\captionsetup{
  hypcap=false,
  font={small},
  labelfont={bf,color=UTPNavy},
  textfont={color=UTPInk},
  labelsep=period,
  justification=centering,
  skip=4pt
}

\newenvironment{utpfig}[3][\relax]
  {\def\utpfigshort{#1}\def\utpfigcap{#2}\def\utpfiglab{#3}%
   \par\addvspace{0.55em}%
   \begin{tcolorbox}[enhanced,breakable,colback=white,colframe=UTPBlue!32,
     boxrule=0.5pt,arc=2mm,left=2.5mm,right=2.5mm,top=2.6mm,bottom=2.2mm]
   \centering}
  {\par\nopagebreak\vspace{0.45em}%
   % The caption is boxed: \captionof starts with a write and a zero
   % rule, and the glue after them would otherwise be a place where the
   % page could break, leaving the caption alone at the top of the next.
   % Inside the box the caption is told it sits below the drawing, which
   % it can no longer see, so the gap stays above it.
   \nointerlineskip
   \vbox{\setlength{\parskip}{0pt}\captionsetup{position=below}%
   \ifx\utpfigshort\relax
     \captionof{figure}{\utpfigcap}%
   \else
     \captionof{figure}[\utpfigshort]{\utpfigcap}%
   \fi
   \label{\utpfiglab}}%
   \end{tcolorbox}%
   \par\addvspace{0.55em}}

\newcommand{\sourceimage}[3][0.9\linewidth]{%
  \begin{utpfigure}[Source visual --- #2]%
    \includegraphics[width=#1,height=0.56\textheight,keepaspectratio]{#3}%
  \end{utpfigure}}
\newcommand{\compactsourceimage}[3][0.78\linewidth]{%
  \begin{utpfigure}[Source visual --- #2]%
    \includegraphics[width=#1,height=0.42\textheight,keepaspectratio]{#3}%
  \end{utpfigure}}

\newcommand{\fignote}[1]{\par\smallskip{\scriptsize\color{UTPMuted}#1\par}}

% --------------------------------------------------------------------
%  12.  LAB SHEETS
% --------------------------------------------------------------------
%  Every lab sheet opens with the banner, a full-width page break that
%  also puts the lab in the contents:
%
%      \labsection{3}{Scheduling Policies}{sec:lab03}
%
%  The "of N" comes from \utplabtotal, so a book states its lab count
%  once in its configuration rather than in every lab file.
%
%  The body of the sheet is a labproject box (section 8), and inside it
%  the same toolkit serves every course:
%
%      \labflow{chapter link}{goal}{what to explain}
%          the three-line brief under the lab title.
%      \labsourcecard{file}{what it does}{what to inspect}
%          the heading over a listing of one supplied source file.
%      \begin{labquestion}{3}{Short title} ... \end{labquestion}
%      \begin{labanswer} ... \expectedoutput ... \end{labanswer}
%          a question as the student meets it, then the worked code, the
%          output it should print, and what that output means.
%      \begin{checkanswers} ... \end{checkanswers}
%          answers to a quickcheck, set apart so they can be skipped.
%      labmode, verification, troubleshoot, labdeliverable (section 8)
%          how much is scaffolded, the evidence it worked, what to do if
%          it did not, and what to hand in.

\newcommand{\labsection}[3]{%
  \clearpage
  \phantomsection
  \addcontentsline{toc}{lab}{Lab #1 of \utplabtotal\space--- #2}%
  \markboth{Lab #1 of \utplabtotal\space--- #2}{}%
  \begin{tcolorbox}[enhanced,width=\linewidth,colback=UTPNavy,colframe=UTPNavy,
    boxrule=0pt,arc=2mm,left=5mm,right=5mm,top=5mm,bottom=5mm,
    before skip=0pt,after skip=1.2em]
    {\large\bfseries\color{UTPTeal!35!white} PRACTICAL MILESTONE\par}
    \vspace{0.25em}
    {\Huge\bfseries\color{white} Lab #1 of \utplabtotal\par}
    \vspace{0.20em}
    {\Large\bfseries\color{white} #2\par}
  \end{tcolorbox}
  \label{#3}%
}

\newcommand{\labflow}[3]{%
  \par\smallskip\noindent
  \textbf{Chapter connection.} #1\par
  \textbf{Goal.} #2\par
  \textbf{Explain.} #3\par\medskip}
\newcommand{\labsourcecard}[3]{%
  \par\medskip\noindent\textbf{#1 --- #2}\par
  \textbf{Inspect.} #3\par\smallskip}

\newtcolorbox{labquestion}[2]{enhanced,breakable,colback=UTPPaleOrange!55,colframe=UTPOrange,
  boxrule=0.8pt,arc=2mm,left=2.6mm,right=2.6mm,top=1.8mm,bottom=1.8mm,
  fonttitle=\bfseries\color{white},colbacktitle=UTPOrange,title={Question #1 --- #2}}
\newtcolorbox{labanswer}{enhanced,breakable,colback=white,colframe=UTPTeal,
  boxrule=0.8pt,arc=2mm,left=2.6mm,right=2.6mm,top=1.8mm,bottom=1.8mm,
  fonttitle=\bfseries\color{white},colbacktitle=UTPTeal,title=Worked answer}
\newcommand{\expectedoutput}{\par\smallskip\noindent\textbf{\color{UTPNavy}Reference output.}\nopagebreak\par\nopagebreak}
\utpnewcallout[boxrule=0.65pt,left=2.5mm,right=2.5mm,top=1.6mm,bottom=1.6mm,colbacktitle=UTPMuted]{checkanswers}{UTPPaleGray!72}{UTPMuted!62}{Quick-check answers}

% --------------------------------------------------------------------
%  13.  STAGE DIVIDER
% --------------------------------------------------------------------
%  A full-page divider used by the longer books to open a part:
%  \substage{3A}{Supervised Workflow}{one-line summary}.
%  Harmless to leave unused.

\newcommand{\substage}[3]{%
  \cleardoublepage
  \thispagestyle{empty}
  \phantomsection
  \addcontentsline{toc}{chapter}{Stage #1 --- #2}
  \begin{center}
    \vspace*{0.20\textheight}
    {\Large\bfseries\color{UTPTeal} Stage #1\par}
    \vspace{0.45em}
    {\Huge\bfseries\color{UTPNavy} #2\par}
    \vspace{0.65em}
    {\color{UTPTeal}\rule{0.46\linewidth}{1.5pt}\par}
    \vspace{1.1em}
    \begin{minipage}{0.76\linewidth}
      \centering\large\color{UTPMuted} #3
    \end{minipage}
  \end{center}
  \cleardoublepage
}

% --------------------------------------------------------------------
%  14.  FRONT MATTER
% --------------------------------------------------------------------
%  The opening every book shares: the cover, the contents in roman
%  numerals, the course map, then chapter one on an arabic page 1.
%
%      \begin{document}
%      \utpfrontmatter
%
%  What it prints comes from the configuration in section 1 --- the
%  cover text from \utpcoursetitle and \utpcourseoutline, the overview
%  chapter from \utpcoursemap --- and the contents page is set as
%  \utpcontentspage (section 5) describes. A book with no overview
%  chapter sets \def\utpcoursemap{} and the step is skipped.
%
%  Page anchors are off across the cover so hyperref does not warn
%  about the unnumbered pages, and back on before the contents links
%  are used.
%
%  The number of front-matter pages is recorded in the .aux file as
%  \utpfrontpages, which is how .shared/build.py turns the page numbers
%  in the contents into physical pages when it splits the book.

\newcommand{\utpfrontmatter}{%
  \hypersetup{pageanchor=false}%
  % ---------- Cover ----------
%  The title page, printed by \utpfrontmatter. Everything that changes
%  between books is a configuration macro from style/configuration.tex:
%  \utpcoursetitle, \utpcourseoutline, \utpoutlineheading, \utpauthor
%  and \utpauthorrole. The colours are the shared palette.

\begin{titlepage}
\thispagestyle{empty}
\begin{tikzpicture}[remember picture,overlay]
  \fill[UTPGold]
    (current page.north west) rectangle
    ([yshift=-0.12cm]current page.north east);
  \fill[UTPGold]
    (current page.south west) rectangle
    ([yshift=0.12cm]current page.south east);
\end{tikzpicture}

\centering
\includegraphics[
  height=2.85cm,
  clip
]{logo-UTP.png}\par

\vspace{0.50cm}
{\fontsize{29}{33}\selectfont\bfseries\color{UTPNavy}\utpcoursetitle\par}
\vspace{0.24cm}
{\Large\color{UTPMuted} \utpdocumentsubtitle\par}
\vspace{0.42cm}
{\color{UTPGold}\rule{0.26\linewidth}{1.3pt}\par}
\vspace{0.72cm}
\begin{tikzpicture}[
  core/.style={circle,draw=UTPGold,fill=UTPNavy,minimum size=1.95cm,
    align=center,font=\small\bfseries\color{white},line width=1pt},
  stage/.style={rounded corners=1.5mm,draw=UTPBlue!65,fill=UTPPaleBlue!42,
    text width=3.15cm,minimum height=1.42cm,align=center,
    font=\small\bfseries\color{UTPNavy},inner xsep=3mm,
    inner ysep=3.5mm,line width=0.65pt},
  cycle/.style={-{Stealth[length=2.1mm]},semithick,draw=UTPMuted,
    rounded corners=4mm,shorten <=7pt,shorten >=7pt}
]
\node[core] (guide) at (0,-2.55) {TA\\WORKFLOW};
\node[stage] (prepare) at (0,0) {SESSION\\PLANNING};
\node[stage] (explain) at (4.75,-2.55) {INSTRUCTIONAL\\SUPPORT};
\node[stage] (facilitate) at (0,-5.10) {STUDENT\\ENGAGEMENT};
\node[stage] (assess) at (-4.75,-2.55) {LEARNING\\ASSESSMENT};

\draw[cycle] (prepare.east) -- (explain.north);
\draw[cycle] (explain.south) -- (facilitate.east);
\draw[cycle] (facilitate.west) -- (assess.south);
\draw[cycle] (assess.north) -- (prepare.west);
\end{tikzpicture}

\vspace{0.72cm}
\begin{tcolorbox}[
  enhanced,width=0.82\linewidth,colback=white,colframe=UTPBlue!18,
  arc=1.5mm,boxrule=0.45pt,left=6mm,right=6mm,top=3.2mm,bottom=3.2mm,
  borderline west={1.6pt}{0pt}{UTPGold}]
\centering\normalsize\color{UTPInk}
{\large\bfseries\color{UTPBlue}\utpoutlineheading\par}
\vspace{0.3em}
\utpcourseoutline
\end{tcolorbox}

\vfill
\begin{minipage}{0.82\textwidth}
\begin{tabular}{@{}>{\RaggedRight\arraybackslash}m{0.42\linewidth}@{\hspace{0.10\linewidth}}>{\RaggedLeft\arraybackslash}m{0.48\linewidth}@{}}
{\normalsize\bfseries\color{UTPBlue}PREPARED BY\par\vspace{0.25em}
 \Large\bfseries\color{UTPNavy}\utpauthor\par\vspace{0.12em}
 \normalfont\small\color{UTPMuted}\utpauthorrole}
&
\includegraphics[height=2.10cm,keepaspectratio]{logo-DC.png}
\end{tabular}
\end{minipage}
\vspace{0.5cm}
\end{titlepage}%
  \cleardoublepage
  \pagenumbering{roman}%
  \setcounter{tocdepth}{0}%
  \utptableofcontents
  \newpage
  \hypersetup{pageanchor=true}%
  \ifx\utpcoursemap\empty\else\input{\utpcoursemap}\fi
  \cleardoublepage
  \ifdefined\ReadonlyShipoutCounter
    \immediate\write\csname @auxout\endcsname
      {\string\gdef\string\utpfrontpages{\the\ReadonlyShipoutCounter}}%
  \fi
  \pagenumbering{arabic}%
}

% --------------------------------------------------------------------
%  15.  LISTING STYLES
% --------------------------------------------------------------------
%  Two base looks, and one named style per language built on them, so
%  the colours and frame are the same in every book:
%
%      utpcode          the base for source code
%      utpterminal      the base for commands typed and output printed
%      utppython        Python
%      utpr             R
%      utpc             C
%      utpcpp           C++
%      utpbash          shell scripts and commands
%      utpcnumbered     C with line numbers, for a complete source file
%      utpbashnumbered  shell with line numbers, likewise
%
%      \begin{lstlisting}[style=utppython] ... \end{lstlisting}
%      \lstinputlisting[style=utpcnumbered]{all_experiments/Lab04/fcfs.c}
%
%  A book that sets most of its listings in one style makes it the
%  default once, after the input:  \lstset{style=utpcpp}

\lstdefinestyle{utpcode}{
  basicstyle=\ttfamily\scriptsize,
  keywordstyle=\bfseries\color{UTPBlue},
  commentstyle=\itshape\color{UTPMuted},
  stringstyle=\color{UTPGreen!70!black},
  backgroundcolor=\color{UTPPaleGray!72},
  frame=single,
  rulecolor=\color{UTPBlue!28},
  framesep=5pt,
  breaklines=true,
  breakatwhitespace=false,
  showstringspaces=false,
  keepspaces=true,
  columns=fullflexible,
  upquote=true,
  tabsize=4,
  aboveskip=0.7em,
  belowskip=0.8em
}
\lstdefinestyle{utpterminal}{
  basicstyle=\ttfamily\scriptsize,
  backgroundcolor=\color{UTPPaleGray!72},
  frame=single,
  rulecolor=\color{UTPMuted!35},
  framesep=5pt,
  breaklines=true,
  showstringspaces=false,
  keepspaces=true,
  columns=fullflexible,
  upquote=true,
  aboveskip=0.6em,
  belowskip=0.7em
}

\lstdefinestyle{utppython}{style=utpcode,language=Python}
\lstdefinestyle{utpr}{style=utpcode,language=R,tabsize=2}
\lstdefinestyle{utpc}{style=utpcode,language=C}
\lstdefinestyle{utpcpp}{style=utpcode,language=C++}
\lstdefinestyle{utpbash}{style=utpcode,language=bash,
  rulecolor=\color{UTPMuted!35},aboveskip=0.6em,belowskip=0.7em}
\lstdefinestyle{utpcnumbered}{style=utpc,numbers=left,numberstyle=\tiny\color{UTPMuted},
  numbersep=7pt,xleftmargin=1.7em,framexleftmargin=1.4em}
\lstdefinestyle{utpbashnumbered}{style=utpbash,numbers=left,numberstyle=\tiny\color{UTPMuted},
  numbersep=7pt,xleftmargin=1.7em,framexleftmargin=1.4em}

% Pseudocode. A book that sets algorithms asks for them once, after the
% input, and gets algorithm2e loaded and coloured to match:
%
%     \utpusealgorithms
%
% \algostyle then sets one algorithm environment a size smaller than the
% text around it.
\newcommand{\utpalgorithmstyle}{%
  \SetAlCapNameFnt{\bfseries\color{UTPNavy}}%
  \SetAlCapFnt{\bfseries\color{UTPNavy}}%
  \SetAlgoCaptionSeparator{.}%
  \SetKwInput{KwIn}{Input}%
  \SetKwInput{KwOut}{Output}%
  \SetNlSty{color}{UTPMuted}{}%
  \SetCommentSty{\itshape\color{UTPMuted}}%
  \SetKwSty{bfseries\color{UTPBlue}}%
  \SetArgSty{}%
  \providecommand{\algostyle}{\small\SetAlgoNlRelativeSize{-1}\DontPrintSemicolon\SetAlgoVlined}%
}
\newcommand{\utpusealgorithms}{%
  \usepackage[linesnumbered,ruled,vlined]{algorithm2e}%
  \utpalgorithmstyle
}


% Use the full Learning Path page without pushing entries onto another page.
\newcommand{\osfullpagetoc}{%
  \begingroup
  \normalsize
  \setlength{\parskip}{0pt}%
  \setlength{\cftbeforepartskip}{1.15em}%
  \setlength{\cftbeforechapskip}{0.38em}%
  \setlength{\cftbeforelabskip}{0.30em}%
  \renewcommand{\cftpartfont}{\large\bfseries\color{UTPNavy}}%
  \renewcommand{\cftpartpagefont}{\large\bfseries\color{UTPNavy}}%
  \renewcommand{\cftchapfont}{\normalsize}%
  \renewcommand{\cftchappagefont}{\normalsize}%
  \renewcommand{\cftlabfont}{\normalsize\bfseries\color{UTPOrange!90!black}}%
  \renewcommand{\cftlabpagefont}{\normalsize\bfseries\color{UTPOrange!90!black}}%
  \tableofcontents
  \endgroup
}
\patchcmd{\utpfrontmatter}
  {\begingroup\small\setlength{\parskip}{0pt}\tableofcontents\endgroup}
  {\osfullpagetoc}{}{}

% ---------- course-only settings ----------
% The shared file already defines osc and osbash. These variants simply
% add line numbers when the guide prints complete experiment source files.
\lstdefinestyle{oslabc}{style=osc,numbers=left,numberstyle=\tiny\color{UTPMuted},
  numbersep=7pt,xleftmargin=1.7em,framexleftmargin=1.4em}
\lstdefinestyle{oslabsh}{style=osbash,numbers=left,numberstyle=\tiny\color{UTPMuted},
  numbersep=7pt,xleftmargin=1.7em,framexleftmargin=1.4em}
\utpdiagramstyles{os}{1.05cm}{0.80cm}

% Short lab helpers kept deliberately lightweight so the practical pages
% follow the same visual rhythm as the Machine Learning guide.
\newcommand{\ossourceintro}[3]{%
  \par\medskip\noindent\textbf{#1 --- #2}\par
  \textbf{Inspect.} #3\par\smallskip}

\newcommand{\ossourcecard}[3]{\ossourceintro{#1}{#2}{#3}}
\newcommand{\oslabflow}[3]{%
  \par\smallskip\noindent
  \textbf{Chapter connection.} #1\par
  \textbf{Goal.} #2\par
  \textbf{Explain.} #3\par\medskip}

\begin{document}
\utpfrontmatter

\part{Orientation and Tools}
\chapter*{Your Linux Toolkit and Environment Setup}
\addcontentsline{toc}{chapter}{Your Linux Toolkit and Environment Setup}
\markboth{Your Linux Toolkit and Environment Setup}{}
\chapterguide{No Linux experience is required. You need a computer where you can use Linux, WSL, or a Linux virtual machine.}{identify the course folder, open a Linux shell, edit source in Vim, compile with GCC, run a program, and prepare Git/SSH for saving course work.}

The practicals use a \textbf{command-line workflow}. Linux is the reference environment because the processes, files, permissions, IPC objects, and other OS concepts used in the course are easy to observe there.

\section*{The course project layout}
Keep the supplied structure unchanged. The lab pages print source directly from \texttt{all\_experiments/}, so the files in the book and the files you edit are the same files.

\begin{lstlisting}[style=osbash]
OPERATING SYSTEMS/
|-- Lab01/
|   |-- directory_list.c
|   |-- fork_demo.c
|-- Lab02/
|   |-- file_copy.c
|   |-- file_info.c
|   |-- mini_grep.c
|   |-- sample.txt
|-- Lab03/
|   |-- events.log
|   |-- log_summary.sh
|   |-- process_summary.sh
|   |-- system_report.sh
|-- Lab04/
|   |-- fcfs.c
|   |-- priority.c
|   |-- round_robin.c
|   |-- sjf.c
|-- Lab05/
|   |-- pipe_demo.c
|   |-- producer_consumer_sim.c
|   |-- shm_reader.c
|   |-- shm_writer.c
|-- Lab06/
|   |-- banker.c
|   |-- threads_demo.c
|-- Lab07/
|   |-- best_fit.c
|   |-- first_fit.c
|   |-- worst_fit.c
|-- Lab08/
|   |-- fifo.c
|   |-- lfu.c
|   |-- lru.c
|-- Lab09/
|   |-- directory_walk.c
|   |-- make_demo_tree.c
|-- Lab10/
|   |-- binary_records.c
|   |-- indexed_allocation.c
|   |-- text_records.c
|-- README.md
\end{lstlisting}

\section*{Choose the Linux environment}
\begin{mltable}
\renewcommand{\arraystretch}{1.18}
\begin{tabularx}{\linewidth}{@{}
  >{\raggedright\arraybackslash}p{0.14\linewidth}
  >{\raggedright\arraybackslash}p{0.32\linewidth}
  >{\raggedright\arraybackslash}X@{}}
\toprule
\rowcolor{UTPPaleBlue!92}
\textbf{Platform} & \textbf{Recommended route} & \textbf{What you use in the labs} \\
\midrule
Linux & Native Linux & Terminal, Vim, GCC, Bash \\
Windows & WSL 2 with Ubuntu & Ubuntu/WSL terminal \\
macOS & macOS Terminal & Terminal and compiler \\
\bottomrule
\end{tabularx}
\end{mltable}

\section*{Windows students: install WSL}
WSL lets Windows run a Linux environment alongside Windows. Your Windows desktop remains unchanged; the \textbf{Ubuntu/WSL terminal} becomes the laboratory environment.

Open \textbf{PowerShell or Windows Terminal as Administrator}:
\begin{lstlisting}[style=osbash]
wsl --install
\end{lstlisting}
Restart if Windows asks. Open \textbf{Ubuntu}, create a Linux username and password, then verify from PowerShell:
\begin{lstlisting}[style=osbash]
wsl --status
wsl --list --verbose
\end{lstlisting}

\begin{mlfigure}
\begin{tikzpicture}[
  box/.style={draw=UTPBlue,fill=UTPPaleBlue,rounded corners=2mm,
              text width=4.25cm,minimum height=1.55cm,inner sep=3mm,
              align=center,font=\small\color{UTPNavy}},
  linux/.style={box,draw=UTPTeal,fill=UTPPaleTeal},
  a/.style={-{Stealth[length=2.2mm]},thick,draw=UTPMuted}
]
\node[box] (win) at (-4.0,0) {\textbf{Windows PowerShell}\\[2pt]
  install and check WSL\\\texttt{wsl --status}};
\node[linux] (wsl) at (4.0,0) {\textbf{Ubuntu / WSL shell}\\[2pt]
  run course commands\\\texttt{gcc}, \texttt{vim}, \texttt{ps}, \texttt{bash}};
\draw[a] (win) -- node[above=3pt,fill=white,inner sep=1.5pt,
  text width=2.6cm,align=center,font=\scriptsize\color{UTPMuted}]
  {launches the\\Linux environment} (wsl);
\end{tikzpicture}
\end{mlfigure}

\begin{pitfall}
Do not type Linux package, compiler, or lab commands into PowerShell. After WSL is installed, run the practical work inside Ubuntu/WSL unless the guide explicitly says \textbf{PowerShell}.
\end{pitfall}

\section*{Install the command-line tools}
On Ubuntu, Debian, or WSL:
\begin{lstlisting}[style=osbash]
sudo apt update
sudo apt install build-essential gdb make vim git openssh-client
\end{lstlisting}
Verify the core tools:
\begin{lstlisting}[style=osbash]
gcc --version
vim --version | head -n 1
git --version
uname -a
\end{lstlisting}

For WSL, keep Linux course work in the Linux home directory rather than under \texttt{/mnt/c}:
\begin{lstlisting}[style=osbash]
mkdir -p ~/operating-systems
cd ~/operating-systems
pwd
\end{lstlisting}

\section*{The CLI commands you need first}
\begin{mltable}
\renewcommand{\arraystretch}{1.18}
\begin{tabularx}{\linewidth}{@{}
  >{\raggedright\arraybackslash}p{0.16\linewidth}
  >{\raggedright\arraybackslash}p{0.28\linewidth}
  >{\raggedright\arraybackslash}X@{}}
\toprule
\rowcolor{UTPPaleBlue!92}
\textbf{Command} & \textbf{Purpose} & \textbf{Example} \\
\midrule
\texttt{pwd} & show the current directory & \texttt{pwd} \\
\texttt{ls} & list files & \texttt{ls -la} \\
\texttt{cd} & change directory & \texttt{cd all\_experiments/Lab01} \\
\texttt{mkdir} & create a directory & \texttt{mkdir practice} \\
\texttt{cp} / \texttt{mv} & copy / move or rename & \texttt{cp a.c b.c}, \texttt{mv old new} \\
\texttt{rm} & remove a file & \texttt{rm program} \\
\texttt{man} & open the manual & \texttt{man ps} \\
\bottomrule
\end{tabularx}
\end{mltable}

\section*{Vim: the editor used in this course}
Vim is a terminal editor. You do not need advanced Vim knowledge; the course uses a small, repeatable set of commands.

\begin{mlfigure}
\begin{tikzpicture}[
  n/.style={draw=UTPBlue,fill=UTPPaleBlue,rounded corners=2mm,
            text width=4.25cm,minimum height=1.55cm,inner sep=3.5mm,
            align=center,font=\small\color{UTPNavy}},
  i/.style={n,draw=UTPTeal,fill=UTPPaleTeal,minimum height=2.05cm},
  c/.style={n,draw=UTPOrange,fill=UTPPaleOrange,minimum height=2.05cm},
  bileft/.style={{Stealth[length=2.5mm,width=2mm]}-{Stealth[length=2.5mm,width=2mm]},
                line width=1.05pt,draw=UTPTeal!85!black,
                shorten <=2pt,shorten >=2pt},
  biright/.style={{Stealth[length=2.5mm,width=2mm]}-{Stealth[length=2.5mm,width=2mm]},
                 line width=1.05pt,draw=UTPOrange!90!black,
                 shorten <=2pt,shorten >=2pt}
]
\node[n] (normal) at (0,2.65) {\textbf{NORMAL MODE}\\[3pt]
  navigate, edit, search, and issue commands};
\node[i] (insert) at (-4.25,-0.75) {\textbf{INSERT MODE}\\[3pt]
  type and edit program text\\[5pt]
  {\scriptsize\textbf{Enter:} \texttt{i}\quad\textbf{Return:} \texttt{Esc}}};
\node[c] (command) at (4.25,-0.75) {\textbf{COMMAND-LINE MODE}\\[3pt]
  save, quit, and change settings\\[5pt]
  {\scriptsize\textbf{Enter:} \texttt{:}\quad\textbf{Return:} \texttt{Enter/Esc}}};

\draw[bileft] (normal.south west)
  to[out=235,in=55,looseness=1.08] (insert.north east);
\draw[biright] (normal.south east)
  to[out=305,in=125,looseness=1.08] (command.north west);
\end{tikzpicture}
\end{mlfigure}

Open a file from the shell:
\begin{lstlisting}[style=osbash]
vim hello.c
\end{lstlisting}

\begin{mltable}
\renewcommand{\arraystretch}{1.18}
\begin{tabularx}{\linewidth}{@{}
  >{\raggedright\arraybackslash}p{0.22\linewidth}
  >{\raggedright\arraybackslash}X@{}}
\toprule
\rowcolor{UTPPaleBlue!92}
\textbf{Vim command} & \textbf{What it does} \\
\midrule
\texttt{i} & enter Insert mode and start typing \\
\texttt{Esc} & return to Normal mode \\
\texttt{:w} & save \\
\texttt{:q} & quit if nothing changed \\
\texttt{:wq} & save and quit \\
\texttt{:q!} & quit and discard unsaved changes \\
\texttt{u} / \texttt{Ctrl+r} & undo / redo \\
\texttt{/word} & search forward for \texttt{word}; press \texttt{n} for the next match \\
\texttt{:set number} & show line numbers \\
\texttt{:25} & jump to line 25, useful for compiler errors \\
\bottomrule
\end{tabularx}
\end{mltable}

For guided practice, run \texttt{vimtutor} once before Lab 1.

\section*{The edit-compile-run cycle}
Create a first program with \texttt{vim hello.c}:
\begin{lstlisting}[style=osc]
#include <stdio.h>

int main(void) {
    printf("My OS lab environment is ready.\n");
    return 0;
}
\end{lstlisting}
Compile and run:
\begin{lstlisting}[style=osbash]
gcc -Wall -Wextra -std=c11 hello.c -o hello
./hello
\end{lstlisting}

\begin{mlfigure}
\begin{tikzpicture}[
  flowcard/.style={draw=UTPBlue!65,fill=white,rounded corners=1.5mm,
                   text width=2.45cm,minimum height=1.35cm,
                   inner xsep=2mm,inner ysep=2.5mm,align=center,
                   font=\small\color{UTPNavy}},
  forward/.style={-{Stealth[length=2.3mm,width=1.8mm]},line width=0.9pt,
                  draw=UTPTeal,shorten <=2pt,shorten >=2pt},
  feedback/.style={-{Stealth[length=2.2mm,width=1.7mm]},semithick,dashed,
                   draw=UTPMuted,rounded corners=4pt,
                   shorten <=2pt,shorten >=2pt},
  feedbacklabel/.style={fill=white,inner xsep=2mm,inner ysep=1pt,
                       font=\scriptsize\color{UTPMuted}}
]
\node[flowcard] (edit) at (-5.2,0) {\textcolor{UTPBlue}{\textbf{01}}\\[2pt]
  \textbf{EDIT SOURCE}\\[3pt]\texttt{vim file.c}};
\node[flowcard] (save) at (-1.75,0) {\textcolor{UTPBlue}{\textbf{02}}\\[2pt]
  \textbf{SAVE}\\[3pt]\texttt{:wq}};
\node[flowcard] (build) at (1.75,0) {\textcolor{UTPBlue}{\textbf{03}}\\[2pt]
  \textbf{COMPILE}\\[3pt]\texttt{gcc ...}};
\node[flowcard] (run) at (5.2,0) {\textcolor{UTPBlue}{\textbf{04}}\\[2pt]
  \textbf{RUN AND INSPECT}\\[3pt]\texttt{./program}};

\draw[forward] (edit.east) -- (save.west);
\draw[forward] (save.east) -- (build.west);
\draw[forward] (build.east) -- (run.west);

\draw[feedback] (run.south) -- ++(0,-1.0)
  -- node[feedbacklabel,below]{revise one thing, then repeat} ++(-10.4,0)
  -- (edit.south);
\end{tikzpicture}
\end{mlfigure}

\section*{Git and SSH setup for course work}
Git records changes to your source files. SSH lets Git authenticate to GitHub or another Git host without using an account password for every push.

\subsection*{1. Install Git and OpenSSH}
\begin{mltable}
\renewcommand{\arraystretch}{1.18}
\begin{tabularx}{\linewidth}{@{}
  >{\raggedright\arraybackslash}p{0.16\linewidth}
  >{\raggedright\arraybackslash}X@{}}
\toprule
\rowcolor{UTPPaleBlue!92}
\textbf{System} & \textbf{Command} \\
\midrule
Linux & \texttt{sudo apt update \&\& sudo apt install git openssh-client} \\
WSL & Run the same Linux command \textbf{inside Ubuntu/WSL}, not PowerShell \\
macOS & \texttt{brew install git} then \texttt{git --version} \\
\bottomrule
\end{tabularx}
\end{mltable}

Confirm:
\begin{lstlisting}[style=osbash]
git --version
ssh -V
\end{lstlisting}

\subsection*{2. Configure your commit identity}
Use the name and email that should appear on your course commits:
\begin{lstlisting}[style=osbash]
git config --global user.name "Your Name"
git config --global user.email "you@example.com"
git config --global init.defaultBranch main
git config --global --list
\end{lstlisting}

\subsection*{3. Create an SSH key}
First check whether a key already exists:
\begin{lstlisting}[style=osbash]
ls -al ~/.ssh
\end{lstlisting}
If \texttt{id\_ed25519} and \texttt{id\_ed25519.pub} are absent, create them:
\begin{lstlisting}[style=osbash]
ssh-keygen -t ed25519 -C "you@example.com"
\end{lstlisting}
Press Enter to accept the default location. A passphrase is recommended.

Load the key:
\begin{lstlisting}[style=osbash]
eval "$(ssh-agent -s)"
ssh-add ~/.ssh/id_ed25519
\end{lstlisting}
On macOS, \texttt{ssh-add --apple-use-keychain \string~/.ssh/id\_ed25519} may be used instead.

\subsection*{4. Copy only the public key}
\begin{mltable}
\renewcommand{\arraystretch}{1.18}
\begin{tabularx}{\linewidth}{@{}
  >{\raggedright\arraybackslash}p{0.16\linewidth}
  >{\raggedright\arraybackslash}X@{}}
\toprule
\rowcolor{UTPPaleBlue!92}
\textbf{System} & \textbf{Copy command} \\
\midrule
Linux & \texttt{cat \string~/.ssh/id\_ed25519.pub} and copy the complete printed line \\
WSL & \texttt{clip.exe < \string~/.ssh/id\_ed25519.pub} \\
macOS & \texttt{pbcopy < \string~/.ssh/id\_ed25519.pub} \\
\bottomrule
\end{tabularx}
\end{mltable}

In GitHub, open \textbf{Settings $>$ SSH and GPG keys $>$ New SSH key}, give the key a recognizable name, and paste the public key.

\begin{pitfall}
Never share or commit \texttt{\string~/.ssh/id\_ed25519}. That is the \textbf{private key}. Only the file ending in \texttt{.pub} is copied to the Git hosting service.
\end{pitfall}

\subsection*{5. Test the SSH connection}
For GitHub:
\begin{lstlisting}[style=osbash]
ssh -T git@github.com
\end{lstlisting}
The first connection may ask you to confirm the host fingerprint. For GitLab, use \texttt{git@gitlab.com}.

\subsection*{6. Connect the course folder}
For a new empty GitHub repository, run these commands from the course root and replace \texttt{USERNAME/REPOSITORY}:
\begin{lstlisting}[style=osbash]
git init -b main
git add .
git status
git commit -m "Start operating-systems course work"
git remote add origin git@github.com:USERNAME/REPOSITORY.git
git push -u origin main
\end{lstlisting}

For later work:
\begin{lstlisting}[style=osbash]
git status
git add .
git commit -m "Complete Lab 1 process observations"
git push
\end{lstlisting}



\chapter{Operating Systems at a Glance}\label{ch:overview}
\chapterguide{You should be comfortable with basic C programs and the idea that a computer has a CPU, memory, and storage.}{Explain why an OS exists, distinguish kernel and user mode, classify OS services, and connect the main course topics into one execution story.}

\section{Start with something you have already done}
Open a terminal and type \texttt{ls}. A list of filenames appears. Nothing about that is obvious once you look closely:

\begin{itemize}
  \item Your shell was running. It had to \emph{stop} doing that and start a different program.
  \item \texttt{ls} needed memory --- but so does every other program currently running, and none of them interfered.
  \item The filenames live as magnetic or flash state on a disk. Something turned that into text.
  \item The output reached your screen, which \texttt{ls} does not own and cannot address directly.
\end{itemize}

Four different subsystems cooperated, and you typed two letters. The \term{operating system (OS)} is the software that made all four happen.

\begin{keyidea}
An OS does two jobs at once. It is an \emph{abstraction layer}, turning awkward hardware into workable objects like processes, files, and memory addresses. It is also a \emph{resource manager}, deciding who gets what, for how long, and with what permission.

Nearly every design decision in this course comes from those two jobs pulling in different directions.
\end{keyidea}

\section{The layers involved}
Your program sits at the top. The hardware sits at the bottom. Everything between is the OS's responsibility.

\begin{mlfigure}
\begin{tikzpicture}[
  layer/.style={draw=MLBlue,rounded corners=2mm,minimum width=11.5cm,minimum height=1.0cm,align=center,font=\bfseries\color{MLNavy}},
  node distance=0.25cm
]
\node[layer,fill=MLPaleBlue] (apps) {Applications: shell, editor, browser, compiler, user programs};
\node[layer,fill=MLPaleTeal,below=of apps] (sys) {System-call interface and libraries};
\node[layer,fill=MLPaleOrange,below=of sys] (kernel) {Kernel: processes, scheduler, memory manager, file systems, device drivers};
\node[layer,fill=MLPaleGray,below=of kernel] (hw) {Hardware: CPU, RAM, storage, timers, network and I/O devices};
\end{tikzpicture}
\end{mlfigure}

Read the diagram as a rule about \emph{reach}. An application cannot skip a layer. It cannot address the disk controller, and it cannot read another process's memory. Every downward step goes through the interface directly beneath it.

\section{Kernel mode and user mode}
That restriction is enforced by the CPU itself, not by good manners.

Modern processors run in at least two privilege levels. Your programs run in \term{user mode}, where privileged instructions and direct hardware access are simply unavailable --- attempting them traps instead of executing. The kernel runs in \term{kernel mode}, where everything is permitted.

Control crosses that boundary in exactly three ways, and it is worth memorising them now because the rest of the course keeps returning to them:

\begin{description}
  \item[System call] The program \emph{asks}. It wants something only the kernel can do.
  \item[Exception] The current instruction \emph{cannot proceed} --- divide by zero, bad memory access, page fault.
  \item[Interrupt] Hardware \emph{demands attention} --- a timer tick, a keypress, a disk finishing a read.
\end{description}

\begin{intuition}
User mode is a blast wall. A buggy program should be free to crash itself, and unable to take down the machine, corrupt another process's memory, or reprogram a disk controller on its way out.

This is why an infinite loop in your code does not freeze the whole computer. The timer interrupt fires regardless of what your program is doing, and control returns to the kernel whether your program likes it or not.
\end{intuition}

\section{System calls are the operating-system API}
A \term{system call} is a request for a service that needs kernel authority. They group into a handful of families:

\begin{mltable}
\begin{tabularx}{\linewidth}{@{}>{\bfseries\leavevmode\color{MLNavy}}p{0.22\linewidth}X p{0.26\linewidth}@{}}
\toprule
\rowcolor{MLPaleBlue!92}
Family & What it does & Seen in this course \\
\midrule
Process control & Create, terminate, wait for, or replace a running program & \texttt{fork()}, \texttt{waitpid()} in Lab 1 \\
File management & Open, read, write, seek, close, inspect metadata & \texttt{fopen()}, \texttt{fgetc()} in Lab 2 \\
Device management & Read and write devices, configure device state & Chapters 9--10 \\
Information & Process identifiers, time, resource limits & \texttt{getpid()} in Lab 1 \\
Communication & Pipes, shared memory, message queues, sockets, signals & Lab 5 \\
Protection & Permissions, identities, access checks & Chapter 2 \\
\bottomrule
\end{tabularx}
\end{mltable}

\begin{pitfall}
\texttt{printf()} is not a system call. It is a C library function that formats your text and then calls \texttt{write()}, which is the system call.

The distinction matters more than it sounds. The library buffers output, so \texttt{printf()} may return long before anything reaches the screen. In Lab 1 you fork a process, and both parent and child inherit a copy of that buffer --- which is how a single \texttt{printf()} can produce output twice. Add \texttt{fflush(stdout)} before \texttt{fork()} and the duplicate disappears.
\end{pitfall}

\section{One command, every subsystem}
Return to \texttt{ls}, and trace it properly. The shell creates a process. The scheduler decides when it runs. Virtual memory hands it an address space that appears to start at zero regardless of where it physically sits. The file system turns the pathname into a location. The storage stack reads blocks. An interrupt announces the read finished. The terminal driver receives the output.

\begin{keyidea}
The topics in this course are not a list of separate features. They are cooperating subsystems, and almost every non-trivial program execution needs all of them at once.

Each later chapter zooms in on one of the steps above.
\end{keyidea}

\section{Program, process, and thread}
These three words get used loosely in conversation and precisely in operating systems. The difference is worth getting right early.

A \term{program} is passive: instructions and data sitting in a file. It does nothing. A \term{process} is a running instance of a program, with state --- registers, program counter, address space, open files, scheduling information, credentials. A \term{thread} is a single flow of execution inside a process.

\begin{analogy}
A program is a recipe. A process is one cook actually working from it, with their own bench space, ingredients, and half-finished pan. A thread is a pair of hands.

Two cooks working from the same recipe do not share a bench --- that is two processes. Two pairs of hands belonging to one cook share everything on that bench --- that is two threads, and it is exactly why threads need locks and processes usually do not.
\end{analogy}

\begin{mltable}
\begin{tabularx}{\linewidth}{@{}>{\bfseries\leavevmode\color{MLNavy}}p{0.20\linewidth}XX@{}}
\toprule
\rowcolor{MLPaleBlue!92}
Concept & Owns / represents & Typical question \\
\midrule
Program & Executable instructions and static data & What code could run? \\
Process & Resource container plus execution state & Which running instance is this? \\
Thread & Schedulable execution flow & Which sequence of instructions runs now? \\
\bottomrule
\end{tabularx}
\end{mltable}

Running the same program twice gives two processes. They have different process IDs, separate memory, and no ability to see each other's variables. You will observe this directly in Lab 1: after \texttt{fork()}, parent and child both continue from the same line with the same variables --- and changing a variable in one has no effect on the other.

\section{The process-state model}
A process is not always running. Most of the time it is waiting for something. The standard five-state model captures that:

\begin{mlfigure}
\begin{tikzpicture}[
  st/.style={draw=MLBlue,fill=MLPaleBlue,rounded corners=2mm,minimum width=2.1cm,minimum height=0.8cm,align=center,font=\small\bfseries\color{MLNavy}},
  arr/.style={-{Stealth[length=2mm]},draw=MLTeal,thick}
]
\node[st] (new) at (0,0) {New};
\node[st] (ready) at (3.6,0) {Ready};
\node[st] (run) at (7.2,0) {Running};
\node[st] (term) at (11.0,0) {Terminated};
\node[st] (wait) at (7.2,-2.3) {Waiting};

\draw[arr] (new)--node[above,font=\scriptsize]{admit}(ready);
\draw[arr] (ready)--node[above,font=\scriptsize]{dispatch}(run);
\draw[arr] (run)--node[above,font=\scriptsize]{exit}(term);
\draw[arr] (run)--node[right,font=\scriptsize]{wait for I/O}(wait);

\draw[arr] (run.south west) to[bend left=38]
  node[below,font=\scriptsize]{preempt} (ready.south east);
\coordinate (eventturn) at (3.6,-2.3);
\draw[arr] (wait.west) -- node[below,font=\scriptsize]{event done}
  (eventturn) -- (ready.south);
\end{tikzpicture}
\end{mlfigure}

Two arrows deserve attention, because students routinely confuse them.

\emph{Running to Ready} is \term{preemption}. The process was perfectly able to continue; the OS stopped it anyway, usually because its time slice expired. This is the arrow Round Robin in Chapter 4 depends on.

\emph{Running to Waiting} is \term{blocking}. The process cannot continue --- it asked for data from disk and the data is not there yet. No scheduling policy can help; only the event can.

\begin{pitfall}
``Ready'' does not mean idle. A ready process wants the CPU right now and cannot have it, because another process is using it. On a busy machine, the ready queue is where processes spend most of their existence.
\end{pitfall}

\begin{quickcheck}
\begin{enumerate}
  \item Why does an application need a system call to create another process?
  \item What do threads in one process share, and what stays private to each?
  \item Name one event that moves a running process into the waiting state.
  \item A process sits in Ready for 200\,ms. Is it blocked? What is it waiting for?
\end{enumerate}

\medskip\noindent\textbf{Answers.}
\begin{enumerate}
  \item Creating a process means allocating an address space, a process-table entry, and an identity. All of those are kernel-owned structures, and letting user code write them directly would remove every protection guarantee the OS makes.
  \item Shared: address space, open-file table, code, globals, and the heap. Private: registers, program counter, and stack. Sharing the heap is exactly why two threads writing one variable need a lock.
  \item A read from disk, waiting on a child with \texttt{waitpid()}, or blocking on a semaphore --- anything where the process cannot proceed until an external event occurs.
  \item Not blocked. Ready means it is able to run and is waiting only for a turn on the CPU. A blocked process waits for an event; a ready one waits for the scheduler.
\end{enumerate}
\end{quickcheck}


\part{Processes and the Command Line}
\chapter{Processes and the Linux Process Namespace}\label{ch:processes}
\chapterguide{The OS primer and the basic terminal workflow.}{distinguish a program from a process, read PID/PPID relationships, explain what \texttt{fork()} creates, and relate directory names to the Linux namespace before Lab 1.}

\section{The picture to remember}
A \term{program} is code stored in a file. A \term{process} is a running instance of that code with its own PID, memory, open resources, and execution state.

\begin{mlfigure}
\begin{tikzpicture}[
  box/.style={draw=UTPBlue,fill=UTPPaleBlue,rounded corners=2mm,minimum width=3.5cm,minimum height=1cm,align=center,font=\small\color{UTPNavy}},
  live/.style={box,draw=UTPTeal,fill=UTPPaleTeal},
  kern/.style={box,draw=UTPNavy,fill=UTPNavy!10},
  ar/.style={-{Stealth[length=2.2mm]},thick,draw=UTPOrange}]
\node[box] (file) at (-4.2,0) {Executable file\\\texttt{fork\_demo}};
\node[live] (proc) at (0,0) {Running process\\PID 4200};
\node[kern] (os) at (4.2,0) {Kernel process table\\state, PID, PPID};
\draw[ar] (file) -- node[above,font=\scriptsize]{run} (proc);
\draw[ar] (proc) -- (os);
\end{tikzpicture}
\end{mlfigure}

\section{PID and PPID: the minimum process identity}
Linux assigns every process a \term{process identifier (PID)}. A process also has a \term{parent process identifier (PPID)}. The shell that launches a command is normally its parent.

\begin{lstlisting}[style=osbash]
ps -o pid,ppid,state,comm
\end{lstlisting}

The exact numbers are temporary. The relationship matters more than the number itself.

\section{What \texttt{fork()} changes}
\texttt{fork()} asks the kernel to create a child process. Both parent and child continue after the call, but the return value tells them which role they are in.

\begin{mlfigure}
\begin{tikzpicture}[
  processnode/.style={draw=UTPBlue,fill=UTPPaleBlue,rounded corners=2mm,
    text width=4.2cm,minimum height=1.65cm,inner sep=3mm,
    align=center,font=\small\color{UTPNavy}},
  childnode/.style={processnode,draw=UTPOrange,fill=UTPPaleOrange},
  branch/.style={-{Stealth[length=2.3mm,width=1.8mm]},
    line width=0.9pt,draw=UTPTeal,rounded corners=2mm},
  connector/.style={line width=0.9pt,draw=UTPTeal}]
\node[processnode] (before) at (0,2.8) {\textbf{Before fork}\\[3pt]
  PID 4200\\[2pt]{\scriptsize one running process}};
\node[processnode] (parent) at (-3.4,-0.8) {\textbf{Parent process}\\[3pt]
  PID 4200 (unchanged)\\[2pt]{\scriptsize\texttt{fork()} returns child PID: 4201}};
\node[childnode] (child) at (3.4,-0.8) {\textbf{Child process}\\[3pt]
  PID 4201 (new)\\[2pt]{\scriptsize\texttt{fork()} returns 0}};

\coordinate (forkpoint) at (0,1.0);
\draw[connector] (before.south) --
  node[right=3pt,font=\scriptsize\color{UTPNavy}]{\texttt{fork()}} (forkpoint);
\draw[branch] (forkpoint) -| (parent.north);
\draw[branch] (forkpoint) -| (child.north);
\end{tikzpicture}
\end{mlfigure}

\begin{keyidea}
After a successful \texttt{fork()}, there are two independently scheduled processes. Output order can change from run to run unless the program adds synchronization.
\end{keyidea}

\section{Waiting for a child}
A parent can use \texttt{waitpid()} to wait for a specific child. This is both synchronization and cleanup: the parent can collect the child's exit status instead of leaving a finished child uncollected.

\begin{lstlisting}[style=osc]
pid_t child = fork();
if (child > 0) {
    int status;
    waitpid(child, &status, 0);
}
\end{lstlisting}

\section{Directories are part of the Linux namespace}
A directory is not just a graphical folder. It is a file-system object that maps names to underlying objects. C programs can inspect directory entries with \texttt{opendir()}, \texttt{readdir()}, and \texttt{closedir()}.

\begin{mlfigure}
\begin{tikzpicture}[
  d/.style={draw=UTPBlue,fill=UTPPaleBlue,rounded corners=2mm,minimum width=3.2cm,minimum height=0.8cm,font=\small\bfseries\color{UTPNavy}},
  f/.style={draw=UTPTeal,fill=UTPPaleTeal,rounded corners=1mm,minimum width=2.6cm,minimum height=0.65cm,font=\scriptsize\color{UTPNavy}},
  ar/.style={-{Stealth[length=1.8mm]},thick,draw=UTPMuted}]
\node[d] (dir) at (0,1.5) {Lab01 directory};
\node[f] (a) at (-3.2,-0.2) {\texttt{fork\_demo.c}};
\node[f] (b) at (0,-0.2) {\texttt{directory\_list.c}};
\node[f] (c) at (3.2,-0.2) {\texttt{observation.txt}};
\draw[ar] (dir) -- (a); \draw[ar] (dir) -- (b); \draw[ar] (dir) -- (c);
\end{tikzpicture}
\end{mlfigure}

Creating, renaming, or deleting a file changes these names. That is exactly what Lab 1 will observe.

\section{Chapter-to-lab bridge}
\begin{mltable}
\renewcommand{\arraystretch}{1.18}
\begin{tabularx}{\linewidth}{@{}
  >{\raggedright\arraybackslash}p{0.24\linewidth}
  >{\raggedright\arraybackslash}X
  >{\raggedright\arraybackslash}p{0.25\linewidth}@{}}
\toprule
\rowcolor{UTPPaleBlue!92}
\textbf{Concept} & \textbf{Commands} & \textbf{Evidence in Lab 1} \\
\midrule
Process identity & \texttt{getpid()}, \texttt{getppid()}, \texttt{ps} & PID/PPID table \\
Process creation & \texttt{fork()}, \texttt{waitpid()} & parent/child output \\
Directory namespace & \texttt{opendir()}, \texttt{readdir()} & before/after names \\
\bottomrule
\end{tabularx}
\end{mltable}

\labsection{1}{Processes and the Linux namespace}{sec:lab01}
\begin{labproject}{From one executable to observable processes and directory names}

\textbf{Coverage.} Chapter~\ref{ch:processes}: PID/PPID, \texttt{fork()}, directory names.\\
\textbf{Goal.} Run the exact supplied source and inspect live OS state.\\
\textbf{Explain.} Explain how one executable becomes two processes and how directory names change.

\courselabflow{Build programs}{Inspect directory}{Fork process}{Wait and verify}



\section*{Source code: read before you run}
\labsourcecard{\texttt{directory\_list.c}}{Open a directory stream and print every entry returned by Linux.}{The acquire/use/release sequence: \texttt{opendir} $\rightarrow$ \texttt{readdir} loop $\rightarrow$ \texttt{closedir}.}
\noindent\textbf{Start the file.} \texttt{directory\_list.c}, lines 1--33.
\courselabtrace{One fork, two return paths}{Parent}{fork returns the positive child PID}{Child}{fork returns zero; getppid identifies its parent}{Join}{waitpid collects the child exit status}

\lstinputlisting[style=utpcnumbered,firstline=1,lastline=33]{all_experiments/Lab01/directory_list.c}

\noindent\textbf{Continue the same file.} \texttt{directory\_list.c}, lines 34--35.
\lstinputlisting[style=utpcnumbered,firstline=34,lastline=35]{all_experiments/Lab01/directory_list.c}

\labsourcecard{\texttt{fork\_demo.c}}{Create a child process, print process identity, optionally keep the child alive, then wait for it.}{The return value from \texttt{fork()}, the child's PPID, and the parent's \texttt{waitpid}.}
\noindent\textbf{Start the file.} \texttt{fork\_demo.c}, lines 1--34.


\textbf{Use the chapter while coding.} Revisit \hyperref[ch:processes]{the earlier chapter explanation}; match its operation to the variables and state changes below.

\textbf{Follow the supplied implementation.} Read the branch condition before attributing a line of output to parent or child. Scheduling can change print order. Use the optional delay for observation, then verify the parent waits rather than inferring behaviour from a single output ordering.
\courselabtrace{Trace fork\_demo and process ownership}{1. Input / state}{fork returns a child PID in the parent and zero in the child; failure is a separate branch.}{2. Operation}{The child prints its identity; the parent retains child\_pid and waits for that specific child.}{3. Check}{waitpid collects the child's termination status; printed PID values vary by run.}


\lstinputlisting[style=utpcnumbered,firstline=1,lastline=34]{all_experiments/Lab01/fork_demo.c}

\noindent\textbf{Continue the same file.} \texttt{fork\_demo.c}, lines 35--51.
\lstinputlisting[style=utpcnumbered,firstline=35,lastline=51]{all_experiments/Lab01/fork_demo.c}

\begin{utptable}
\renewcommand{\arraystretch}{1.18}
\begin{tabularx}{\linewidth}{@{}>{\bfseries\leavevmode\color{UTPNavy}}m{0.34\linewidth}
  >{\raggedright\arraybackslash}X@{}}
\toprule
\rowcolor{UTPPaleBlue!92}
\utpth{Code region} & \utpth{What it means} \\
\midrule
\texttt{opendir} / \texttt{readdir} / \texttt{closedir} & Linux exposes directory entries through a stream-like API. \\
\texttt{child\_pid = fork()} & One running process becomes parent + child. \\
\texttt{child\_pid == 0} & This branch is executing in the child. \\
\texttt{waitpid(...)} & The parent waits and collects the child's exit status. \\
\bottomrule
\end{tabularx}
\end{utptable}

\section*{Task 1 --- Build a clean baseline}
\begin{lstlisting}[style=utpbash]
cd all_experiments/Lab01
gcc -Wall -Wextra -std=c11 directory_list.c -o directory_list
gcc -Wall -Wextra -std=c11 fork_demo.c -o fork_demo
./directory_list .
./fork_demo
\end{lstlisting}
Do not change the source yet. Record the parent PID, child PID, and child PPID from one run.

\section*{Task 2 --- Observe the namespace changing}
Run \texttt{./directory\_list .} after each command:
\begin{lstlisting}[style=utpbash]
touch observation.txt
mv observation.txt renamed.txt
rm renamed.txt
\end{lstlisting}
Create a four-row table: initial state, after \texttt{touch}, after \texttt{mv}, after \texttt{rm}. Record only the name that changed.

\section*{Task 3 --- Observe the parent and child while alive}
\begin{lstlisting}[style=utpbash]
./fork_demo 10 &
ps -C fork_demo -o pid,ppid,state,comm
wait
\end{lstlisting}
Match the child's PPID to the parent's PID. The state letter can vary because scheduling is real and timing-dependent.

\section*{Task 4 --- Controlled edit in Vim}
Open \texttt{fork\_demo.c} with Vim. Change the child message so it includes the words \texttt{child branch observed}. Save, recompile, and run again.

Before running, answer: \emph{Will the parent print the new child message too?} Then explain the observed result using the idea of two processes, not only the \texttt{if} statement.

\begin{quickcheck}
\begin{enumerate}
  \item Why does \texttt{fork()} produce two different return-value cases from one call?
  \item How can a child process identify its parent?
  \item Why can creating or renaming a file change a directory listing without changing the directory program itself?
\end{enumerate}
\end{quickcheck}

\textbf{Chapter connection.} Revisit \hyperref[ch:processes]{the chapter explanation} and map its states to the steps below.
% Reused from chapters/02-linux-processes-files-and-system-calls.tex
\begin{labvisual}{Recall the chapter model}
\begin{labgraphic}
\begin{tikzpicture}[
  box/.style={draw=UTPBlue,fill=UTPPaleBlue,rounded corners=2mm,minimum width=3.5cm,minimum height=1cm,align=center,font=\small\color{UTPNavy}},
  live/.style={box,draw=UTPTeal,fill=UTPPaleTeal},
  kern/.style={box,draw=UTPNavy,fill=UTPNavy!10},
  ar/.style={-{Stealth[length=2.2mm]},thick,draw=UTPOrange}]
\node[box] (file) at (-4.2,0) {Executable file\\\texttt{fork\_demo}};
\node[live] (proc) at (0,0) {Running process\\PID 4200};
\node[kern] (os) at (4.2,0) {Kernel process table\\state, PID, PPID};
\draw[ar] (file) -- node[above,font=\scriptsize]{run} (proc);
\draw[ar] (proc) -- (os);
\end{tikzpicture}
\end{labgraphic}\end{labvisual}


\begin{verification}
\textbf{Explain and extend your result.} Start a harmless long-running command, record its PID and parent, terminate only that process and confirm its exit. Explain why a second process with the same command name must remain unaffected.


\begin{itemize}
  \item Both programs compile with \texttt{-Wall -Wextra} and run from the unchanged source first.
  \item The child's PPID matches the parent's PID in the \texttt{ps} output.
  \item The namespace table records exactly one changed name after each of \texttt{touch}, \texttt{mv} and \texttt{rm}.
  \item Your prediction for the Vim edit is written down before you run it, and the result is explained with two processes.
\end{itemize}
\end{verification}

\begin{troubleshoot}
If \texttt{ps -C} shows nothing, rerun \texttt{./fork\_demo 10 \&} and issue \texttt{ps} before the child exits. Different PID numbers are expected on every run.
\end{troubleshoot}
\end{labproject}


\chapter{Files, Metadata, and Reliable Command-Line Utilities}\label{ch:files-cli}
\chapterguide{Chapter 1: processes, system calls, and the Linux command line.}{separate file contents from metadata, interpret file descriptors and exit status, and explain the read/write pattern used by the utilities in Lab 2.}

\section{Two views of a file}
A regular file has \textbf{contents} and \textbf{metadata}. The contents are bytes. Metadata describes the object: size, permissions, timestamps, inode number, and link count.

\begin{utpfig}[Contents and metadata]{Two views of a file: its contents are the bytes read with \texttt{read()} or \texttt{fread()}; its metadata, such as size, mode, inode, and link count, is read with \texttt{stat()}.}{fig:contents-metadata}
\begin{tikzpicture}[
  card/.style={draw=UTPBlue,fill=UTPPaleBlue,rounded corners=2mm,text width=5.0cm,minimum height=2.0cm,align=center,font=\small\color{UTPNavy}},
  meta/.style={card,draw=UTPOrange,fill=UTPPaleOrange}]
\node[card] at (-3.0,0) {\textbf{CONTENTS}\\bytes read with \texttt{read()} or \texttt{fread()}\\e.g. text, source code, records};
\node[meta] at (3.0,0) {\textbf{METADATA}\\size, mode, inode, links\\read with \texttt{stat()}};
\end{tikzpicture}
\end{utpfig}

\section{The file-I/O ladder}
The C standard library adds buffering and formatting; POSIX uses integer file descriptors closer to the OS interface.

\begin{utptable}
\renewcommand{\arraystretch}{1.18}
\begin{tabularx}{\linewidth}{@{}>{\bfseries\leavevmode\color{UTPNavy}}m{0.22\linewidth}
  >{\raggedright\arraybackslash}m{0.34\linewidth}
  >{\raggedright\arraybackslash}X@{}}
\toprule
\rowcolor{UTPPaleBlue!92}
\utpth{Layer} & \utpth{Typical functions} & \utpth{Use in this chapter} \\
\midrule
C streams & \texttt{fopen}, \texttt{fgets}, \texttt{fclose} & line-oriented text search \\
POSIX file descriptors & \texttt{open}, \texttt{read}, \texttt{write}, \texttt{close} & explicit file copy loop \\
Metadata & \texttt{stat} & inspect size, inode, links, permissions \\
\bottomrule
\end{tabularx}
\end{utptable}

\section{The acquire-use-release pattern}
Reliable systems code repeatedly follows the same shape:

\begin{utpfig}[Acquire, use, release]{The acquire-use-release pattern. A resource is opened, used, and then closed, and each step has its own failure to check.}{fig:acquire-use-release}
\begin{tikzpicture}[
  s/.style={draw=UTPBlue,fill=UTPPaleBlue,rounded corners=2mm,minimum width=3.0cm,minimum height=0.85cm,align=center,font=\small\color{UTPNavy}},
  ar/.style={-{Stealth[length=2mm]},thick,draw=UTPTeal}]
\node[s] (a) at (-4,0) {Acquire\\\texttt{open}/\texttt{fopen}};
\node[s] (b) at (0,0) {Use\\\texttt{read}/\texttt{fgets}};
\node[s] (c) at (4,0) {Release\\\texttt{close}/\texttt{fclose}};
\draw[ar] (a) -- (b); \draw[ar] (b) -- (c);
\end{tikzpicture}
\end{utpfig}

Every acquire step can fail, every loop needs a clear stopping rule, and every successfully acquired resource should eventually be released.

\section{A correct read/write loop}
\texttt{read()} can return fewer bytes than requested. Similarly, \texttt{write()} may write fewer bytes than supplied. A robust copy therefore checks return values rather than assuming one call completes the whole job.

\begin{keyidea}
For systems programming, the return value is part of the data. Do not discard it: it tells you whether an operation succeeded, reached end-of-file, or failed.
\end{keyidea}

\section{Exit status makes a CLI tool composable}
A command-line utility should communicate to both humans and other programs. The normal convention is:
\begin{itemize}
  \item \texttt{0}: success or requested condition found;
  \item non-zero: a different outcome or an error.
\end{itemize}
The supplied \texttt{mini\_grep.c} uses 0 for a match, 1 for no match, and 2 for incorrect use or a file error. This is why a shell script can make decisions without parsing sentences.

\section{Chapter-to-lab bridge}
\begin{utptable}
\renewcommand{\arraystretch}{1.18}
\begin{tabularx}{\linewidth}{@{}>{\bfseries\leavevmode\color{UTPNavy}}m{0.24\linewidth}
  >{\raggedright\arraybackslash}m{0.22\linewidth}
  >{\raggedright\arraybackslash}X@{}}
\toprule
\rowcolor{UTPPaleBlue!92}
\utpth{Concept} & \utpth{Source file} & \utpth{What you will explain} \\
\midrule
Metadata & \texttt{file\_info.c} & content is different from inode/permissions/links \\
Low-level file I/O & \texttt{file\_copy.c} & why loops check \texttt{read()} and \texttt{write()} \\
CLI contract & \texttt{mini\_grep.c} & output + exit status make a reliable tool \\
\bottomrule
\end{tabularx}
\end{utptable}

\chapter{Shell Scripting and Command-Line Automation}\label{ch:shell}
\chapterguide{Chapter 2: files, command-line utilities, and exit status.}{use variables, quoting, redirection, pipelines, conditions, and loops to automate repeatable OS observations in Lab 3.}

\section{Why the shell belongs in an operating-systems course}
A shell such as Bash is an ordinary user-space program, but it is unusually useful for observing an operating system. It starts other programs, connects their input and output, reports whether they succeeded, and gives us a small programming language for repeating tasks.

This makes the shell a convenient bridge between theory and experiment. Instead of writing C code every time we want to inspect a process or count log messages, we can combine existing tools such as \texttt{ps}, \texttt{grep}, \texttt{wc}, \texttt{df}, and \texttt{uname}.

\begin{keyidea}
A good shell script is usually not doing everything itself. It coordinates small programs and lets the operating system connect them through processes, files, and standard streams.
\end{keyidea}

\section{What happens when you run a command}
When you enter an external command such as \texttt{ps}, the shell performs a simple sequence conceptually:
\begin{enumerate}
  \item interpret the command line and expand variables;
  \item prepare redirection or pipes if requested;
  \item start the requested program;
  \item wait for it if it is a foreground command;
  \item collect its exit status.
\end{enumerate}
You do not need to implement these steps in Bash, but recognising them helps connect shell work to processes, file descriptors, \texttt{fork()}, and \texttt{exec()}.

\section{Variables, arguments, and quoting}
A shell variable stores text. Script arguments are available as \texttt{\$1}, \texttt{\$2}, and so on. Quoting matters because the shell processes text before the final command runs.

\begin{lstlisting}[style=utpbash]
file=${1:-report.txt}
echo "Writing report to: $file"
df -h . > "$file"
\end{lstlisting}

The double quotes around \texttt{"\$file"} preserve a filename containing spaces as one argument.

\begin{pitfall}
Unquoted variables are a common source of scripting errors. If \texttt{file="OS Lab.txt"}, then \texttt{cat \$file} can be interpreted as two filenames. Prefer \texttt{cat "\$file"}.
\end{pitfall}

\section{Standard input, output, and error}
Most command-line programs begin with three standard streams:
\begin{itemize}
  \item \textbf{standard input} --- where normal input is read;
  \item \textbf{standard output} --- where normal results are written;
  \item \textbf{standard error} --- where diagnostics are written.
\end{itemize}
The shell can redirect these streams without changing the program.

\begin{utptable}
\renewcommand{\arraystretch}{1.18}
\begin{tabularx}{\linewidth}{@{}>{\bfseries\leavevmode\color{UTPNavy}}m{0.10\linewidth}
  >{\raggedright\arraybackslash}X
  >{\raggedright\arraybackslash}m{0.38\linewidth}@{}}
\toprule
\rowcolor{UTPPaleBlue!92}
\utpth{Syntax} & \utpth{Meaning} & \utpth{Example} \\
\midrule
\texttt{>} & replace a file with standard output & \texttt{ps > processes.txt} \\
\texttt{>>} & append standard output & \texttt{date >> history.txt} \\
\texttt{<} & take standard input from a file & \texttt{wc -l < events.log} \\
\texttt{2>} & redirect standard error & \texttt{cmd 2> errors.txt} \\
\texttt{|} & connect one program's output to another's input & \texttt{ps | grep bash} \\
\bottomrule
\end{tabularx}
\end{utptable}

\section{Pipelines are process composition}
Consider:
\begin{lstlisting}[style=utpbash]
ps -eo pid,ppid,stat,comm | grep bash
\end{lstlisting}
The first process produces text. The shell connects that output to a pipe. The second process reads from the pipe and keeps selected lines. The programs do not need to know each other's implementation.

\begin{intuition}
A pipeline is like an assembly line for bytes. Each program performs one small job, while the shell and operating system arrange the connections.
\end{intuition}

\section{Exit status is part of the interface}
Every command returns an integer status. By convention, 0 means success and a non-zero value means failure or a false result. The shell stores the most recent status in \texttt{\$?}.

\begin{lstlisting}[style=utpbash]
grep kernel notes.txt
echo $?
\end{lstlisting}
This is why commands can be composed inside \texttt{if}, \texttt{\&\&}, and \texttt{||} expressions.

\begin{workedexample}
A search tool can use status 0 for ``match found'', 1 for ``no match'', and 2 for ``could not read the file''. A script can then react to the condition without trying to interpret printed sentences.
\end{workedexample}

\section{Conditions and loops for repeatable observations}
Shell conditions and loops are most useful here when they automate repeated OS observations. For example, a script can count several log levels using one loop:
\begin{lstlisting}[style=utpbash]
for level in INFO WARN ERROR; do
    count=$(grep -c "^$level " events.log || true)
    echo "$level: $count"
done
\end{lstlisting}
The mathematics is simple; the important idea is repeatability. If the log changes, rerunning the same script applies the same method again.

\section{Background work and process observation}
Adding \texttt{\&} runs a command in the background:
\begin{lstlisting}[style=utpbash]
sleep 30 &
ps -o pid,ppid,state,comm
wait
\end{lstlisting}
This provides an easy way to create a temporary process and then inspect it. The \texttt{wait} command asks the shell to wait for its outstanding background job.

\section{A professional scripting workflow}
For a student lab script, a reliable workflow is:
\begin{enumerate}
  \item define the intended input and output;
  \item make one normal case work;
  \item quote variables and check required arguments;
  \item add one failure or boundary test;
  \item return a meaningful exit status;
  \item add a short usage message if another person will run the script.
\end{enumerate}

\begin{keyidea}
The goal of Bash in this course is not to memorise syntax. It is to make OS observations reproducible: the same commands, in the same order, with failures handled clearly.
\end{keyidea}

\begin{quickcheck}
\begin{enumerate}
  \item Why should a filename variable usually be quoted?
  \item What OS idea is visible in a pipeline such as \texttt{ps | grep bash}?
  \item Why is exit status more useful to another program than a human-readable sentence alone?
  \item What does appending \texttt{\&} change about how the shell treats a command?
\end{enumerate}
\end{quickcheck}


\part{CPU, Concurrency, and Resource Management}
\chapter{CPU Scheduling}\label{ch:scheduling}
\chapterguide{Chapter 1: processes, plus comfort reading simple tables and timelines.}{calculate waiting, turnaround, response, and completion time; trace FCFS, SJF, Priority, and Round Robin; and compare them fairly in Lab 4.}

\section{The decision the scheduler makes}
Three processes are ready. One CPU is free. Which one runs?

The OS answers this repeatedly. The \term{scheduler} chooses a process; the
\term{dispatcher} gives it the CPU.

There is no single best choice. Short jobs reduce average waiting, small turns
improve responsiveness, and priorities favour important work.

\section{Five quantities, and what each one means to a user}
Use $A$ for arrival, $B$ for CPU burst, $S$ for first start, and $C$ for
completion. The three rules are: $T=C-A$, $W=T-B$, and $R=S-A$.

Each measures a different delay.

\begin{mltable}
\renewcommand{\arraystretch}{1.18}
\begin{tabularx}{\linewidth}{@{}
  >{\raggedright\arraybackslash}p{0.19\linewidth}
  >{\raggedright\arraybackslash}X
  >{\raggedright\arraybackslash}X@{}}
\toprule
\rowcolor{UTPPaleBlue!92}
\textbf{Quantity} & \textbf{What it measures} & \textbf{What the user notices} \\
\midrule
Turnaround $T$ & Total time in the system, start to finish & ``My batch job took 9 seconds'' \\
Waiting $W$ & Time sat in the ready queue doing nothing & Pure overhead --- the part a better policy can reduce \\
Response $R$ & Delay before \emph{any} progress is visible & ``The window froze when I clicked'' \\
\bottomrule
\end{tabularx}
\end{mltable}

\begin{keyidea}
Waiting and turnaround are closely related because the burst time is fixed.
Response time is different: it measures how quickly a process starts, not how
quickly it finishes.
\end{keyidea}

\begin{workedexample}
A process arrives at time 2, first receives the CPU at time 5, needs 4 CPU units in total, and finishes at time 11.

$R=3$, $T=9$, and $W=5$. A negative waiting time always signals an error.
\end{workedexample}

\section{One workload, four policies}
Everything below uses the same three processes, all arriving at time 0. Using one workload throughout is what makes the policies comparable.

\begin{mltable}
\renewcommand{\arraystretch}{1.18}
\begin{tabularx}{\linewidth}{@{}
  >{\centering\arraybackslash}p{0.16\linewidth}
  >{\centering\arraybackslash}p{0.16\linewidth}
  >{\centering\arraybackslash}p{0.16\linewidth}
  >{\centering\arraybackslash}X@{}}
\toprule
\rowcolor{UTPPaleBlue!92}
\textbf{Process} & \textbf{Arrival $A$} & \textbf{Burst $B$} & \textbf{Priority (lower = more urgent)} \\
\midrule
P1 & 0 & 5 & 3 \\
P2 & 0 & 3 & 1 \\
P3 & 0 & 1 & 2 \\
\bottomrule
\end{tabularx}
\end{mltable}

\subsection{First-Come, First-Served}
\term{FCFS} is non-preemptive and runs processes in arrival order. Here that is P1, P2, P3.

\begin{mlfigure}
\begin{tikzpicture}[x=0.95cm,y=0.7cm]
  \fill[MLPaleBlue,draw=MLBlue] (0,0) rectangle (5,1); \node at (2.5,0.5) {\small\bfseries P1};
  \fill[MLPaleTeal,draw=MLBlue] (5,0) rectangle (8,1); \node at (6.5,0.5) {\small\bfseries P2};
  \fill[MLPaleOrange,draw=MLBlue] (8,0) rectangle (9,1); \node at (8.5,0.5) {\small\bfseries P3};
  \foreach \x in {0,5,8,9} { \draw[MLMuted] (\x,0) -- (\x,-0.25); \node[below,font=\scriptsize] at (\x,-0.25) {\x}; }
\end{tikzpicture}
\end{mlfigure}

\begin{mltable}
\renewcommand{\arraystretch}{1.18}
\begin{tabularx}{\linewidth}{@{}
  >{\centering\arraybackslash}p{0.13\linewidth}
  *{5}{>{\centering\arraybackslash}X}@{}}
\toprule
\rowcolor{UTPPaleBlue!92}
\textbf{Process} & $\boldsymbol{B}$ & $\boldsymbol{C}$ & $\boldsymbol{T=C-A}$ & $\boldsymbol{W=T-B}$ & $\boldsymbol{R=S-A}$ \\
\midrule
P1 & 5 & 5 & 5 & 0 & 0 \\
P2 & 3 & 8 & 8 & 5 & 5 \\
P3 & 1 & 9 & 9 & 8 & 8 \\
\midrule
\multicolumn{3}{@{}l}{\bfseries Average} & \bfseries 7.33 & \bfseries 4.33 & \bfseries 4.33 \\
\bottomrule
\end{tabularx}
\end{mltable}

P3 needs one unit of CPU and waits eight to get it. That is the \term{convoy effect}: short jobs stuck behind a long one, like cars behind a lorry on a single-lane road.

\subsection{Shortest Job First}
\term{SJF} picks the smallest remaining burst: P3, then P2, then P1.

\begin{mlfigure}
\begin{tikzpicture}[x=0.95cm,y=0.7cm]
  \fill[MLPaleOrange,draw=MLBlue] (0,0) rectangle (1,1); \node at (0.5,0.5) {\small\bfseries P3};
  \fill[MLPaleTeal,draw=MLBlue] (1,0) rectangle (4,1); \node at (2.5,0.5) {\small\bfseries P2};
  \fill[MLPaleBlue,draw=MLBlue] (4,0) rectangle (9,1); \node at (6.5,0.5) {\small\bfseries P1};
  \foreach \x in {0,1,4,9} { \draw[MLMuted] (\x,0) -- (\x,-0.25); \node[below,font=\scriptsize] at (\x,-0.25) {\x}; }
\end{tikzpicture}
\end{mlfigure}

\begin{mltable}
\renewcommand{\arraystretch}{1.18}
\begin{tabularx}{\linewidth}{@{}
  >{\raggedright\arraybackslash}p{0.13\linewidth}
  *{5}{>{\raggedright\arraybackslash}X}@{}}
\toprule
\rowcolor{UTPPaleBlue!92}
\textbf{Process} & $\boldsymbol{B}$ & $\boldsymbol{C}$ & $\boldsymbol{T}$ & $\boldsymbol{W}$ & $\boldsymbol{R}$ \\
\midrule
P3 & 1 & 1 & 1 & 0 & 0 \\
P2 & 3 & 4 & 4 & 1 & 1 \\
P1 & 5 & 9 & 9 & 4 & 4 \\
\midrule
\multicolumn{3}{@{}l}{\bfseries Average} & \bfseries 4.67 & \bfseries 1.67 & \bfseries 1.67 \\
\bottomrule
\end{tabularx}
\end{mltable}

Average waiting fell from 4.33 to 1.67 --- and no work was done faster. The CPU was busy for exactly nine units in both schedules. Only the \emph{order} changed.

\begin{keyidea}
For a fixed set of jobs all available at once, SJF provably minimises average waiting time. Running a short job early delays only that short job's successors briefly; running a long job early delays everyone behind it.

The catch is that $B_i$ is the burst the process is \emph{about to} need, which the OS cannot know. Real systems estimate it from recent history, and the estimate is often wrong.
\end{keyidea}

\subsection{Priority scheduling}
Highest priority first. With the convention that a lower number is more urgent, the order is P2, P3, P1.

\begin{mlfigure}
\begin{tikzpicture}[x=0.95cm,y=0.7cm]
  \fill[MLPaleTeal,draw=MLBlue] (0,0) rectangle (3,1); \node at (1.5,0.5) {\small\bfseries P2};
  \fill[MLPaleOrange,draw=MLBlue] (3,0) rectangle (4,1); \node at (3.5,0.5) {\small\bfseries P3};
  \fill[MLPaleBlue,draw=MLBlue] (4,0) rectangle (9,1); \node at (6.5,0.5) {\small\bfseries P1};
  \foreach \x in {0,3,4,9} { \draw[MLMuted] (\x,0) -- (\x,-0.25); \node[below,font=\scriptsize] at (\x,-0.25) {\x}; }
\end{tikzpicture}
\end{mlfigure}

Averages: $T = 5.33$, $W = 2.33$. Better than FCFS, worse than SJF --- because priority happens to correlate with burst length here, but is not chosen to optimise it.

\begin{pitfall}
State your priority convention before you compute anything. ``Priority 1'' means \emph{most} urgent in most textbooks and lab manuals, and \emph{least} urgent in some real systems --- Linux \texttt{nice} values run from $-20$ (most favourable) to $19$.

Half of all lost marks on scheduling questions come from silently assuming the wrong direction. Write the convention at the top of your answer.
\end{pitfall}

Priority scheduling can starve a low-priority process indefinitely if urgent work keeps arriving. The standard remedy is \term{aging}: gradually raise the priority of anything that has waited a long time, so every process eventually becomes urgent enough to run.

\subsection{Round Robin}
\term{Round Robin} is preemptive. Each process gets at most a \term{time quantum} $q$, then goes to the back of the queue if unfinished. With $q = 2$ and queue order P1, P2, P3:

\begin{mlfigure}
\begin{tikzpicture}[x=0.95cm,y=0.7cm]
  \fill[MLPaleBlue,draw=MLBlue]   (0,0) rectangle (2,1); \node at (1,0.5) {\small\bfseries P1};
  \fill[MLPaleTeal,draw=MLBlue]   (2,0) rectangle (4,1); \node at (3,0.5) {\small\bfseries P2};
  \fill[MLPaleOrange,draw=MLBlue] (4,0) rectangle (5,1); \node at (4.5,0.5) {\small\bfseries P3};
  \fill[MLPaleBlue,draw=MLBlue]   (5,0) rectangle (7,1); \node at (6,0.5) {\small\bfseries P1};
  \fill[MLPaleTeal,draw=MLBlue]   (7,0) rectangle (8,1); \node at (7.5,0.5) {\small\bfseries P2};
  \fill[MLPaleBlue,draw=MLBlue]   (8,0) rectangle (9,1); \node at (8.5,0.5) {\small\bfseries P1};
  \foreach \x in {0,2,4,5,7,8,9} { \draw[MLMuted] (\x,0) -- (\x,-0.25); \node[below,font=\scriptsize] at (\x,-0.25) {\x}; }
\end{tikzpicture}
\end{mlfigure}

Note that P3 finishes early at time 5 despite only getting the CPU third: it needed one unit, which is less than the quantum, so it ran to completion in its first turn.

\begin{mltable}
\renewcommand{\arraystretch}{1.18}
\begin{tabularx}{\linewidth}{@{}
  >{\centering\arraybackslash}p{0.13\linewidth}
  *{5}{>{\centering\arraybackslash}X}@{}}
\toprule
\rowcolor{UTPPaleBlue!92}
\textbf{Process} & $\boldsymbol{B}$ & $\boldsymbol{C}$ & $\boldsymbol{T}$ & $\boldsymbol{W}$ & $\boldsymbol{R}$ \\
\midrule
P1 & 5 & 9 & 9 & 4 & 0 \\
P2 & 3 & 8 & 8 & 5 & 2 \\
P3 & 1 & 5 & 5 & 4 & 4 \\
\midrule
\multicolumn{3}{@{}l}{\bfseries Average} & \bfseries 7.33 & \bfseries 4.33 & \bfseries 2.00 \\
\bottomrule
\end{tabularx}
\end{mltable}

\section{Reading the comparison honestly}
\begin{mltable}
\renewcommand{\arraystretch}{1.18}
\begin{tabularx}{\linewidth}{@{}
  >{\raggedright\arraybackslash}p{0.26\linewidth}
  *{3}{>{\raggedright\arraybackslash}X}@{}}
\toprule
\rowcolor{UTPPaleBlue!92}
\textbf{Policy} & \textbf{Avg turnaround} & \textbf{Avg waiting} & \textbf{Avg response} \\
\midrule
FCFS & 7.33 & 4.33 & 4.33 \\
SJF & \bfseries 4.67 & \bfseries 1.67 & 1.67 \\
Priority & 5.33 & 2.33 & 2.33 \\
Round Robin ($q=2$) & 7.33 & 4.33 & \bfseries 2.00 \\
\bottomrule
\end{tabularx}
\end{mltable}

Round Robin improves first response without making the CPU faster; it changes
how waiting is shared. SJF performs best here, but it needs future burst lengths
and can delay long jobs indefinitely.

\begin{intuition}
The quantum controls the trade-off. A large quantum makes Round Robin behave
like FCFS. A very small quantum improves response but causes more context
switches.
\end{intuition}

\section{Drawing a Gantt chart that survives checking}
Work in this order, every time:

\begin{mlsteps}
  \item Draw the timeline first, marking every instant the CPU changes hands. Do not compute anything yet.
  \item Read the completion time $C_i$ off the chart --- the right-hand edge of that process's \emph{last} block.
  \item Compute $T_i = C_i - A_i$.
  \item Compute $W_i = T_i - B_i$. Never count queue intervals by hand; the subtraction is what stops you double-counting.
  \item Read $S_i$ off the chart for $R_i$ --- the left edge of that process's \emph{first} block.
\end{mlsteps}

\begin{pitfall}
The most common Round Robin error is computing waiting time by adding up gaps in the chart. With several preemptions there are several gaps per process, and one is always missed.

Use $W = T - B$. It cannot miss a gap, because it never looks at them.
\end{pitfall}

\begin{tryit}
Redo the same workload with $q = 1$, then with $q = 5$. Confirm that $q = 5$ reproduces the FCFS chart exactly, and explain in one sentence why.

Then change the arrivals to $A_1 = 0$, $A_2 = 1$, $A_3 = 2$ and redo FCFS. The chart now contains idle time if a process has not arrived yet --- make sure your $C$ values account for it.
\end{tryit}

\begin{quickcheck}
\begin{enumerate}
  \item Why can Round Robin have the same average waiting time as FCFS but much better response time?
  \item A process has $T = 12$ and $B = 15$. What went wrong?
  \item Why does SJF need an estimate rather than a measurement?
  \item What does aging fix, and which policy needs it?
\end{enumerate}

\medskip\noindent\textbf{Answers.}
\begin{enumerate}
  \item Waiting time depends only on total time queued, which the same total work distributes similarly. Response time depends on the \emph{first} turn, and RR guarantees every process one within $n \times q$ of arriving.
  \item An arithmetic error. Waiting time would be $-3$, and a process cannot exist for less time than it spent running. Recheck $C$, or check that $B$ is total burst rather than remaining burst.
  \item $B_i$ is the burst the process is about to request. It is only measurable after the fact, so schedulers predict it from recent history --- typically an exponential average of previous bursts.
  \item Starvation. Priority scheduling needs it: without aging, a steady stream of urgent work means a low-priority process may never run at all.
\end{enumerate}
\end{quickcheck}

\labsection{4}{CPU scheduling: compare policies on one workload}{sec:lab04}
\begin{labproject}{Make scheduling policy visible through order, response, and waiting time}

\oslabflow{Chapter 4: scheduling metrics and policy rules.}{Run four policies on the same bursts so only the policy changes.}{Use the trace to explain why order or time quantum changes a metric.}

\section*{Source code}
\ossourcecard{\texttt{fcfs.c}}{Run processes in the supplied order.}{Waiting time is the accumulated burst time of earlier processes.}
\lstinputlisting[style=oslabc]{all_experiments/Lab04/fcfs.c}

\ossourcecard{\texttt{sjf.c}}{Sort the same workload by shortest burst first.}{The sorting step changes execution order; the metric loop is otherwise very similar to FCFS.}
\lstinputlisting[style=oslabc]{all_experiments/Lab04/sjf.c}

\ossourcecard{\texttt{priority.c}}{Sort by priority, where a smaller number means higher priority.}{The code uses PID only as a tie-breaker so equal priorities are deterministic.}
\lstinputlisting[style=oslabc]{all_experiments/Lab04/priority.c}

\ossourcecard{\texttt{round\_robin.c}}{Give each unfinished process a bounded time slice.}{\texttt{remaining[]} tracks work left; \texttt{first\_run[]} gives response time when all arrivals are 0.}
\lstinputlisting[style=oslabc]{all_experiments/Lab04/round_robin.c}

\section*{Task 1 --- Compile all policies}
\begin{lstlisting}[style=osbash]
cd all_experiments/Lab04
for f in fcfs sjf priority round_robin; do
  gcc -Wall -Wextra -std=c11 "$f.c" -o "$f"
done
\end{lstlisting}

\section*{Task 2 --- Baseline comparison}
\begin{lstlisting}[style=osbash]
./fcfs
./sjf
./priority
./round_robin 2
\end{lstlisting}
Complete a table with execution order, first process to finish, average waiting time when printed, and one observation about responsiveness.

\begin{mltable}
\renewcommand{\arraystretch}{1.18}
\begin{tabularx}{\linewidth}{@{}
  >{\raggedright\arraybackslash}p{0.18\linewidth}
  *{3}{>{\raggedright\arraybackslash}X}@{}}
\toprule
\rowcolor{UTPPaleBlue!92}
\textbf{Policy} & \textbf{First execution order} & \textbf{Waiting/response evidence} & \textbf{What the rule favours} \\
\midrule
FCFS & & & \\
SJF & & & \\
Priority & & & \\
RR, $q=2$ & & & \\
\bottomrule
\end{tabularx}
\end{mltable}

\section*{Task 3 --- Verify one result by hand}
For FCFS, draw a one-line Gantt chart for bursts 8, 4, 2, 1. Compute the waiting time of each process before comparing your arithmetic with the program.

\section*{Task 4 --- Change only the Round Robin quantum}
\begin{lstlisting}[style=osbash]
./round_robin 1
./round_robin 2
./round_robin 4
\end{lstlisting}
Record the first response time of each process for each quantum. Explain the trade-off visible in the Gantt output: smaller turns provide earlier rotation but create more slices.

\begin{quickcheck}
\begin{enumerate}
  \item Which scheduling metric measures time spent waiting in the ready queue?
  \item Why can Round Robin improve response time while increasing the number of context switches?
  \item Why must all policies be compared using the same workload?
\end{enumerate}
\end{quickcheck}

\begin{labdeliverable}
Submit the four-policy comparison table, your hand-checked FCFS Gantt chart, and a small $q=1,2,4$ response-time table. Your explanation should identify a policy rule first and a metric consequence second.
\end{labdeliverable}
\end{labproject}


\chapter{Interprocess Communication and Synchronization}\label{ch:ipc}
\chapterguide{Chapter 1: processes and Chapter 3: standard input/output and pipelines.}{compare pipes and shared memory, identify a critical section and race condition, and explain why shared state requires synchronization before Lab 5.}

\section{Concurrency creates more than one valid execution order}
When two activities can make progress at the same time, their individual instructions may interleave in several valid ways. The operating system may preempt one process or thread, run another, and later return to the first.

If the activities touch only private state, this may not matter. If they both read and update the same state, the exact interleaving can change the result.

\section{Race conditions and critical sections}
A \term{race condition} occurs when correctness depends on timing that the program has not controlled. Suppose two workers both execute ``read counter, add one, write counter.'' If both read the same old value before either writes, one increment is lost.

The code that must behave as one indivisible logical action is called a \term{critical section}. A synchronization mechanism prevents unsafe overlap.

\begin{keyidea}
Synchronization is not about forcing the whole program to run one step at a time. It protects only the state transition that must not be interleaved incorrectly.
\end{keyidea}

\section{Mutexes and semaphores in plain language}
A \term{mutex} is commonly used to give one thread at a time exclusive access to a critical section. A \term{semaphore} maintains a count that can represent available units of a resource or coordinate producer-consumer conditions.

\begin{mltable}
\renewcommand{\arraystretch}{1.18}
\begin{tabularx}{\linewidth}{@{}
  >{\raggedright\arraybackslash}p{0.22\linewidth}
  >{\raggedright\arraybackslash}p{0.32\linewidth}
  >{\raggedright\arraybackslash}X@{}}
\toprule
\rowcolor{UTPPaleBlue!92}
\textbf{Mechanism} & \textbf{Think of it as} & \textbf{Typical use} \\
\midrule
Mutex & one key to a protected region & protect a shared counter or data structure \\
Counting semaphore & a controlled count of available units & free slots, queued items, limited resources \\
\bottomrule
\end{tabularx}
\end{mltable}

\section{Producer-consumer and the bounded buffer}
A bounded buffer with capacity $N$ can be reasoned about using two counts:
\[
\texttt{empty} + \texttt{full} = N.
\]
A producer requires an empty slot before adding an item. A consumer requires a full slot before removing one. In a real concurrent implementation, checking and changing these synchronization values must itself be performed safely.

\begin{workedexample}
For capacity 3, the initial state may be \texttt{empty=3, full=0}. After one successful production it becomes \texttt{empty=2, full=1}. A fourth production without any intervening consumption must wait because no empty slot remains.
\end{workedexample}

\section{Communication and synchronization are different questions}
Processes have separate address spaces, so they need an interprocess communication (IPC) mechanism to exchange data. Common models include:

\begin{mltable}
\renewcommand{\arraystretch}{1.18}
\begin{tabularx}{\linewidth}{@{}
  >{\raggedright\arraybackslash}p{0.19\linewidth}
  >{\raggedright\arraybackslash}p{0.35\linewidth}
  >{\raggedright\arraybackslash}X@{}}
\toprule
\rowcolor{UTPPaleBlue!92}
\textbf{IPC model} & \textbf{Main idea} & \textbf{Introductory trade-off} \\
\midrule
Pipe & ordered byte stream between endpoints & simple flow, usually directional \\
Message queue & discrete messages managed by the kernel & message boundaries are preserved \\
Shared memory & multiple processes map the same memory region & fast sharing, but concurrent writers need coordination \\
Socket & endpoint-based communication & works locally or across networks \\
\bottomrule
\end{tabularx}
\end{mltable}

\section{A pipe moves data; shared memory exposes common data}
With a pipe, one process writes bytes and another reads them. Data conceptually travels through a kernel-managed channel.

With shared memory, processes attach the same kernel-managed memory object into their separate address spaces. The bytes do not need to be copied through a read/write stream each time, but the processes must agree on when data is ready and who may modify it.

\begin{pitfall}
Shared memory solves the ``where can both processes see the data?'' problem. It does not by itself solve the ``who may update it now?'' problem.
\end{pitfall}

\section{System V shared-memory lifecycle}
The lab uses System V shared memory because its lifecycle is visible:
\begin{enumerate}
  \item create or locate a key;
  \item create/find a segment with \texttt{shmget};
  \item attach it with \texttt{shmat};
  \item read or write the shared bytes;
  \item detach with \texttt{shmdt};
  \item remove the segment with \texttt{shmctl(..., IPC\_RMID, ...)}.
\end{enumerate}

The \texttt{ipcs -m} command lets students observe these kernel-managed objects while they exist.

\section{How to reason about synchronization problems}
For an introductory problem, ask four questions:
\begin{enumerate}
  \item What state is shared?
  \item What invariant should always remain true?
  \item Which operation changes that state?
  \item What must wait when the required condition is not available?
\end{enumerate}
This approach is more reliable than memorising a sequence of \texttt{wait()} and \texttt{signal()} calls without understanding the state they protect.

\begin{quickcheck}
\begin{enumerate}
  \item Why can shared memory require synchronization even after both processes attach successfully?
  \item What invariant holds for \texttt{empty} and \texttt{full} in a fixed-capacity buffer?
  \item What is one conceptual difference between a pipe and shared memory?
\end{enumerate}
\end{quickcheck}

\labsection{5}{Synchronization and interprocess communication}{sec:lab05}
\begin{labproject}{Trace bounded-buffer state and compare a pipe with shared memory}
\textbf{Suggested time:} 100--120 minutes \hfill
\textbf{Level:} intermediate\par
\textbf{Starter files:} \texttt{producer\_consumer\_sim.c}, \texttt{pipe\_demo.c}, \texttt{shm\_writer.c}, \texttt{shm\_reader.c}.\par
\textbf{Evidence to keep:} invariant trace, IPC comparison table, and shared-memory lifecycle evidence.

\begin{labmode}
Two questions stay separate throughout this lab: \emph{how do processes exchange data?} and \emph{how do concurrent activities coordinate safe access to shared state?}
\end{labmode}

\medskip\noindent\textbf{Learning outcomes.}
\begin{enumerate}
  \item use a simple invariant to reason about producer-consumer state;
  \item explain the communication model of a pipe and of shared memory;
  \item identify create/attach/use/detach/remove stages for System V shared memory.
\end{enumerate}

\medskip\noindent\textbf{Part A --- Producer-consumer as a state problem.}
Compile and run \texttt{producer\_consumer\_sim.c}. For each action record \texttt{empty}, \texttt{full}, and whether the action was allowed. With capacity 3, check after every successful operation that
\[
\texttt{empty} + \texttt{full} = 3.
\]
Explain why producing into a full buffer and consuming from an empty buffer must wait or fail rather than blindly changing the counters.

\begin{pitfall}
This first program is a sequential model of semaphore state. It helps you understand the rules, but the integers are not themselves a real thread-safe semaphore implementation.
\end{pitfall}

\medskip\noindent\textbf{Part B --- A pipe moves bytes through the kernel.}
\begin{lstlisting}[style=osbash]
gcc -Wall -Wextra pipe_demo.c -o pipe_demo
./pipe_demo
\end{lstlisting}
Find the call that creates the pipe, the parent's write end, and the child's read end. Explain why the unused ends are closed in each process.

\medskip\noindent\textbf{Part C --- Shared memory exposes the same bytes to two processes.}
\begin{lstlisting}[style=osbash]
gcc -Wall -Wextra shm_writer.c -o shm_writer
gcc -Wall -Wextra shm_reader.c -o shm_reader
ipcs -m
./shm_writer
ipcs -m
./shm_reader
ipcs -m
\end{lstlisting}
Record the shared-memory segment ID and compare the three \texttt{ipcs -m} snapshots. Map each program operation to the lifecycle: create/find $\rightarrow$ attach $\rightarrow$ access $\rightarrow$ detach $\rightarrow$ remove.

\begin{workedexample}
A shared-memory segment is a kernel-managed IPC object. Each process attaches that object into its own address space. Detaching one process does not automatically mean the object has been removed from the system.
\end{workedexample}

\medskip\noindent\textbf{Part D --- Compare the mechanisms.}
Complete a small comparison with rows for data movement, lifetime, setup complexity, and need for synchronization. A pipe provides an ordered byte stream; shared memory gives processes access to a common region and therefore usually needs an additional coordination mechanism when both sides can write.

\medskip\noindent\textbf{Part E --- Failure case.}
Run the reader before the writer and record the error. Then run the normal writer/reader sequence again. A useful systems program should make failure visible instead of silently continuing with invalid state.

\begin{verification}
Submit the producer-consumer trace, one paragraph comparing pipe and shared memory, and terminal evidence showing a shared segment appear and later disappear.
\end{verification}

\begin{troubleshoot}
If a segment remains after a crash, inspect only your own IPC objects with \texttt{ipcs -m}. Remove a known stale segment with \texttt{ipcrm -m ID}. On a shared university machine, never remove another user's IPC object.
\end{troubleshoot}
\end{labproject}


\chapter{Threads, Resource Allocation, and Deadlocks}\label{ch:deadlock}
\chapterguide{Chapter 5: shared state, race conditions, and synchronization.}{distinguish processes from threads, use a mutex to protect a critical section, identify the four deadlock conditions, and reason about safe states in Lab 6.}

\section{POSIX threads}
The POSIX threads API creates multiple execution streams inside one process. \texttt{pthread\_create()} starts a thread, \texttt{pthread\_self()} identifies the calling thread, \texttt{pthread\_equal()} compares thread identifiers, and \texttt{pthread\_join()} waits for completion.

\begin{lstlisting}[style=utpc]
pthread_t t;
int err = pthread_create(&t, NULL, worker, NULL);
if (err != 0) {
    /* handle error */
}
pthread_join(t, NULL);
\end{lstlisting}

\begin{pitfall}
Thread identifiers are opaque values. Portable code compares them with \texttt{pthread\_equal()} rather than assuming they are ordinary integers with a specific representation.
\end{pitfall}

\section{Mutexes}
A \term{mutex} protects a critical section so only one thread owns the lock at a time. The normal lifecycle is initialize, lock, access shared state, unlock, then destroy after all users are finished.

\begin{utpfig}[What threads share]{What the threads of one process share. Code, globals, and the heap exist once and are seen by every thread; each thread has its own stack and registers, which need no synchronization.}{fig:thread-memory}
\begin{tikzpicture}[
  seg/.style={draw=UTPBlue,fill=UTPPaleBlue,minimum width=3.0cm,minimum height=0.8cm,
              font=\scriptsize\color{UTPNavy}},
  shared/.style={seg,fill=UTPPaleOrange,draw=UTPOrange},
  stk/.style={draw=UTPTeal,fill=UTPPaleTeal,minimum width=1.7cm,minimum height=0.8cm,
              font=\scriptsize\color{UTPNavy}},
  lab/.style={font=\scriptsize\color{UTPMuted},align=center},
  warn/.style={font=\scriptsize\color{UTPOrange},align=center}]

\draw[draw=UTPNavy,rounded corners=2mm,line width=0.9pt] (-2.1,-2.75) rectangle (8.7,1.75);
\node[font=\small\bfseries\color{UTPNavy},anchor=north west] at (-2.0,1.65) {one process};

\node[seg]    at (-0.4,0.85) {code (text)};
\node[shared] at (-0.4,-0.35) {globals};
\node[shared] at (-0.4,-1.55) {heap};
\node[lab,text width=2.6cm] at (-0.4,-2.30) {one copy, seen by every thread};

\foreach \i/\x in {1/2.9, 2/5.2, 3/7.5} {
  \node[stk] at (\x,0.85) {stack \i};
  \node[stk] at (\x,-0.35) {registers \i};
  \node[lab] at (\x,-1.08) {thread \i};
}
\node[lab,text width=5.6cm] at (5.2,-2.05) {private to each thread --- no synchronisation needed};

\node[warn,text width=11cm] at (3.3,-3.55)
  {The orange region is the whole reason a mutex exists. Two threads writing a
   global without one is a data race, and the outcome depends on scheduling ---
   which is why such a bug reproduces only sometimes.};

\end{tikzpicture}
\end{utpfig}

\section{What deadlock means}
A \term{deadlock} is a state in which a set of processes or threads are permanently blocked because each waits for an event that only another member of the set can cause.

Four Coffman conditions must hold simultaneously for resource deadlock:
\begin{enumerate}
  \item \textbf{Mutual exclusion}: at least one resource cannot be shared concurrently.
  \item \textbf{Hold and wait}: a process holds resources while requesting more.
  \item \textbf{No preemption}: resources cannot simply be taken away.
  \item \textbf{Circular wait}: a cycle of processes each waits for a resource held by the next.
\end{enumerate}
Breaking any one condition prevents this class of deadlock.

\begin{utpfig}[A circular wait]{A circular wait. P1 holds R1 and requests R2, while P2 holds R2 and requests R1, so neither can proceed.}{fig:circular-wait}
\begin{tikzpicture}[
  proc/.style={circle,draw=UTPBlue,fill=UTPPaleBlue,minimum size=1.05cm,
               font=\small\bfseries\color{UTPNavy}},
  res/.style={draw=UTPNavy,fill=white,minimum width=1.05cm,minimum height=0.9cm,
              font=\small\bfseries\color{UTPNavy}},
  cyc/.style={-{Stealth[length=2.3mm]},line width=1pt,draw=UTPRed},
  lab/.style={font=\scriptsize\color{UTPMuted}},
  note/.style={font=\scriptsize\color{UTPInk},align=center}]

\node[proc] (p1) at (0,2.5)   {P1};
\node[res]  (r2) at (5.0,2.5) {R2};
\node[proc] (p2) at (5.0,0)   {P2};
\node[res]  (r1) at (0,0)     {R1};

\draw[cyc] (p1) -- node[lab,above]{requests} (r2);
\draw[cyc] (r2) -- node[lab,right]{assigned to} (p2);
\draw[cyc] (p2) -- node[lab,below]{requests} (r1);
\draw[cyc] (r1) -- node[lab,left]{assigned to} (p1);

\node[note,text width=4.2cm] at (2.5,1.25)
  {the cycle \emph{is} the circular wait --- neither process can move, and
   nothing outside the cycle will move them};

\node[note,text width=11cm,text=UTPMuted] at (2.5,-1.15)
  {Each resource has one instance, so this cycle is a deadlock. Remove any one
   arrow --- by preemption, by releasing before requesting, or by ordering the
   resources --- and the cycle opens.};

\end{tikzpicture}
\end{utpfig}

\section{Prevention, avoidance, and detection}
\begin{utptable}
\renewcommand{\arraystretch}{1.18}
\begin{tabularx}{\linewidth}{@{}>{\bfseries\leavevmode\color{UTPNavy}}m{0.22\linewidth}
  >{\raggedright\arraybackslash}X
  >{\raggedright\arraybackslash}X@{}}
\toprule
\rowcolor{UTPPaleBlue!92}
\utpth{Strategy} & \utpth{Core idea} & \utpth{Trade-off} \\
\midrule
Prevention & Design the protocol so a Coffman condition cannot hold & Can reduce utilization or flexibility \\
Avoidance & Grant a request only if the resulting state remains safe & Requires future maximum-demand information \\
Detection/recovery & Allow deadlock, detect it later, then recover & Detection and recovery can be expensive \\
\bottomrule
\end{tabularx}
\end{utptable}

\section{Banker's algorithm}
Banker's algorithm avoids unsafe allocations. First calculate
\texttt{Need = Max - Allocation}. Start with \texttt{Work = Available}, then
find a process whose needs fit within Work. Assume it finishes and returns its
allocation to Work. If every process can finish in some order, the state is
\term{safe}.

\begin{workedexample}
Suppose Available is $(1,1)$, P1 has Need $(1,0)$ and Allocation $(1,1)$, while P2 has Need $(2,1)$ and Allocation $(1,0)$. P1 can finish first because $(1,0)\le(1,1)$. Work becomes $(2,2)$, allowing P2 to finish. Therefore a safe sequence is P1, P2.
\end{workedexample}

\begin{utpfig}[The Banker's safety check]{The Banker's safety check for the worked example. P1's need fits within Work, so it can finish and release its allocation; the larger Work then lets P2 finish, so P1, P2 is a safe sequence.}{fig:bankers}
\begin{tikzpicture}[
  wk/.style={draw=UTPTeal,fill=UTPPaleTeal,rounded corners=1.5mm,minimum width=1.5cm,
             minimum height=0.75cm,font=\small\bfseries\color{UTPNavy}},
  test/.style={draw=UTPBlue,fill=UTPPaleBlue,rounded corners=1.5mm,text width=3.0cm,
               minimum height=1.15cm,align=center,font=\scriptsize\color{UTPNavy}},
  go/.style={-{Stealth[length=2.2mm]},thick,draw=UTPTeal},
  wcap/.style={font=\scriptsize\bfseries\color{UTPMuted}},
  rel/.style={font=\scriptsize\color{UTPGreen},align=center}]

\node[wcap] at (0,1.15) {Work = Available};
\node[wk] (w0) at (0,0.35) {(1,\,1)};

\node[test] (t1) at (3.5,0.35) {P1 needs (1,\,0)\\$(1,0)\le(1,1)$ --- P1 can finish};
\node[wk]   (w1) at (7.0,0.35) {(2,\,2)};
\node[rel]  at (7.0,-0.62) {P1 releases (1,\,1)};

\node[test] (t2) at (10.5,0.35) {P2 needs (2,\,1)\\$(2,1)\le(2,2)$ --- P2 can finish};
\node[wk]   (w2) at (14.0,0.35) {(3,\,2)};

\draw[go] (w0) -- (t1);
\draw[go] (t1) -- (w1);
\draw[go] (w1) -- (t2);
\draw[go] (t2) -- (w2);

\node[font=\scriptsize\color{UTPInk},align=center,text width=13cm] at (7.0,-1.55)
  {Every process fits into the order, so the state is \textbf{safe} and the
   safe sequence is P1, P2. Had P1 failed its test as well, no order would
   exist and the request would be refused --- not because deadlock had
   happened, but because it could no longer be ruled out.};

\end{tikzpicture}
\end{utpfig}

\begin{pitfall}
An \term{unsafe state} is not identical to an already existing deadlock. Unsafe means the OS can no longer guarantee that every future maximum claim can be satisfied in some order. Deadlock is an actual circular waiting condition.
\end{pitfall}

\section{Detection versus the supplied exercise}
A textbook deadlock-detection algorithm normally uses current outstanding requests rather than each process's declared future maximum. The supplied lab calculates \texttt{Need = Max - Allocation}; this is close to a safety/avoidance test. Treat the code as a useful matrix exercise, but keep the conceptual distinction clear.


\part{Memory Management}
\chapter{Contiguous Memory Allocation}\label{ch:allocation}
\chapterguide{Know that a running process needs memory for code, data, heap, and stack.}{Explain dynamic contiguous allocation, distinguish internal from external fragmentation, trace first-fit/best-fit/worst-fit, and explain how freeing and coalescing change the hole list.}

\section{The allocation problem}
In a simple contiguous-allocation model, each request must be placed into one sufficiently large free region, often called a \term{hole}. If a 120-byte request is placed into a 500-byte hole, the remaining 380 bytes become a smaller free hole.

This model is deliberately simpler than modern virtual memory, but it is useful because it makes fragmentation and online allocation decisions visible.

\section{Internal and external fragmentation}
\term{Internal fragmentation} is unused space inside an allocated region. \term{External fragmentation} is free memory split into separate holes such that a large request cannot fit even when the total free space is sufficient.

\begin{workedexample}
Suppose the free holes are 70, 40, and 60. Total free memory is 170, but a request for 100 cannot be placed because the largest individual hole is only 70. That is external fragmentation.
\end{workedexample}

\section{First fit}
First fit scans from the beginning and chooses the first hole large enough for the request. Its main appeal is simple search behaviour: it can stop as soon as a usable hole is found.

\section{Best fit}
Best fit chooses the smallest hole that can satisfy the current request. This minimizes the immediate remainder, but a very small remainder may be difficult to use later.

\section{Worst fit}
Worst fit chooses the largest available hole, leaving a large remainder. This can preserve medium-size holes but may also consume the largest flexible region too aggressively.

\begin{keyidea}
All three algorithms make an \emph{online} decision. They choose using the requests seen so far; they do not know the complete future workload.
\end{keyidea}

\section{One worked comparison}
For holes 120, 500, and 250 with a request of 200:
\begin{itemize}
  \item first fit chooses 500 because 120 is too small and 500 is the first sufficient hole;
  \item best fit chooses 250 because it is the smallest sufficient hole;
  \item worst fit chooses 500 because it is the largest hole.
\end{itemize}
The policies already create different future hole lists after only one request.

\section{Freeing memory and coalescing holes}
Real allocators must also handle deallocation. If two free holes become physically adjacent, the allocator may \term{coalesce} them into one larger hole.

For example, imagine allocated block A lies between two free regions of 100 and 80. If A is later freed and has size 50, all three adjacent regions can potentially become one 230-unit hole. Without coalescing, the allocator would see several smaller pieces instead.

\section{Compaction is conceptually different}
Coalescing combines only free regions that are already adjacent. \term{Compaction} moves allocated regions so scattered free space becomes contiguous. Moving live memory is expensive and requires address-management support, which is one reason paging becomes attractive in later chapters.

\section{How to compare placement policies fairly}
Keep the initial holes and request order identical. Record after every request:
\begin{itemize}
  \item selected hole;
  \item hole size before allocation;
  \item request size;
  \item remainder;
  \item complete free-hole list.
\end{itemize}
Useful summary measures include total free memory, the largest remaining hole, number of holes, and whether a request failed.

\begin{pitfall}
A result from one request sequence is evidence about that workload, not proof that one policy is universally superior.
\end{pitfall}

\section{What the Lab 7 simulator leaves out}
The lab intentionally ignores many details of production allocators: metadata overhead, alignment, multiple free lists, coalescing strategies, page-based allocation, and concurrency. Its purpose is to make the placement decision itself easy to see.

\begin{quickcheck}
\begin{enumerate}
  \item Can total free memory be greater than a request and still fail? Why?
  \item How is coalescing different from compaction?
  \item Why does ``best fit'' not mean ``best final result''?
\end{enumerate}
\end{quickcheck}

\labsection{7}{Contiguous memory allocation and fragmentation}{sec:lab07}
\begin{labproject}{Compare First Fit, Best Fit, and Worst Fit on the same requests}

\textbf{Coverage.} Chapter~\ref{ch:allocation}: holes, requests, placement rules, external fragmentation.\\
\textbf{Goal.} Watch the free-hole list after every allocation.\\
\textbf{Explain.} Explain how an early placement changes whether later requests can fit.

\courselabflow{Set holes}{Allocate requests}{Trace leftovers}{Compare policies}



\section*{Source code}
\labsourcecard{\texttt{first\_fit.c}}{Choose the first hole that is large enough.}{The search stops immediately when a suitable hole is found.}
\noindent\textbf{Start the file.} \texttt{first\_fit.c}, lines 1--28.
\courselabtrace{First request 212; same starting holes}{First Fit}{500 becomes 288; first suitable hole wins}{Best Fit}{300 becomes 88; smallest suitable hole wins}{Worst Fit}{600 becomes 388; largest suitable hole wins}



\textbf{Use the chapter while coding.} Revisit \hyperref[ch:allocation]{the earlier chapter explanation}; match its operation to the variables and state changes below.

\textbf{Follow the supplied implementation.} choice is a zero-based index and starts at -1 for each request. The break selects the first eligible hole. Update only the chosen hole, and compare the unchanged map before and after the failed request.
\courselabtrace{Trace first\_fit using the supplied requests}{1. Input / state}{Holes start [100, 500, 200, 300, 600]; requests are [212, 417, 112, 426].}{2. Operation}{212 uses hole 1, leaving 288. 417 uses hole 4, leaving 183. 112 then uses hole 1, leaving 176.}{3. Check}{426 cannot fit. Final holes: [100, 176, 200, 300, 183]; no hole is large enough.}


\lstinputlisting[style=utpcnumbered,firstline=1,lastline=28]{all_experiments/Lab07/first_fit.c}

\noindent\textbf{Continue the same file.} \texttt{first\_fit.c}, lines 29--41.
\lstinputlisting[style=utpcnumbered,firstline=29,lastline=41]{all_experiments/Lab07/first_fit.c}

\labsourcecard{\texttt{best\_fit.c}}{Choose the smallest hole that can satisfy the request.}{The whole hole list must be scanned to find the tightest fit.}
\noindent\textbf{Start the file.} \texttt{best\_fit.c}, lines 1--28.
\lstinputlisting[style=utpcnumbered,firstline=1,lastline=28]{all_experiments/Lab07/best_fit.c}

\noindent\textbf{Continue the same file.} \texttt{best\_fit.c}, lines 29--41.
\lstinputlisting[style=utpcnumbered,firstline=29,lastline=41]{all_experiments/Lab07/best_fit.c}

\labsourcecard{\texttt{worst\_fit.c}}{Choose the largest suitable hole.}{The policy tries to leave a relatively large remainder, but the future workload still decides whether that helps.}
\noindent\textbf{Start the file.} \texttt{worst\_fit.c}, lines 1--28.
\lstinputlisting[style=utpcnumbered,firstline=1,lastline=28]{all_experiments/Lab07/worst_fit.c}

\noindent\textbf{Continue the same file.} \texttt{worst\_fit.c}, lines 29--41.
\lstinputlisting[style=utpcnumbered,firstline=29,lastline=41]{all_experiments/Lab07/worst_fit.c}

\section*{Task 1 --- Compile and run}
\begin{lstlisting}[style=utpbash]
cd all_experiments/Lab07
for f in first_fit best_fit worst_fit; do
  gcc -Wall -Wextra -std=c11 "$f.c" -o "$f"
  ./$f
  echo
 done
\end{lstlisting}

For each policy, copy the hole list after every request. Do not compare only the final line; the intermediate state is the algorithm.

\section*{Task 2 --- Visual comparison}
Use the same starting holes: 100, 500, 200, 300, 600. Shade the selected hole for the first request (212) under each policy.

\begin{utptable}
\renewcommand{\arraystretch}{1.18}
\begin{tabularx}{\linewidth}{@{}>{\bfseries\leavevmode\color{UTPNavy}}m{0.19\linewidth}
  *{3}{>{\raggedright\arraybackslash}X}@{}}
\toprule
\rowcolor{UTPPaleBlue!92}
\utpth{Policy} & \utpth{Hole chosen for 212} & \utpth{Remainder created} & \utpth{Later consequence} \\
\midrule
First Fit & & & \\
Best Fit & & & \\
Worst Fit & & & \\
\bottomrule
\end{tabularx}
\end{utptable}

\section*{Task 3 --- Independent verification}
Before looking at the program's second line of output, manually choose a hole for request 417 using each policy. Then compare your prediction with the code.

\section*{Task 4 --- Controlled edit}
Pick one of the three programs. Open it in Vim and change the final request from 426 to 250. Predict whether that request will now fit and which hole will be selected. Recompile and check your prediction.

\begin{keyidea}
``Best Fit'' is the name of a rule, not a guarantee of the best final memory layout. Every policy is making a decision without knowing future requests.
\end{keyidea}

\begin{quickcheck}
\begin{enumerate}
  \item What is external fragmentation?
  \item Why can two placement policies leave different hole lists after the same request?
  \item Why does a request sometimes fail even when the total free memory is large enough?
\end{enumerate}
\end{quickcheck}

\begin{verification}
\textbf{Explain and extend your result.} Trace the first two requests by hand for each policy. Report remaining hole sizes and distinguish failed placement from exhaustion of total free memory.


\begin{itemize}
  \item The hole list is recorded after every request for each policy, not only the final line.
  \item The hole chosen for the first request is compared across the three policies.
  \item Your manual prediction for request 417 matches the program.
  \item The before/after result of the 426 $\rightarrow$ 250 edit is explained.
\end{itemize}
\end{verification}
\end{labproject}


\chapter{Virtual Memory and Page Replacement}\label{ch:vm}
\chapterguide{Understand contiguous allocation from Chapter~\ref{ch:allocation}.}{Explain pages and frames, translate a virtual page through a page table, define page faults, and trace FIFO, LRU, and LFU replacement.}

\section{Why virtual memory changes the problem}
Paging divides virtual memory into fixed-size \term{pages} and physical memory into equal-size \term{frames}. A page table records which frame currently contains each resident virtual page. Because pages need not occupy adjacent frames, paging removes external fragmentation from physical allocation, though the last page of an object can still create internal fragmentation.

\section{Address translation}
Divide the virtual address by the page size. The quotient is the page number;
the remainder is the offset. The page table replaces the page number with a
frame number, while the offset stays unchanged. In compact form:
\texttt{physical address = frame $\times$ page size + offset}.

\begin{mlfigure}
\centering
\begin{tikzpicture}[
  cell/.style={draw=MLBlue,fill=MLPaleBlue,minimum width=1.6cm,minimum height=0.7cm,
               font=\small\color{MLNavy},align=center},
  off/.style={draw=MLOrange,fill=MLPaleOrange,minimum width=1.6cm,minimum height=0.7cm,
              font=\small\color{MLNavy},align=center},
  pte/.style={draw=MLTeal,fill=MLPaleTeal,minimum width=1.2cm,minimum height=0.58cm,
              font=\scriptsize\color{MLNavy}},
  lab/.style={font=\scriptsize\color{MLMuted},align=center},
  ar/.style={-{Stealth[length=2.2mm]},thick,draw=MLTeal}
]
% The two address fields are STACKED, not side by side, so the page-number
% path and the offset path never cross each other or the page table.

% ---- headings ----
\node[lab] at (0,3.5)   {\textbf{Virtual address}\\ $v = 9\,000$, \ $P = 4\,096$};
\node[lab] at (4.9,4.0) {\textbf{Page table}};
\node[lab] at (9.6,3.5) {\textbf{Physical address}};

% ---- page table, positioned so entry 2 lines up with the page row ----
\node[pte] (p0) at (4.9,3.2) {0 $\to$ 7};
\node[pte] (p1) at (4.9,2.6) {1 $\to$ 3};
\node[pte,draw=MLRed,fill=MLPaleRed,very thick] (p2) at (4.9,2.0) {2 $\to$ 5};
\node[pte] (p3) at (4.9,1.4) {3 $\to$ 1};

% ---- page row (y = 2.0) and offset row (y = 0.5) ----
\node[cell] (vp) at (0,2.0) {page 2};
\node[off]  (vo) at (0,0.5) {offset 808};
\node[cell] (pf) at (9.6,2.0) {frame 5};
\node[off]  (po) at (9.6,0.5) {offset 808};

% ---- arithmetic, below its own column ----
\node[lab] at (0,-0.35)   {$\lfloor 9000/4096 \rfloor = 2$\\ $9000 \bmod 4096 = 808$};
\node[lab] at (9.6,-0.35) {$5 \times 4096 + 808$\\ $= 21\,288$};

% ---- translated: page number in, frame out ----
\draw[ar] (vp.east) -- node[lab,above]{look up} (p2.west);
\draw[ar] (p2.east) -- node[lab,above]{gives frame} (pf.west);

% ---- untranslated: the offset passes straight under the table ----
\draw[ar,draw=MLOrange] (vo.east) --
      node[lab,above]{copied unchanged --- it never enters the page table} (po.west);
\end{tikzpicture}

\smallskip
{\scriptsize Only the \emph{page number} is translated. The offset passes through
untouched, which is why the page size fixes how many bits each part uses --- and
why page size is always a power of two.}
\end{mlfigure}

\begin{keyidea}
Translation replaces the page number and copies the offset. If a virtual address is contiguous inside one page, its physical address is contiguous too; across a page boundary it need not be, because the next page may live in any frame.
\end{keyidea}

\section{Page faults}
A \term{page fault} occurs when a referenced virtual page is not currently resident in a usable physical frame. The OS must locate the page on backing storage (or otherwise materialize it), obtain a frame, update mapping metadata, and restart/resume the faulting instruction as appropriate.

\section{Replacement policies}
When no free frame exists, the OS needs a victim page.

\begin{mltable}
\begin{tabularx}{\linewidth}{@{}>{\bfseries\leavevmode\color{MLNavy}}p{0.16\linewidth}p{0.29\linewidth}X@{}}
\toprule
\rowcolor{MLPaleBlue!92}
Policy & Victim rule & State required \\
\midrule
FIFO & Page resident for the longest time & Arrival/load order \\
LRU & Page whose most recent use is farthest in the past & Recency information \\
LFU & Page with the smallest reference count & Frequency counts plus a tie-breaker \\
\bottomrule
\end{tabularx}
\end{mltable}

\section{FIFO}
FIFO treats loaded pages like a queue. It is easy to implement but does not use actual access behavior. FIFO can exhibit \term{Belady's anomaly}: for some reference strings, adding more frames can increase the number of faults.

\section{LRU}
LRU approximates the principle of temporal locality: a page used recently is likely to be used again soon. Exact LRU can be expensive, so real systems often use hardware-assisted approximations.

\section{LFU}
LFU retains frequently referenced pages. Pure cumulative counts can preserve stale historical popularity, so practical frequency-based schemes often age counts or use bounded windows.

\begin{workedexample}
With three frames and reference string 1, 2, 3, 1, 4: the first three references fault and fill the frames. The second reference to 1 is a hit. On reference 4, FIFO evicts 1 if 1 was loaded earliest, whereas LRU evicts 2 because 2 is the least recently used at that moment.
\end{workedexample}

\section{Locality explains why replacement policies can work}
Programs often reuse recently accessed instructions and data (temporal locality) and access addresses close to one another (spatial locality). Page-replacement policies try to exploit these patterns indirectly. LRU uses recent history; LFU uses frequency history; FIFO deliberately ignores reference history and uses load order only.

\begin{keyidea}
A page-replacement algorithm is a prediction policy: it tries to choose a resident page whose removal will hurt future execution as little as possible. The policy never knows the future perfectly, so the workload matters.
\end{keyidea}

\section{How to trace replacement professionally}
A trustworthy trace shows the reference, frame contents after the reference, hit/fault result, and enough policy state to justify a victim. Totals are useful only after individual decisions can be checked.

\begin{quickcheck}
\begin{enumerate}
  \item Which part of a virtual address remains unchanged during page-table translation?
  \item What is the difference between a page hit and a page fault?
  \item Why can LRU and LFU choose different victims even on the same reference string?
\end{enumerate}
\end{quickcheck}

\labsection{8}{Virtual memory and page replacement}{sec:lab08}
\begin{labproject}{Trace FIFO, LRU, and LFU using the same reference string}

\oslabflow{Chapter 8: page hits, page faults, replacement state.}{Run three policies on one reference string so the victim rule is the only major change.}{Justify at least one replacement decision from each policy's stored state.}

\section*{Source code}
\ossourcecard{\texttt{fifo.c}}{Replace the page that has been resident the longest.}{\texttt{next} rotates through frame positions; page-use history is ignored.}
\lstinputlisting[style=oslabc]{all_experiments/Lab08/fifo.c}

\ossourcecard{\texttt{lru.c}}{Replace the least recently used resident page.}{\texttt{last\_used[]} stores the most recent reference step for each frame.}
\lstinputlisting[style=oslabc]{all_experiments/Lab08/lru.c}

\ossourcecard{\texttt{lfu.c}}{Replace the least frequently used resident page; use recency to break ties.}{Two pieces of policy state are tracked: \texttt{frequency[]} and \texttt{last\_used[]}.}
\lstinputlisting[style=oslabc]{all_experiments/Lab08/lfu.c}

\section*{Task 1 --- Compile and collect traces}
\begin{lstlisting}[style=osbash]
cd all_experiments/Lab08
for f in fifo lru lfu; do
  gcc -Wall -Wextra -std=c11 "$f.c" -o "$f"
  ./$f
  echo
 done
\end{lstlisting}
Fill in the final hit/fault totals.

\begin{mltable}
\renewcommand{\arraystretch}{1.18}
\begin{tabularx}{\linewidth}{@{}
  >{\raggedright\arraybackslash}p{0.20\linewidth}
  *{2}{>{\raggedright\arraybackslash}X}@{}}
\toprule
\rowcolor{UTPPaleBlue!92}
\textbf{Policy} & \textbf{Hits} & \textbf{Faults} \\
\midrule
FIFO & & \\
LRU & & \\
LFU & & \\
\bottomrule
\end{tabularx}
\end{mltable}

\section*{Task 2 --- Explain one victim from each policy}
Choose a step where all three frames are full and a fault occurs. For FIFO identify the oldest frame position; for LRU compare last-use steps; for LFU compare frequencies and then the tie-break rule if needed.

\section*{Task 3 --- Verify the beginning by hand}
Manually trace the first six references: 7, 0, 1, 2, 0, 3. Compare your frame table with all three programs. The first few decisions should make the policies' state differences visible.

\section*{Task 4 --- Controlled edit}
Choose one policy. In Vim change the final reference in its array from 2 to 7. Predict whether that final access is a hit or a fault for that policy, recompile, and verify.

\begin{quickcheck}
\begin{enumerate}
  \item What event causes a page fault?
  \item What state does LRU need that FIFO does not?
  \item Why can two replacement policies choose different victims for the same reference?
\end{enumerate}
\end{quickcheck}

\begin{labdeliverable}
Submit the three totals, one justified victim decision per policy, your six-reference hand trace, and the prediction/result for the final-reference edit.
\end{labdeliverable}
\end{labproject}



\part{File Systems and Storage}
\chapter{Directories and the File-System Namespace}\label{ch:directories}
\chapterguide{Chapter 2: files and metadata, especially inode numbers and permissions.}{read absolute and relative paths, distinguish directory entries from inodes, compare hard and symbolic links, and follow recursive traversal in Lab 9.}

\section{Why directories exist}
Storage blocks do not naturally come with human-friendly names such as \texttt{report.txt}. A file system therefore provides a \term{namespace}: a structured way to associate names with file-system objects.

A directory is part of that namespace. In a Unix-like file system, directory entries associate names with underlying objects identified by metadata such as inode numbers.

\section{From flat names to a hierarchy}
Textbooks often describe several directory organizations:
\begin{itemize}
  \item \textbf{single-level:} every file shares one directory;
  \item \textbf{two-level:} each user receives a separate directory;
  \item \textbf{tree-structured:} directories contain files and subdirectories;
  \item \textbf{graph-like structures:} links allow an object to be reachable through more than one path.
\end{itemize}
Modern general-purpose systems use hierarchical trees with linking mechanisms.

\begin{mlfigure}
\begin{tikzpicture}[
  grow=down,
  level distance=1.55cm,
  level 1/.style={sibling distance=4.2cm},
  level 2/.style={sibling distance=2.35cm},
  dir/.style={draw=MLBlue,fill=MLPaleBlue,rounded corners=2mm,minimum width=2.6cm,minimum height=0.75cm,font=\small\bfseries\color{MLNavy}},
  file/.style={draw=MLTeal,fill=MLPaleTeal,rounded corners=1mm,minimum width=1.9cm,minimum height=0.65cm,font=\scriptsize\color{MLNavy}},
  edge from parent/.style={draw=MLMuted,thick,-{Stealth[length=1.8mm]}}
]
\node[dir] {root}
 child {node[dir] {home} child {node[file] {notes.txt}} child {node[file] {code.c}}}
 child {node[dir] {etc} child {node[file] {config}}}
 child {node[dir] {tmp}};
\end{tikzpicture}
\end{mlfigure}

\section{Absolute and relative paths}
An \term{absolute path} begins at the root directory, for example \texttt{/home/student/notes.txt}. A \term{relative path} begins at the process's current working directory.

Two special names appear frequently:
\begin{itemize}
  \item \texttt{.} means the current directory;
  \item \texttt{..} means the parent directory.
\end{itemize}
The shell maintains a current working directory for itself, and a child process normally inherits that starting directory.

\section{Pathname resolution is repeated lookup}
To resolve a path such as \texttt{/home/student/notes.txt}, the system conceptually starts at the root, finds the entry named \texttt{home}, then inside that directory finds \texttt{student}, and finally finds \texttt{notes.txt}.

This is why permissions on intermediate directories matter: reaching a file requires traversing the path, not merely knowing the final filename.

\section{Name, metadata, and data are different layers}
A filename is not the complete file. Metadata can include object type, owner, permissions, timestamps, size, link count, and pointers or mappings to data blocks.

\begin{keyidea}
The namespace tells us \emph{how an object is named}. Metadata tells us \emph{what the object is and its properties}. Storage allocation tells us \emph{where its bytes are represented on storage}.
\end{keyidea}

\section{Hard links}
A hard link gives another directory entry to the same underlying inode. The two names are peers: neither is inherently the ``real'' name. Removing one name does not remove the file data while another hard link still refers to the inode.

The inode's link count helps the file system track how many directory entries refer to it.

\section{Symbolic links}
A symbolic link is a separate small file-system object whose content is a pathname. Following the link causes another pathname lookup.

\begin{mltable}
\renewcommand{\arraystretch}{1.18}
\begin{tabularx}{\linewidth}{@{}
  >{\raggedright\arraybackslash}p{0.28\linewidth}
  >{\raggedright\arraybackslash}X
  >{\raggedright\arraybackslash}X@{}}
\toprule
\rowcolor{UTPPaleBlue!92}
\textbf{Feature} & \textbf{Hard link} & \textbf{Symbolic link} \\
\midrule
Underlying inode & same inode as target & separate inode \\
Stores a pathname & no & yes \\
Can become dangling & not in the same way & yes, if target path disappears \\
Recursive traversal concern & ordinary name to same object & can create path-following cycles \\
\bottomrule
\end{tabularx}
\end{mltable}

\section{Recursive traversal}
A directory walker opens a directory, reads each entry, and recursively repeats the operation for child directories. Two practical rules matter:
\begin{enumerate}
  \item do not recurse into \texttt{.} and \texttt{..};
  \item decide deliberately whether symbolic links will be followed.
\end{enumerate}
Blindly following links can revisit the same locations indefinitely if links form a cycle.

\section{Directory operations and consistency}
Common namespace operations include create, list, rename, link, unlink, and traverse. A real file system must update its metadata consistently even if two processes act at once or a crash interrupts an update.

The lab stays at an observable introductory level: you create a small tree, inspect inode identity and links, and compare the shell's view with a recursive C walker.

\begin{quickcheck}
\begin{enumerate}
  \item What is the starting point of an absolute path? A relative path?
  \item Why can two different names report the same inode number?
  \item Why might a directory walker choose not to follow symbolic links?
\end{enumerate}
\end{quickcheck}

\labsection{9}{File-system namespace, paths, and links}{sec:lab09}
\begin{labproject}{Build a directory tree and investigate names, inode identity, and links}
\textbf{Suggested time:} 90--100 minutes \hfill
\textbf{Level:} introductory/intermediate\par
\textbf{Starter files:} \texttt{make\_demo\_tree.sh}, \texttt{directory\_walk.c}.\par
\textbf{Evidence to keep:} hand-drawn tree, pathname examples, inode/link table, and recursive-walk output.

\begin{labmode}
Use the real Linux/WSL file system. No graphics library is required. The goal is to understand what a directory tree represents, not to draw a decorative visualization.
\end{labmode}

\medskip\noindent\textbf{Learning outcomes.}
\begin{enumerate}
  \item distinguish absolute and relative pathnames;
  \item explain why a filename is not the same thing as the file's metadata/inode;
  \item distinguish a hard link from a symbolic link at an introductory level;
  \item explain how recursive directory traversal follows the namespace.
\end{enumerate}

\medskip\noindent\textbf{Part A --- Create and draw a known hierarchy.}
\begin{lstlisting}[style=osbash]
cd experiments/Lab09
bash make_demo_tree.sh
find demo_tree -maxdepth 4 -print
\end{lstlisting}
Draw the tree by hand. Label directories, regular files, the hard link, and the symbolic link. Your diagram should reflect the names reported by \texttt{find}.

\medskip\noindent\textbf{Part B --- Navigate using different path forms.}
Move into \texttt{demo\_tree/home/student}. Display \texttt{notes.txt} using: (1) a relative path, (2) an absolute path, and (3) a relative path containing \texttt{..}. Run \texttt{realpath notes.txt} and explain what it resolves.

\medskip\noindent\textbf{Part C --- Compare names with inode identity.}
\begin{lstlisting}[style=osbash]
ls -li notes.txt notes-hardlink.txt code-link.c
stat notes.txt
stat notes-hardlink.txt
readlink code-link.c
\end{lstlisting}
Record inode numbers and link counts. Explain why the hard-linked names refer to the same inode while a symbolic link stores a pathname to another object.

\begin{keyidea}
A directory gives names to objects in a namespace. The inode/metadata identifies properties of the underlying file object; storage blocks hold the file's data. These are related layers, not the same thing.
\end{keyidea}

\medskip\noindent\textbf{Part D --- Walk the tree recursively.}
\begin{lstlisting}[style=osbash]
gcc -Wall -Wextra directory_walk.c -o directory_walk
./directory_walk demo_tree
\end{lstlisting}
Find the recursive function call in the source. Explain the stopping rule and why the program deliberately does not recurse through symbolic links.

\medskip\noindent\textbf{Part E --- Controlled namespace change.}
Rename \texttt{notes.txt} while the hard link still exists. Inspect both names again. Then remove one hard link and inspect the remaining link count. Record the observations before deleting the demo tree.

\begin{verification}
Submit the tree, three pathname examples, the inode/link comparison, and a short explanation of why following symbolic links recursively without a policy could create cycles.
\end{verification}
\end{labproject}


\chapter{File Organization and Disk-Block Allocation}\label{ch:filesystems}
\chapterguide{Read Chapter~\ref{ch:directories} first.}{Distinguish application record formats from file-system block allocation, compare basic access methods, and explain contiguous, linked, and indexed allocation at an introductory level.}

\section{A file is a byte sequence; applications add meaning}
To the operating system, a regular file can be treated as a sequence of bytes plus metadata. Applications decide how those bytes represent records such as students, employees, images, or database pages.

This gives us two layers to keep separate:
\begin{itemize}
  \item \textbf{application organization:} how fields and records are encoded inside the file;
  \item \textbf{file-system allocation:} how logical file blocks are mapped onto storage.
\end{itemize}

\section{Text and binary records}
Text formats convert values into readable characters. Binary formats can store compact machine-oriented representations.

\begin{mltable}
\begin{tabularx}{\linewidth}{@{}>{\bfseries\leavevmode\color{MLNavy}}p{0.19\linewidth}XX@{}}
\toprule
\rowcolor{MLPaleBlue!92}
Format & Strength & Caution \\
\midrule
Text & readable with ordinary tools, easy to exchange & parsing and formatting add overhead \\
Raw C struct & compact and easy for a controlled exercise & padding, endianness, and type sizes can reduce portability \\
\bottomrule
\end{tabularx}
\end{mltable}

\begin{pitfall}
``Binary'' does not automatically mean ``better.'' A durable binary format normally defines its byte layout explicitly instead of simply dumping a compiler-dependent in-memory structure.
\end{pitfall}

\section{Sequential access and sequential search}
A small record file can be searched by reading one record after another until the key is found. This is easy to implement and understand, but the worst case may inspect every record.

Larger systems introduce indexes so a key can point more directly to the needed location. The idea is similar to a book index: extra structure is maintained so lookup does not require scanning every page.

\section{Logical blocks and physical placement}
A file system commonly views a file as logical blocks numbered 0, 1, 2, and so on. The allocation method determines where those logical blocks are stored physically.

\section{Contiguous allocation}
All file blocks occupy consecutive storage blocks. Sequential and random access are simple, but file growth can be difficult when the following physical blocks are already occupied.

\section{Linked allocation}
Each allocated block identifies the next block in the file. Files can grow without requiring adjacent space, but following a long chain is inefficient for direct/random access and damaged pointers are a concern.

\section{Indexed allocation}
A separate index block stores references to the file's data blocks. Logical block 0 can be found through the first index entry, logical block 1 through the second, and so on. The data blocks themselves do not need to be adjacent.

\begin{mlfigure}
\begin{tikzpicture}[
  b/.style={draw=MLBlue,fill=MLPaleBlue,rounded corners=1mm,minimum width=2.2cm,minimum height=0.7cm,font=\small\bfseries\color{MLNavy}},
  d/.style={draw=MLTeal,fill=MLPaleTeal,rounded corners=1mm,minimum width=1.7cm,minimum height=0.65cm,font=\scriptsize\color{MLNavy}},
  arr/.style={-{Stealth[length=2mm]},draw=MLMuted,thick}
]
\node[b] (idx) {Index block 7};
\node[d,right=2.2cm of idx,yshift=1.0cm] (d1) {block 12};
\node[d,right=2.2cm of idx] (d2) {block 19};
\node[d,right=2.2cm of idx,yshift=-1.0cm] (d3) {block 31};
\draw[arr] (idx)--(d1); \draw[arr] (idx)--(d2); \draw[arr] (idx)--(d3);
\end{tikzpicture}
\end{mlfigure}

\begin{mltable}
\begin{tabularx}{\linewidth}{@{}>{\bfseries\leavevmode\color{MLNavy}}p{0.18\linewidth}p{0.27\linewidth}X@{}}
\toprule
\rowcolor{MLPaleBlue!92}
Method & Main advantage & Main limitation \\
\midrule
Contiguous & simple, efficient direct access & growth and external fragmentation \\
Linked & easy growth without contiguous space & poor direct access \\
Indexed & direct access without contiguous data blocks & index-space overhead \\
\bottomrule
\end{tabularx}
\end{mltable}

\section{Free-space management}
A file system also needs to know which storage blocks are free. A bitmap is one common idea: each bit records whether one block is free or allocated. Other structures can use lists or grouped free blocks.

The Lab 10 simulator uses a simple integer array as an easy-to-see allocation map. It is a model, not a complete disk file system.

\section{Validate first, commit second}
Suppose a request needs one index block and three data blocks. If the program marks the first two blocks as used and only later discovers that the third is unavailable, it has partially changed state.

A safer introductory pattern is:
\begin{enumerate}
  \item validate the complete request without modifying shared state;
  \item if every block is valid and free, commit all allocation changes;
  \item otherwise, leave the map unchanged.
\end{enumerate}
This transaction-like reasoning appears throughout systems software.

\section{Durability and crash consistency}
Real file systems must survive interruptions during metadata updates. Techniques such as journaling, copy-on-write, checksums, and careful write ordering help prevent inconsistent states after a crash. These mechanisms are beyond the Lab 10 implementation, but the validate-then-commit habit is an accessible first step toward the same design concern.

\section{Keep the layers separate}
When an application stores a student record, it decides which bytes represent the ID, name, and GPA. When the file system stores the file, it decides how file offsets map to blocks. One file can change its physical allocation without changing its application record format.

\begin{quickcheck}
\begin{enumerate}
  \item Which layer decides how an employee record is encoded?
  \item Which allocation method uses a dedicated block of pointers?
  \item Why should a multi-block allocation request be fully checked before state is modified?
\end{enumerate}
\end{quickcheck}

\labsection{10}{File records and indexed block allocation}{sec:lab10}
\begin{labproject}{Separate application record formats from file-system block allocation}
\textbf{Suggested time:} 100--110 minutes \hfill
\textbf{Level:} intermediate\par
\textbf{Starter files:} \texttt{text\_records.c}, \texttt{binary\_records.c}, \texttt{indexed\_allocation.c}.\par
\textbf{Evidence to keep:} text/binary comparison, successful and missing search cases, and allocation-state trace.

\begin{labmode}
This practical deliberately works at two layers. First the application decides how records become bytes. Then a separate simulator shows how a file system could map a file through an index block to storage blocks.
\end{labmode}

\medskip\noindent\textbf{Learning outcomes.}
\begin{enumerate}
  \item compare text and raw binary record encodings;
  \item explain why sequential record search is simple but may require scanning many records;
  \item describe indexed allocation and explain why validate-then-commit protects allocation state.
\end{enumerate}

\medskip\noindent\textbf{Part A --- Text records are easy to inspect.}
\begin{lstlisting}[style=osbash]
cd experiments/Lab10
gcc -Wall -Wextra text_records.c -o text_records
./text_records
cat employees.txt
ls -l employees.txt
\end{lstlisting}
Identify the record delimiter and field delimiter. Change the target ID to a missing value, rerun, and ensure the program reports the missing record clearly.

\medskip\noindent\textbf{Part B --- Binary records preserve machine representation.}
\begin{lstlisting}[style=osbash]
gcc -Wall -Wextra binary_records.c -o binary_records
./binary_records
ls -l students.bin
hexdump -C students.bin | head
\end{lstlisting}
If \texttt{hexdump} is unavailable, use \texttt{od -Ax -tx1 students.bin | head}. Compare readability and size with the text file.

\begin{pitfall}
Writing a C structure directly with \texttt{fwrite()} is convenient for a controlled lab, but it is not a portable long-term interchange format. Structure padding, integer representation, endianness, and compiler choices can differ.
\end{pitfall}

\medskip\noindent\textbf{Part C --- Trace indexed allocation.}
\begin{lstlisting}[style=osbash]
gcc -Wall -Wextra indexed_allocation.c -o indexed_allocation
./indexed_allocation
\end{lstlisting}
Draw file A as one index block pointing to its three data blocks. Then explain why file B is rejected and verify that the printed allocation map is unchanged after that rejected request.

\medskip\noindent\textbf{Part D --- Explain validate first, commit second.}
Read the two functions \texttt{request\_is\_valid()} and \texttt{allocate\_file()}. In 4--6 sentences explain why checking the whole request before marking any block avoids a partially committed allocation after an error.

\medskip\noindent\textbf{Part E --- Connect the two layers.}
Draw two small diagrams: one showing fields inside an employee/student record, and one showing logical file blocks reached through an index block. Label which decisions belong to the application and which belong to the file system.

\begin{verification}
A complete submission includes a text/binary comparison, one successful and one missing record search, an indexed-allocation diagram, and evidence that a rejected overlapping allocation leaves block state unchanged.
\end{verification}

\begin{tryit}
Optional extension: add a deletion operation to the allocation simulator that frees an index block and all data blocks belonging to one simulated file. Print the allocation map before and after deletion.
\end{tryit}
\end{labproject}


\chapter*{Revision Map and Capstone Practice}
\addcontentsline{toc}{chapter}{Revision Map and Capstone Practice}
\markboth{Revision Map and Capstone Practice}{}

\section*{The five exam-solving habits}
Most Operating Systems questions become manageable if you use a repeatable method:
\begin{enumerate}
  \item \textbf{Name the state.} Write the queue, matrices, frame contents, free blocks, or directory nodes before applying a policy.
  \item \textbf{Name the event.} Arrival, dispatch, quantum expiry, resource request, memory request, page reference, or file-block allocation.
  \item \textbf{Name the policy.} FCFS, SJF, Priority, RR, Banker's safety rule, first/best/worst fit, FIFO/LRU/LFU.
  \item \textbf{Update exactly once.} Change the state according to the event and policy; do not mentally skip intermediate states.
  \item \textbf{Compute metrics last.} Waiting time, turnaround, page faults, fragmentation, or safe sequence should follow from the completed trace.
\end{enumerate}

\begin{utptable}
\renewcommand{\arraystretch}{1.18}
\begin{tabularx}{\linewidth}{@{}>{\bfseries\leavevmode\color{UTPNavy}}m{0.24\linewidth}
  >{\raggedright\arraybackslash}m{0.30\linewidth}
  >{\raggedright\arraybackslash}X@{}}
\toprule
\rowcolor{UTPPaleBlue!92}
\utpth{Question type} & \utpth{State to draw} & \utpth{Final evidence} \\
\midrule
CPU scheduling & Ready queue + Gantt chart & Completion, turnaround, waiting, response \\
Banker's algorithm & Allocation, Max, Need, Work, Finish & Safe sequence or reason none exists \\
Memory fit & Free-block table & Chosen block + fragment for each request \\
Page replacement & Frame table after every reference & Hit/fault count and victim sequence \\
File allocation & Free/used block map & Valid block mapping with conflicts rejected \\
\bottomrule
\end{tabularx}
\end{utptable}

\section*{Capstone practical}
Build one command-line simulator that lets a user choose between CPU scheduling, memory allocation, or page replacement. Separate \emph{policy} from \emph{state}: for example, one common frame-state representation can be reused with FIFO, LRU, and LFU victim-selection functions.

\begin{miniproject}
\textbf{Minimum features.}
\begin{itemize}
  \item Scheduling: FCFS, SJF, Priority, RR with a printed Gantt trace and average waiting/turnaround time.
  \item Memory: first, best, worst fit with fragmentation results.
  \item Virtual memory: FIFO, LRU, LFU with per-reference frames, hits, and faults.
  \item Input validation and deterministic tie-breaking rules.
  \item A short README explaining assumptions and how to compile/run the program.
\end{itemize}
\end{miniproject}

\begin{decisionmemo}
When two algorithms give different answers, do not ask only ``which is better?'' Ask which objective matters: responsiveness, throughput, fairness, memory utilization, implementation overhead, locality, or predictable behavior.
\end{decisionmemo}

\section*{Final checklist}
\begin{quickcheck}
\begin{enumerate}
  \item Can you explain the parent/child return values of \texttt{fork()} without running code?
  \item Can you derive waiting and turnaround time from a Gantt chart?
  \item Can you distinguish a race condition from deadlock?
  \item Can you compute Need and search for a Banker's safe sequence?
  \item Can you distinguish internal from external fragmentation?
  \item Can you trace FIFO, LRU, and LFU one reference at a time?
  \item Can you explain why directory hierarchy is not the same thing as physical disk layout?
  \item Can you compare text/binary records and contiguous/linked/indexed allocation?
\end{enumerate}
\end{quickcheck}


\end{document}
